\input{preamble}

% OK, start here.
%
\begin{document}

\title{Champs}


\maketitle

\phantomsection
\label{section-phantom}

\tableofcontents

\section{Introduction}
\label{section-introduction}

\noindent
Dans ce très court chapitre, nous introduisons les champs et les
champs en groupoïdes. Voir \cite{DM} et \cite{Vis2}.


\section{Préfaisceaux de morphismes associés aux catégories fibrées}
\label{section-morphisms}

\noindent
Soit $\mathcal{C}$ une catégorie.
Soit $p : \mathcal{S} \to \mathcal{C}$ une catégorie fibrée,
voir Catégories, section \ref{categories-section-fibred-categories}.
Supposons que $x, y\in \Ob(\mathcal{S}_U)$ soient des
objets de la catégorie fibre au-dessus de $U$. Nous allons définir
un foncteur
$$
\mathit{Mor}(x, y) : (\mathcal{C}/U)^{opp} \longrightarrow \textit{Sets}.
$$
Autrement dit, ce sera un préfaisceau sur $\mathcal{C}/U$, voir
Sites, définition \ref{sites-definition-presheaf}.
Faisons un choix d'images inverses comme dans
Catégories,
définition \ref{categories-definition-pullback-functor-fibred-category}.
Alors, pour $f : V \to U$, nous posons
$$
\mathit{Mor}(x, y)(f : V \to U) =
\Mor_{\mathcal{S}_V}(f^\ast x, f^\ast y).
$$
Soit $f' : V' \to U$ un second objet de $\mathcal{C}/U$.
Nous devons aussi définir l'application de restriction correspondant à un
morphisme $g : V'/U  \to V/U$ dans $\mathcal{C}/U$,
autrement dit $g : V' \to V$ et $f' = f \circ g$.
Ce sera une application
$$
\Mor_{\mathcal{S}_V}(f^\ast x, f^\ast y)
\longrightarrow
\Mor_{\mathcal{S}_{V'}}({f'}^\ast x, {f'}^\ast y), \quad
\phi \longmapsto \phi|_{V'}
$$
Cette application sera essentiellement $g^\ast$, à ceci près qu'elle transforme
un élément $\phi$ du membre de gauche en un élément
$g^\ast \phi$
de $\Mor_{\mathcal{S}_{V'}}(g^\ast f^\ast x, g^\ast f^\ast y)$.
À ce stade, nous utilisons la transformation $\alpha_{g, f}$ du
lemme \ref{categories-lemma-fibred} de Catégories.
En formule, l'application de restriction est décrite par
$$
\phi|_{V'} =
(\alpha_{g, f})_y^{-1} \circ
g^\ast \phi \circ
(\alpha_{g, f})_x.
$$
Bien entendu, personne ne conçoit cette application de restriction de cette
manière. Nous ne le ferons qu'une fois, afin de vérifier le lemme suivant.

\begin{lemma}
\label{lemma-painful}
Cela définit effectivement un préfaisceau.
\end{lemma}

\begin{proof}
Soient $g : V'/U \to V/U$ comme ci-dessus et, de même,
$g' : V''/U \to V'/U$ des morphismes de $\mathcal{C}/U$.
Ainsi $f' = f \circ g$ et $f'' = f' \circ g' = f \circ g \circ g'$.
Soit $\phi \in \Mor_{\mathcal{S}_V}(f^\ast x, f^\ast y)$.
Nous avons alors
\begin{eqnarray*}
& &
(\alpha_{g \circ g', f})_y^{-1} \circ
(g \circ g')^\ast \phi \circ
(\alpha_{g \circ g', f})_x
\\
& = &
(\alpha_{g \circ g', f})_y^{-1} \circ
(\alpha_{g', g})_{f^*y}^{-1} \circ
(g')^*g^\ast \phi \circ
(\alpha_{g', g})_{f^*x} \circ
(\alpha_{g \circ g', f})_x
\\
& = &
(\alpha_{g', f'})_y^{-1} \circ
(g')^*(\alpha_{g, f})_y^{-1} \circ
(g')^* g^\ast \phi \circ
(g')^*(\alpha_{g, f})_x
\circ
(\alpha_{g', f'})_x
\\
& = &
(\alpha_{g', f'})_y^{-1} \circ
(g')^*\Big(
(\alpha_{g, f})_y^{-1} \circ
g^\ast \phi \circ
(\alpha_{g, f})_x
\Big) \circ
(\alpha_{g', f'})_x
\end{eqnarray*}
ce qui est précisément l'égalité voulue,
$\phi|_{V''} = (\phi|_{V'})|_{V''}$.
La première égalité vaut parce que
$\alpha_{g', g}$ est une transformation de foncteurs, de sorte que
$$
\xymatrix{
(g \circ g')^*f^*x
\ar[rr]_{(g \circ g')^\ast \phi}
\ar[d]_{(\alpha_{g', g})_{f^*x}} & &
(g \circ g')^*f^*y
\ar[d]^{(\alpha_{g', g})_{f^*y}} \\
(g')^*g^*f^*x
\ar[rr]^{(g')^*g^\ast \phi} & &
(g')^*g^*f^*y
}
$$
est commutatif. La deuxième égalité résulte de la propriété (d) d'un
pseudo-foncteur, puisque $f' = f \circ g$ (voir
Catégories, définition \ref{categories-definition-functor-into-2-category}).
La dernière égalité découle du fait que $(g')^*$ est un foncteur.
\end{proof}

\noindent
Désormais, nous omettrons souvent de mentionner les transformations
$\alpha_{g, f}$ et identifierons simplement les foncteurs
$g^* \circ f^*$ et $(f \circ g)^*$. En particulier,
pour $g : V'/U \to V/U$, les applications de restriction
du préfaisceau $\mathit{Mor}(x, y)$ seront parfois notées
$\phi \mapsto g^*\phi$. Nous formalisons cette construction dans une définition.

\begin{definition}
\label{definition-mor-presheaf}
Soit $\mathcal{C}$ une catégorie.
Soit $p : \mathcal{S} \to \mathcal{C}$ une catégorie fibrée,
voir Catégories, section \ref{categories-section-fibred-categories}.
Étant donnés un objet $U$ de $\mathcal{C}$ et des objets
$x$, $y$ de la catégorie fibre, le {\it préfaisceau
des morphismes de $x$ vers $y$} est le préfaisceau
$$
(f : V \to U) \longmapsto \Mor_{\mathcal{S}_V}(f^*x, f^*y)
$$
décrit ci-dessus. On le note $\mathit{Mor}(x, y)$.
Le sous-préfaisceau $\mathit{Isom}(x, y)$ dont la valeur
au-dessus de $V$ est l'ensemble des isomorphismes
$f^*x \to f^*y$ dans la catégorie fibre $\mathcal{S}_V$
est appelé le {\it préfaisceau des isomorphismes de $x$ vers $y$}.
\end{definition}

\noindent
Si $\mathcal{S}$ est fibrée en groupoïdes, alors bien entendu
$\mathit{Isom}(x, y) = \mathit{Mor}(x, y)$, et il est
d'usage d'employer la notation $\mathit{Isom}$.

\begin{lemma}
\label{lemma-presheaf-mor-map-fibred-categories}
Soit $F : \mathcal{S}_1 \to \mathcal{S}_2$ un $1$-morphisme de catégories
fibrées au-dessus de la catégorie $\mathcal{C}$. Soient
$U \in \Ob(\mathcal{C})$ et $x, y\in \Ob((\mathcal{S}_1)_U)$.
Alors $F$ définit un morphisme canonique de préfaisceaux
$$
\mathit{Mor}_{\mathcal{S}_1}(x, y)
\longrightarrow
\mathit{Mor}_{\mathcal{S}_2}(F(x), F(y))
$$
sur $\mathcal{C}/U$.
\end{lemma}

\begin{proof}
D'après la définition \ref{categories-definition-fibred-categories-over-C}
de Catégories, le foncteur $F$ envoie les morphismes fortement cartésiens sur
des morphismes fortement cartésiens. Ainsi, si $f : V \to U$ est un morphisme
de $\mathcal{C}$, il existe des isomorphismes canoniques
$\alpha_V : f^*F(x) \to F(f^*x)$,
$\beta_V : f^*F(y) \to F(f^*y)$ tels que $f^*F(x) \to F(f^*x) \to F(x)$
soit le morphisme canonique $f^*F(x) \to F(x)$, et de même pour $\beta_V$.
Nous pouvons donc définir
$$
\xymatrix{
\mathit{Mor}_{\mathcal{S}_1}(x, y)(f : V \to U) \ar@{=}[r] &
\Mor_{\mathcal{S}_{1, V}}(f^\ast x, f^\ast y) \ar[d] \\
\mathit{Mor}_{\mathcal{S}_2}(F(x), F(y))(f : V \to U) \ar@{=}[r] &
\Mor_{\mathcal{S}_{2, V}}(f^\ast F(x), f^\ast F(y))
}
$$
par $\phi \mapsto \beta_V^{-1} \circ F(\phi) \circ \alpha_V$.
Nous omettons la vérification de la compatibilité avec les applications de
restriction.
\end{proof}

\begin{remark}
\label{remark-alternative}
Supposons que $p : \mathcal{S} \to \mathcal{C}$ soit fibrée en groupoïdes.
Dans ce cas, nous pouvons démontrer le lemme \ref{lemma-painful}
à l'aide du lemme \ref{categories-lemma-fibred-strict} de Catégories,
qui dit que $\mathcal{S} \to \mathcal{C}$ est équivalente à la catégorie
associée à un foncteur contravariant
$F : \mathcal{C} \to \textit{Groupoids}$.
Dans le cas de la catégorie fibrée associée à $F$,
nous avons exactement $g^* \circ f^* = (f \circ g)^*$,
et il n'est pas nécessaire d'utiliser les applications $\alpha_{g, f}$.
Dans ce cas, le lemme est (encore plus) trivial. Bien entendu, on utilise alors
le fait que le préfaisceau $\mathit{Mor}(x, y)$ reste inchangé lorsque l'on
passe à une catégorie fibrée équivalente, ce qui découle du
lemme \ref{lemma-presheaf-mor-map-fibred-categories}.
\end{remark}

\begin{lemma}
\label{lemma-isom-as-2-fibre-product}
Soit $\mathcal{C}$ une catégorie.
Soit $p : \mathcal{S} \to \mathcal{C}$ une catégorie fibrée,
voir Catégories, section \ref{categories-section-fibred-categories}.
Soit $U \in \Ob(\mathcal{C})$ et soient $x, y \in \Ob(\mathcal{S}_U)$.
Notons encore $x, y : \mathcal{C}/U \to \mathcal{S}$ les
$1$-morphismes correspondants, voir
Catégories, lemme \ref{categories-lemma-yoneda-2category}.
Alors
\begin{enumerate}
\item le $2$-produit fibré
$\mathcal{S} \times_{\mathcal{S} \times \mathcal{S}, (x, y)} \mathcal{C}/U$
est fibré en setoïdes au-dessus de $\mathcal{C}/U$, et
\item $\mathit{Isom}(x, y)$ est le préfaisceau d'ensembles correspondant
à cette catégorie fibrée en setoïdes, voir
Catégories, lemme \ref{categories-lemma-2-category-fibred-setoids}.
\end{enumerate}
\end{lemma}

\begin{proof}
Omis. Indication : les objets du $2$-produit fibré sont
$(a : V \to U, z, (\alpha, \beta))$, où
$\alpha : z \to a^*x$ et $\beta : z \to a^*y$ sont des isomorphismes
dans $\mathcal{S}_V$. Le lien avec $\mathit{Isom}(x, y)$
s'obtient donc en associant à un tel objet l'isomorphisme
$\beta \circ \alpha^{-1}$.
\end{proof}







\section{Données de descente dans les catégories fibrées}
\label{section-descent-data}

\noindent
Dans cette section, nous définissons la notion de donnée de descente
dans le cadre abstrait d'une catégorie fibrée. Signalons auparavant qu'elle
est entièrement analogue aux données de descente pour les faisceaux
quasi-cohérents (Descente, section \ref{descent-section-equivalence})
et aux données de descente pour les schémas au-dessus de schémas
(Descente, section \ref{descent-section-descent-datum}).

\medskip\noindent
Nous adopterons la convention selon laquelle les projections
$\text{pr}_i : X \times \ldots \times X \to X$
sont numérotées à partir de $i = 0$. Ainsi, nous avons
$\text{pr}_0, \text{pr}_1 : X \times X  \to X$,
$\text{pr}_0, \text{pr}_1, \text{pr}_2 : X \times X \times X  \to X$,
etc.

\begin{definition}
\label{definition-descent-data}
Soit $\mathcal{C}$ une catégorie.
Soit $p : \mathcal{S} \to \mathcal{C}$ une catégorie fibrée.
Faisons un choix d'images inverses comme dans Catégories,
définition \ref{categories-definition-pullback-functor-fibred-category}.
Soit $\mathcal{U} = \{f_i : U_i \to U\}_{i \in I}$
une famille de morphismes de $\mathcal{C}$. Supposons que tous les produits
fibrés $U_i \times_U U_j$ et $U_i \times_U U_j \times_U U_k$ existent.
\begin{enumerate}
\item Une {\it donnée de descente $(X_i, \varphi_{ij})$ dans $\mathcal{S}$
relative à la famille $\{f_i : U_i \to U\}$} consiste en un objet $X_i$
de $\mathcal{S}_{U_i}$ pour chaque $i \in I$, et en un isomorphisme
$\varphi_{ij} : \text{pr}_0^*X_i \to \text{pr}_1^*X_j$
dans $\mathcal{S}_{U_i \times_U U_j}$ pour chaque paire $(i, j) \in I^2$,
tels que, pour tout triplet d'indices $(i, j, k) \in I^3$, le diagramme
$$
\xymatrix{
\text{pr}_0^*X_i \ar[rd]_{\text{pr}_{01}^*\varphi_{ij}}
\ar[rr]_{\text{pr}_{02}^*\varphi_{ik}} & &
\text{pr}_2^*X_k \\
& \text{pr}_1^*X_j \ar[ru]_{\text{pr}_{12}^*\varphi_{jk}} &
}
$$
dans la catégorie $\mathcal{S}_{U_i \times_U U_j \times_U U_k}$
soit commutatif. C'est ce qu'on appelle la {\it condition de cocycle}.
\item Un {\it morphisme $\psi : (X_i, \varphi_{ij}) \to
(X'_i, \varphi'_{ij})$ de données de descente} consiste en une famille
$\psi = (\psi_i)_{i\in I}$ de morphismes
$\psi_i : X_i \to X'_i$ dans $\mathcal{S}_{U_i}$
telle que tous les diagrammes
$$
\xymatrix{
\text{pr}_0^*X_i \ar[r]_{\varphi_{ij}} \ar[d]_{\text{pr}_0^*\psi_i}
& \text{pr}_1^*X_j \ar[d]^{\text{pr}_1^*\psi_j} \\
\text{pr}_0^*X'_i \ar[r]^{\varphi'_{ij}} &
\text{pr}_1^*X'_j \\
}
$$
dans les catégories $\mathcal{S}_{U_i \times_U U_j}$ soient commutatifs.
\item La catégorie des données de descente relatives à
$\mathcal{U}$ est notée $DD(\mathcal{U})$.
\end{enumerate}
\end{definition}

\noindent
Les produits fibrés $U_i \times_U U_j$ et
$U_i \times_U U_j \times_U U_k$ existent si chacun des morphismes
$f_i : U_i \to U$ est {\it représentable}, voir
Catégories, définition \ref{categories-definition-representable-morphism}.
Rappelons que, dans un site, l'une des conditions imposées à un recouvrement
$\{U_i \to U\}$ est que chacun des morphismes soit représentable, voir
Sites, définition \ref{sites-definition-site}, partie (3).
En fait, la définition ci-dessus présente surtout un intérêt lorsque
$\mathcal{C}$ est un site et que $\{U_i \to U\}$ est un recouvrement de
$\mathcal{C}$. Une donnée de descente n'est toutefois qu'une construction
abstraite que l'on peut définir comme ci-dessus. C'est utile : par exemple,
étant donnée une catégorie fibrée au-dessus de $\mathcal{C}$, on peut considérer
l'ensemble des familles relativement auxquelles les données de descente sont
effectives et tenter de les employer comme famille de recouvrements d'un site.

\begin{remarks}
\label{remarks-definition-descent-datum}
Deux remarques au sujet de la définition \ref{definition-descent-data}
s'imposent. Soit $p : \mathcal{S} \to \mathcal{C}$ une catégorie fibrée.
Soient $\{f_i : U_i \to U\}_{i \in I}$ et $(X_i, \varphi_{ij})$
comme dans la définition \ref{definition-descent-data}.
\begin{enumerate}
\item Il existe un morphisme diagonal $\Delta : U_i \to U_i \times_U U_i$.
Nous pouvons tirer $\varphi_{ii}$ en arrière par ce morphisme et obtenir un
automorphisme $\Delta^\ast \varphi_{ii} \in \text{Aut}_{U_i}(X_i)$.
En tirant la condition de cocycle pour le triplet $(i, i, i)$ en arrière par
$\Delta_{123} : U_i \to U_i \times_U U_i \times_U U_i$, nous déduisons que
$\Delta^\ast \varphi_{ii} \circ \Delta^\ast \varphi_{ii} =
\Delta^\ast \varphi_{ii}$ ; ainsi $\Delta^\ast \varphi_{ii} =
\text{id}_{X_i}$.
\item Il existe un morphisme
$\Delta_{13}: U_i \times_U U_j \to U_i \times_U U_j \times_U U_i$,
et nous pouvons tirer en arrière la condition de cocycle pour le triplet
$(i, j, i)$ afin d'obtenir l'identité
$(\sigma^\ast \varphi_{ji}) \circ \varphi_{ij} =
\text{id}_{\text{pr}_0^\ast X_i}$, où
$\sigma : U_i \times_U U_j \to U_j \times_U U_i$ est le morphisme de permutation.
\end{enumerate}
\end{remarks}

\begin{lemma}
\label{lemma-pullback}
(Changement de base des données de descente.)
Soit $\mathcal{C}$ une catégorie.
Soit $p : \mathcal{S} \to \mathcal{C}$ une catégorie fibrée.
Faisons un choix d'images inverses comme dans Catégories,
définition \ref{categories-definition-pullback-functor-fibred-category}.
Soient $\mathcal{U} = \{f_i : U_i \to U\}_{i \in I}$ et
$\mathcal{V} = \{V_j \to V\}_{j \in J}$
des familles de morphismes de $\mathcal{C}$ ayant chacune un but fixé.
Supposons que tous les produits fibrés
$U_i \times_U U_{i'}$, $U_i \times_U U_{i'} \times_U U_{i''}$,
$V_j \times_V V_{j'}$ et $V_j \times_V V_{j'} \times_V V_{j''}$ existent.
Les données $\alpha : I \to J$, $h : U \to V$ et
$g_i : U_i \to V_{\alpha(i)}$ définissent un morphisme de familles
d'applications à but fixé, voir
Sites, définition \ref{sites-definition-morphism-coverings}.
\begin{enumerate}
\item Soit $(Y_j, \varphi_{jj'})$ une donnée de descente relative à la
famille $\{V_j \to V\}$. Le système
$$
\left(
g_i^*Y_{\alpha(i)},
(g_i \times g_{i'})^*\varphi_{\alpha(i)\alpha(i')}
\right)
$$
est une donnée de descente relative à $\mathcal{U}$.
\item Cette construction définit un foncteur des données de descente relatives
à $\mathcal{V}$ vers les données de descente relatives à $\mathcal{U}$.
\item Étant donné un second morphisme de familles d'applications à but fixé
formé de $\alpha' : I \to J$, $h' : U \to V$ et
$g'_i : U_i \to V_{\alpha'(i)}$, si $h = h'$, alors les deux foncteurs
obtenus entre catégories de données de descente sont canoniquement isomorphes.
\end{enumerate}
\end{lemma}

\begin{proof}
Omis.
\end{proof}

\begin{definition}
\label{definition-pullback-functor}
Avec $\mathcal{U} = \{U_i \to U\}_{i \in I}$,
$\mathcal{V} = \{V_j \to V\}_{j \in J}$,
$\alpha : I \to J$, $h : U \to V$ et
$g_i : U_i \to V_{\alpha(i)}$ comme dans le lemme \ref{lemma-pullback},
le foncteur
$$
(Y_j, \varphi_{jj'}) \longmapsto
(g_i^*Y_{\alpha(i)}, (g_i \times g_{i'})^*\varphi_{\alpha(i)\alpha(i')})
$$
construit dans ce lemme est appelé le {\it foncteur de changement de base}
sur les données de descente.
\end{definition}

\noindent
Étant donné $h : U \to V$, s'il existe un morphisme
$\tilde h : \mathcal{U} \to \mathcal{V}$ au-dessus de $h$,
alors $\tilde h^*$ ne dépend pas du choix de $\tilde h$, comme nous l'avons vu
dans le lemme \ref{lemma-pullback}. Nous écrirons donc parfois simplement
$h^*$ pour désigner le foncteur de changement de base.

\begin{definition}
\label{definition-effective-descent-datum}
Soit $\mathcal{C}$ une catégorie.
Soit $p : \mathcal{S} \to \mathcal{C}$ une catégorie fibrée.
Faisons un choix d'images inverses comme dans Catégories,
définition \ref{categories-definition-pullback-functor-fibred-category}.
Soit $\mathcal{U} = \{f_i : U_i \to U\}_{i \in I}$ une famille de morphismes
de but $U$. Supposons que tous les produits fibrés
$U_i \times_U U_j$ et $U_i \times_U U_j \times_U U_k$ existent.
\begin{enumerate}
\item Étant donné un objet $X$ de $\mathcal{S}_U$, la {\it donnée de descente
triviale} est la donnée de descente $(X, \text{id}_X)$ relative à la famille
$\{\text{id}_U : U \to U\}$.
\item Étant donné un objet $X$ de $\mathcal{S}_U$, nous obtenons une
{\it donnée de descente canonique} sur la famille des objets $f_i^*X$ en
tirant en arrière la donnée de descente triviale $(X, \text{id}_X)$ par
l'application évidente
$\{f_i : U_i \to U\} \to \{\text{id}_U : U \to U\}$.
Nous notons cette donnée de descente $(f_i^*X, can)$.
\item Une donnée de descente $(X_i, \varphi_{ij})$
relative à $\{f_i : U_i \to U\}$ est dite {\it effective}
s'il existe un objet $X$ de $\mathcal{S}_U$ tel que
$(X_i, \varphi_{ij})$ soit isomorphe à $(f_i^*X, can)$.
\end{enumerate}
\end{definition}

\noindent
Remarquons que la règle qui associe à $X \in \mathcal{S}_U$ sa donnée de
descente canonique relative à $\mathcal{U}$ définit un foncteur
$$
\mathcal{S}_U \longrightarrow DD(\mathcal{U}).
$$
Une donnée de descente est effective si et seulement si elle appartient à
l'image essentielle de ce foncteur. Explicitons maintenant la donnée de
descente canonique.

\begin{lemma}
\label{lemma-trivial-cocycle}
Dans la situation de la partie (2) de la
définition \ref{definition-effective-descent-datum}, les applications
$can_{ij} : \text{pr}_0^*f_i^*X \to \text{pr}_1^*f_j^*X$ sont égales à
$(\alpha_{\text{pr}_1, f_j})_X \circ (\alpha_{\text{pr}_0, f_i})_X^{-1}$,
où $\alpha_{\cdot, \cdot}$ est comme dans le
lemme \ref{categories-lemma-fibred} de Catégories et où nous utilisons
l'égalité $f_i \circ \text{pr}_0 = f_j \circ \text{pr}_1$
des applications $U_i \times_U U_j \to U$.
\end{lemma}

\begin{proof}
Omis.
\end{proof}

\begin{lemma}
\label{lemma-compare-descent-condition}
Soit $\mathcal{C}$ une catégorie. Soit
$\mathcal{V} = \{V_j \to U\}_{j \in J} \to
\mathcal{U} = \{U_i \to U\}_{i \in I}$
un morphisme de familles d'applications à but fixé de $\mathcal{C}$, donné par
$\text{id} : U \to U$, $\alpha : J \to I$ et
$f_j : V_j \to U_{\alpha(j)}$. Soit
$p : \mathcal{S} \to \mathcal{C}$ une catégorie fibrée. Si
\begin{enumerate}
\item pour $0 \leq p \leq 3$ et $0 \leq q \leq 3$, avec $p + q \geq 2$,
et pour $i_1, \ldots, i_p \in I$ et $j_1, \ldots, j_q \in J$,
les produits fibrés $U_{i_1} \times_U \ldots \times_U U_{i_p} \times_U
V_{j_1} \times_U \ldots \times_U V_{j_q}$ existent,
\item le foncteur $\mathcal{S}_U \to DD(\mathcal{V})$
est une équivalence,
\item pour tout $i \in I$, le foncteur
$\mathcal{S}_{U_i} \to DD(\mathcal{V}_i)$
est pleinement fidèle, et
\item pour tous $i, i' \in I$, le foncteur
$\mathcal{S}_{U_i \times_U U_{i'}} \to DD(\mathcal{V}_{ii'})$
est fidèle.
\end{enumerate}
Ici $\mathcal{V}_i = \{U_i \times_U V_j \to U_i\}_{j \in J}$ et
$\mathcal{V}_{ii'} =
\{U_i \times_U U_{i'} \times_U V_j \to U_i \times_U U_{i'}\}_{j \in J}$.
Alors $\mathcal{S}_U \to DD(\mathcal{U})$ est une équivalence.
\end{lemma}

\begin{proof}
La condition (1) garantit l'existence d'un nombre suffisant de produits fibrés
pour que l'énoncé ait un sens. Nous allons montrer que le foncteur
$\mathcal{S}_U \to DD(\mathcal{U})$ est essentiellement surjectif.
Supposons donnée une donnée de descente $(X_i, \varphi_{ii'})$ relative à
$\mathcal{U}$. D'après le lemme \ref{lemma-pullback}, nous pouvons la tirer en
arrière en une donnée de descente $(X_j, \varphi_{jj'})$ pour $\mathcal{V}$.
D'après l'hypothèse (2), cette donnée de descente est effective ; nous obtenons
donc un objet $X$ de $\mathcal{S}_U$ tel que l'image inverse de la donnée de
descente triviale $(X, \text{id}_X)$ par le morphisme
$\mathcal{V} \to \{U \to U\}$ soit isomorphe à
$(X_j, \varphi_{jj'})$. Observons ensuite que nous avons un diagramme
$$
\xymatrix{
\mathcal{V}_i \ar[r] \ar[d] &
\mathcal{V} \ar[r] &
\mathcal{U} \ar[d] \\
\{U_i \to U_i\} \ar[rr] \ar[rru] & &
\{U \to U\}
}
$$
de morphismes de familles d'applications à but fixé de $\mathcal{C}$.
Ce diagramme n'est pas commutatif, mais, d'après le
lemme \ref{lemma-pullback}, les foncteurs de changement de base ainsi obtenus
sur les données de descente sont canoniquement isomorphes. Par conséquent,
$(X, \text{id}_X)$ et $(X_i, \text{id}_{X_i})$ ont des images inverses
isomorphes dans $DD(\mathcal{V}_i)$. L'hypothèse (3) fournit donc un
isomorphisme $(U_i \to U)^*X \to X_i$ dans la catégorie
$\mathcal{S}_{U_i}$. Nous omettons de vérifier que ces flèches sont compatibles
avec les morphismes $\varphi_{ii'}$ ; indication : utiliser la fidélité des
foncteurs de la condition (4). Nous omettons également de vérifier que le
foncteur $\mathcal{S}_U \to DD(\mathcal{U})$ est pleinement fidèle.
\end{proof}











\section{Champs}
\label{section-definition}

\noindent
Voici la définition d'un champ. Elle combine la notion de catégorie
fibrée et celle de descente.

\begin{definition}
\label{definition-stack}
Soit $\mathcal{C}$ un site. Un {\it champ} sur $\mathcal{C}$
est une catégorie $p : \mathcal{S} \to \mathcal{C}$ au-dessus de $\mathcal{C}$ qui
satisfait aux conditions suivantes :
\begin{enumerate}
\item $p : \mathcal{S} \to \mathcal{C}$ est une catégorie fibrée, voir
Catégories, définition \ref{categories-definition-fibred-category},
\item pour tout $U \in \Ob(\mathcal{C})$ et tous $x, y \in \mathcal{S}_U$,
le préfaisceau $\mathit{Mor}(x, y)$ (voir
Définition \ref{definition-mor-presheaf}) est un faisceau sur
le site $\mathcal{C}/U$, et
\item pour tout recouvrement $\mathcal{U} = \{f_i : U_i \to U\}_{i \in I}$
du site $\mathcal{C}$, toute donnée de descente dans $\mathcal{S}$
relativement à $\mathcal{U}$ est effective.
\end{enumerate}
\end{definition}

\noindent
La formulation ci-dessus nous paraît être la façon la plus commode de concevoir
un champ. En effet, étant donné une catégorie au-dessus de $\mathcal{C}$, pour vérifier
qu'il s'agit d'un champ, on vérifie les propriétés (1), (2) et
(3), dans cet ordre. Bien entendu, les propriétés (2) et (3) n'ont pas de sens
si la catégorie n'est pas fibrée. Sans (2), nous ne pouvons pas prouver que la
descente de (3) est unique à un unique isomorphisme près et fonctorielle.

\medskip\noindent
Le lemme suivant fournit une autre définition.

\begin{lemma}
\label{lemma-stack-equivalences}
Soit $\mathcal{C}$ un site.
Soit $p : \mathcal{S} \to \mathcal{C}$ une catégorie fibrée
au-dessus de $\mathcal{C}$. Les assertions suivantes sont équivalentes :
\begin{enumerate}
\item $\mathcal{S}$ est un champ sur $\mathcal{C}$, et
\item pour tout recouvrement $\mathcal{U} = \{f_i : U_i \to U\}_{i \in I}$
du site $\mathcal{C}$, le foncteur
$$
\mathcal{S}_U \longrightarrow DD(\mathcal{U})
$$
qui associe à un objet sa donnée de descente canonique est une équivalence.
\end{enumerate}
\end{lemma}

\begin{proof}
Démonstration omise.
\end{proof}

\begin{lemma}
\label{lemma-substack}
Soit $p : \mathcal{S} \to \mathcal{C}$ un champ sur le site $\mathcal{C}$.
Soit $\mathcal{S}'$ une sous-catégorie de $\mathcal{S}$.
Supposons que
\begin{enumerate}
\item si $\varphi : y \to x$ est un morphisme fortement cartésien
de $\mathcal{S}$ et si
$x$ est un objet de $\mathcal{S}'$, alors $y$ est isomorphe à un
objet de $\mathcal{S}'$,
\item $\mathcal{S}'$ est une sous-catégorie pleine de $\mathcal{S}$, et
\item si $\{f_i : U_i \to U\}$ est un recouvrement de $\mathcal{C}$
et si $x$ est un objet de $\mathcal{S}$ au-dessus de $U$ tel que $f_i^*x$
soit isomorphe à un objet de $\mathcal{S}'$ pour tout $i$,
alors $x$ est isomorphe à un objet de $\mathcal{S}'$.
\end{enumerate}
Alors $\mathcal{S}' \to \mathcal{C}$ est un champ.
\end{lemma}

\begin{proof}
Démonstration omise. Indications :
la première condition garantit que $\mathcal{S}'$ est une catégorie fibrée.
La deuxième garantit que les préfaisceaux $\mathit{Isom}$
de $\mathcal{S}'$ sont des faisceaux (puisqu'ils sont identiques à leurs homologues
dans $\mathcal{S}$). La troisième garantit que la condition de descente
est satisfaite dans $\mathcal{S}'$, car nous pouvons d'abord effectuer la descente dans
$\mathcal{S}$, puis (3) implique que l'objet obtenu est isomorphe à un objet de
$\mathcal{S}'$.
\end{proof}

\begin{lemma}
\label{lemma-stack-equivalent}
Soit $\mathcal{C}$ un site.
Soient $\mathcal{S}_1$, $\mathcal{S}_2$ des catégories au-dessus de $\mathcal{C}$.
Supposons que $\mathcal{S}_1$ et $\mathcal{S}_2$ soient équivalentes
comme catégories au-dessus de $\mathcal{C}$.
Alors $\mathcal{S}_1$ est un champ sur $\mathcal{C}$ si et seulement si
$\mathcal{S}_2$ est un champ sur $\mathcal{C}$.
\end{lemma}

\begin{proof}
Soient $F : \mathcal{S}_1 \to \mathcal{S}_2$,
$G : \mathcal{S}_2 \to \mathcal{S}_1$ des foncteurs au-dessus de $\mathcal{C}$, et soient
$i : F \circ G \to \text{id}_{\mathcal{S}_2}$,
$j : G \circ F \to \text{id}_{\mathcal{S}_1}$ des isomorphismes de
foncteurs au-dessus de $\mathcal{C}$. D'après
Catégories, lemme \ref{categories-lemma-fibred-equivalent},
nous voyons que $\mathcal{S}_1$ est fibrée si et seulement si $\mathcal{S}_2$
est fibrée au-dessus de $\mathcal{C}$. Nous pouvons donc supposer que
$\mathcal{S}_1$ et $\mathcal{S}_2$ sont toutes deux fibrées. De plus, la démonstration de
Catégories, lemme \ref{categories-lemma-fibred-equivalent},
montre que $F$ et $G$ envoient les morphismes fortement cartésiens sur des morphismes
fortement cartésiens, c'est-à-dire que $F$ et $G$ sont des $1$-morphismes de catégories
fibrées au-dessus de $\mathcal{C}$. Cela signifie que, étant donnés
$U \in \Ob(\mathcal{C})$ et $x, y \in \mathcal{S}_{1, U}$, les préfaisceaux
$$
\mathit{Mor}_{\mathcal{S}_1}(x, y),
\mathit{Mor}_{\mathcal{S}_1}(F(x), F(y)) :
(\mathcal{C}/U)^{opp} \longrightarrow \textit{Ensembles}.
$$
s'identifient, voir
lemme \ref{lemma-presheaf-mor-map-fibred-categories}. Par conséquent,
le premier est un faisceau si et seulement si le second est un faisceau.
Enfin, nous devons montrer que, si toute donnée de descente dans $\mathcal{S}_1$
est effective, il en va de même de toute donnée de descente dans $\mathcal{S}_2$.
Pour cela, soit $(X_i, \varphi_{ii'})$ une donnée de descente
dans $\mathcal{S}_2$ relativement au recouvrement $\{U_i \to U\}$ du site
$\mathcal{C}$. Alors $(G(X_i), G(\varphi_{ii'}))$ est une donnée de descente
dans $\mathcal{S}_1$ relativement au recouvrement $\{U_i \to U\}$.
Soit $X$ un objet de $\mathcal{S}_{1, U}$ tel que la
donnée de descente $(f_i^*X, can)$ soit isomorphe à
$(G(X_i), G(\varphi_{ii'}))$. Alors $F(X)$ est un objet de $\mathcal{S}_{2, U}$
tel que la donnée de descente $(f_i^*F(X), can)$ soit isomorphe à
$(F(G(X_i)), F(G(\varphi_{ii'})))$, laquelle est à son tour isomorphe à
la donnée de descente initiale $(X_i, \varphi_{ii'})$ au moyen de $i$.
\end{proof}

\noindent
La $2$-catégorie des champs sur $\mathcal{C}$
est définie comme suit.

\begin{definition}
\label{definition-stacks-over-C}
Soit $\mathcal{C}$ un site.
La {\it $2$-catégorie des champs sur $\mathcal{C}$}
est la sous-$2$-catégorie de la $2$-catégorie des catégories fibrées
au-dessus de $\mathcal{C}$ (voir
Catégories, définition \ref{categories-definition-fibred-categories-over-C})
définie comme suit :
\begin{enumerate}
\item Ses objets sont les champs $p : \mathcal{S} \to \mathcal{C}$.
\item Ses $1$-morphismes $(\mathcal{S}, p) \to (\mathcal{S}', p')$
sont les foncteurs $G : \mathcal{S} \to \mathcal{S}'$ tels que
$p' \circ G = p$ et tels que $G$ envoie les morphismes fortement cartésiens
sur des morphismes fortement cartésiens.
\item Ses $2$-morphismes $t : G \to H$, pour
$G, H : (\mathcal{S}, p) \to (\mathcal{S}', p')$,
sont les morphismes de foncteurs
tels que $p'(t_x) = \text{id}_{p(x)}$
pour tout $x \in \Ob(\mathcal{S})$.
\end{enumerate}
\end{definition}

\begin{lemma}
\label{lemma-2-product-stacks}
Soit $\mathcal{C}$ un site.
La $(2, 1)$-catégorie des champs sur $\mathcal{C}$
admet des 2-produits fibrés, et ceux-ci sont décrits comme dans
Catégories, lemme \ref{categories-lemma-2-product-categories-over-C}.
\end{lemma}

\begin{proof}
Soient $f : \mathcal{X} \to \mathcal{S}$ et
$g : \mathcal{Y} \to \mathcal{S}$ des
$1$-morphismes de champs sur $\mathcal{C}$
au sens de la définition ci-dessus. La catégorie
$\mathcal{X} \times_\mathcal{S} \mathcal{Y}$
décrite dans
Catégories, lemme \ref{categories-lemma-2-product-categories-over-C}, est une
catégorie fibrée d'après
Catégories, lemme \ref{categories-lemma-2-product-fibred-categories-over-C}.
(C'est ici que nous utilisons le fait que $f$ et $g$ préservent les morphismes
fortement cartésiens.) Il reste à montrer que les préfaisceaux de morphismes sont des faisceaux
et que la descente relativement aux recouvrements de $\mathcal{C}$ est effective.

\medskip\noindent
Rappelons qu'un objet de $\mathcal{X} \times_\mathcal{S} \mathcal{Y}$
est donné par un quadruplet $(U, x, y, \phi)$.
Il se trouve au-dessus de l'objet
$U$ de $\mathcal{C}$. Soit ensuite $(U, x', y', \phi')$ un second
objet au-dessus de $U$.
Rappelons que $\phi : f(x) \to g(y)$ et $\phi' : f(x') \to g(y')$
sont des isomorphismes dans la catégorie $\mathcal{S}_U$. Utilisons ces
isomorphismes pour identifier $z = f(x) = g(y)$ et
$z' = f(x') = g(y')$. Avec ces identifications,
il est clair que
$$
\mathit{Mor}((U, x, y, \phi), (U, x', y', \phi'))
=
\mathit{Mor}(x, x')
\times_{\mathit{Mor}(z, z')}
\mathit{Mor}(y, y')
$$
comme préfaisceaux. Or, puisque le produit fibré dans la catégorie des
préfaisceaux préserve les faisceaux (Sites, lemme \ref{sites-lemma-limit-sheaf}),
nous voyons que c'est un faisceau.

\medskip\noindent
Soit $\mathcal{U} = \{f_i : U_i \to U\}_{i \in I}$ un recouvrement du site
$\mathcal{C}$. Soit $(X_i, \chi_{ij})$ une donnée de descente
dans $\mathcal{X} \times_\mathcal{S} \mathcal{Y}$ relativement à $\mathcal{U}$.
Écrivons $X_i = (U_i, x_i, y_i, \phi_i)$ comme ci-dessus. Écrivons
$\chi_{ij} = (\varphi_{ij}, \psi_{ij})$ comme dans la définition de
la catégorie $\mathcal{X} \times_\mathcal{S} \mathcal{Y}$ (voir
Catégories, lemme \ref{categories-lemma-2-product-categories-over-C}).
Il est clair que $(x_i, \varphi_{ij})$ est une donnée de descente dans
$\mathcal{X}$ et que $(y_i, \psi_{ij})$ est une donnée de descente dans
$\mathcal{Y}$. Puisque $\mathcal{X}$ et $\mathcal{Y}$ sont des champs, ces
données de descente sont effectives. Nous obtenons donc
$x \in \Ob(\mathcal{X}_U)$ et $y \in \Ob(\mathcal{Y}_U)$
avec $x_i = x|_{U_i}$ et $y_i = y|_{U_i}$, de façon compatible avec les données de descente.
Posons $z = f(x)$ et $z' = g(y)$, qui sont tous deux des objets de $\mathcal{S}_U$.
Les morphismes $\phi_i$ sont des éléments de
$\mathit{Isom}(z, z')(U_i)$ vérifiant
$\phi_i|_{U_i \times_U U_j} = \phi_j|_{U_i \times_U U_j}$.
Par la propriété de faisceau de $\mathit{Isom}(z, z')$,
nous obtenons donc un isomorphisme $\phi : z = f(x) \to z' = g(y)$.
Nous omettons la vérification que la donnée de descente canonique associée à
l'objet $(U, x, y, \phi)$ de
$(\mathcal{X} \times_\mathcal{S} \mathcal{Y})_U$ est isomorphe
à la donnée de descente dont nous sommes partis.
\end{proof}

\begin{lemma}
\label{lemma-characterize-ff}
Soit $\mathcal{C}$ un site.
Soient $\mathcal{S}_1$, $\mathcal{S}_2$ des champs sur $\mathcal{C}$.
Soit $F : \mathcal{S}_1 \to \mathcal{S}_2$ un $1$-morphisme.
Alors les assertions suivantes sont équivalentes :
\begin{enumerate}
\item $F$ est pleinement fidèle,
\item pour tout $U \in \Ob(\mathcal{C})$ et tous
$x, y \in \Ob(\mathcal{S}_{1, U})$, l'application
$$
F :
\mathit{Mor}_{\mathcal{S}_1}(x, y)
\longrightarrow
\mathit{Mor}_{\mathcal{S}_2}(F(x), F(y))
$$
est un isomorphisme de faisceaux sur $\mathcal{C}/U$.
\end{enumerate}
\end{lemma}

\begin{proof}
Supposons (1). Pour $U, x, y$ comme dans (2), l'application affichée $F$ s'évalue en
$F : \Mor_{\mathcal{S}_{1, V}}(x|_V, y|_V) \to
\Mor_{\mathcal{S}_{2, V}}(F(x|_V), F(y|_V))$
sur un objet $V$ de $\mathcal{C}$ au-dessus de $U$.
Or, puisque $F$ est pleinement fidèle, l'application correspondante
$\Mor_{\mathcal{S}_1}(x|_V, y|_V) \to \Mor_{\mathcal{S}_2}(F(x|_V), F(y|_V))$
est une bijection. Les morphismes de la catégorie fibre $\mathcal{S}_{1, V}$ sont
exactement les morphismes entre $x|_V$ et $y|_V$ dans $\mathcal{S}_1$ qui sont
au-dessus de $\text{id}_V$. De même, les morphismes de la catégorie fibre
$\mathcal{S}_{2, V}$ sont exactement les morphismes entre $F(x|_V)$ et
$F(y|_V)$ dans $\mathcal{S}_2$ qui sont au-dessus de $\text{id}_V$. Nous constatons donc que $F$
induit également une bijection entre ceux-ci. Par conséquent, (2) est satisfaite.

\medskip\noindent
Supposons (2). Soient $U$, $V$ des objets de $\mathcal{C}$, ainsi que
$x \in \Ob(\mathcal{S}_{1, U})$ et
$y \in \Ob(\mathcal{S}_{1, V})$. Pour montrer que $F$ est pleinement fidèle,
il suffit de prouver qu'il induit une bijection sur les
morphismes au-dessus d'un $f : U \to V$ fixé. Choisissons un morphisme fortement cartésien
$f^*y \to y$ dans $\mathcal{S}_1$ au-dessus de $f$. Nous obtenons ainsi une
bijection entre l'ensemble des morphismes $x \to y$ dans $\mathcal{S}_1$ au-dessus
de $f$ et $\Mor_{\mathcal{S}_{1, U}}(x, f^*y)$. Puisque $F$ préserve
les morphismes fortement cartésiens en tant que $1$-morphisme dans la $2$-catégorie
des champs sur $\mathcal{C}$, nous obtenons également une bijection
entre l'ensemble des morphismes $F(x) \to F(y)$ dans $\mathcal{S}_2$ au-dessus
de $f$ et $\Mor_{\mathcal{S}_{2, U}}(F(x), F(f^*y))$.
Comme $F$ induit une bijection
$\Mor_{\mathcal{S}_{1, U}}(x, f^*y) \to
\Mor_{\mathcal{S}_{2, U}}(F(x), F(f^*y))$,
nous concluons que (1) est satisfaite.
\end{proof}

\begin{lemma}
\label{lemma-characterize-essentially-surjective-when-ff}
Soit $\mathcal{C}$ un site.
Soient $\mathcal{S}_1$, $\mathcal{S}_2$ des champs sur $\mathcal{C}$.
Soit $F : \mathcal{S}_1 \to \mathcal{S}_2$ un $1$-morphisme
pleinement fidèle. Alors les assertions suivantes sont équivalentes :
\begin{enumerate}
\item $F$ est une équivalence,
\item pour tout $U \in \Ob(\mathcal{C})$ et tout
$x \in \Ob(\mathcal{S}_{2, U})$, il existe un recouvrement
$\{f_i : U_i \to U\}$ tel que $f_i^*x$ appartienne à l'image essentielle
du foncteur $F : \mathcal{S}_{1, U_i} \to \mathcal{S}_{2, U_i}$.
\end{enumerate}
\end{lemma}

\begin{proof}
L'implication (1) $\Rightarrow$ (2) est immédiate.
Pour voir que (2) implique (1), nous devons montrer que tout
$x$ comme dans (2) appartient à l'image essentielle du foncteur $F$.
Pour cela, choisissons un recouvrement comme dans (2),
$x_i \in \Ob(\mathcal{S}_{1, U_i})$, et des
isomorphismes $\varphi_i : F(x_i) \to f_i^*x$. Nous obtenons alors une donnée de
descente pour $\mathcal{S}_1$ relativement à $\{f_i : U_i \to U\}$
en prenant pour
$$
\varphi_{ij} :
x_i|_{U_i \times_U U_j}
\longrightarrow
x_j|_{U_i \times_U U_j}
$$
la flèche telle que $F(\varphi_{ij}) = \varphi_j^{-1} \circ \varphi_i$.
Cette donnée de descente est effective par les axiomes d'un champ ; nous
obtenons donc un objet $x_1$ de $\mathcal{S}_1$ au-dessus de $U$. Nous omettons la
vérification que $F(x_1)$ est isomorphe à $x$ au-dessus de $U$.
\end{proof}

\begin{remark}
\label{remark-stack-make-small}
(Réduction d'un « grand » champ pour obtenir un champ.)
Soit $\mathcal{C}$ un site. Supposons que $p : \mathcal{S} \to \mathcal{C}$
soit un foncteur d'une « grande » catégorie vers $\mathcal{C}$, c'est-à-dire supposons
que la collection des objets de $\mathcal{S}$ forme une classe propre.
Supposons enfin que $p : \mathcal{S} \to \mathcal{C}$ satisfasse
aux conditions (1), (2), (3) de la
Définition \ref{definition-stack}.
En général, il n'existe aucun moyen de remplacer $p : \mathcal{S} \to \mathcal{C}$
par une catégorie équivalente de façon à obtenir un champ. La raison en est qu'il
peut arriver qu'une catégorie fibre $\mathcal{S}_U$ ait une classe propre
de classes d'isomorphisme d'objets.
Supposons en revanche que
\begin{enumerate}
\item[(4)] pour tout $U \in \Ob(\mathcal{C})$, il existe un ensemble
$S_U \subset \Ob(\mathcal{S}_U)$ tel que tout objet de
$\mathcal{S}_U$ soit isomorphe dans $\mathcal{S}_U$ à un élément de $S_U$.
\end{enumerate}
Dans ce cas, nous pouvons trouver une sous-catégorie pleine $\mathcal{S}_{small}$
de $\mathcal{S}$ telle qu'en posant $p_{small} = p|_{\mathcal{S}_{small}}$,
nous ayons
\begin{enumerate}
\item[(a)] le foncteur $p_{small} : \mathcal{S}_{small} \to \mathcal{C}$
définit un champ, et
\item[(b)] l'inclusion $\mathcal{S}_{small} \to \mathcal{S}$
est pleinement fidèle et essentiellement surjective.
\end{enumerate}
(Indication : pour tout $U \in \Ob(\mathcal{C})$,
notons $\alpha(U)$ le plus petit ordinal tel que
$\Ob(\mathcal{S}_U) \cap V_{\alpha(U)}$ s'envoie surjectivement sur l'ensemble
des classes d'isomorphisme de $\mathcal{S}_U$, et posons
$\alpha = \sup_{U \in \Ob(\mathcal{C})} \alpha(U)$.
Prenons alors
$\Ob(\mathcal{S}_{small}) = \Ob(\mathcal{S}) \cap V_\alpha$.
Pour les notations utilisées, voir Ensembles, section \ref{sets-section-sets-hierarchy}.)
\end{remark}







\section{Champs en groupoïdes}
\label{section-stacks-in-groupoids}

\noindent
Parmi les champs, ceux qui sont fibrés en groupoïdes sont un peu plus faciles
à appréhender. Nous les redéfinissons comme suit.

\begin{definition}
\label{definition-stack-in-groupoids}
Un {\it champ en groupoïdes} sur un site $\mathcal{C}$ est une
catégorie $p : \mathcal{S} \to \mathcal{C}$ au-dessus de $\mathcal{C}$
telle que
\begin{enumerate}
\item $p : \mathcal{S} \to \mathcal{C}$ est fibrée
en groupoïdes au-dessus de $\mathcal{C}$ (voir
Catégories, définition \ref{categories-definition-fibred-groupoids}),
\item pour tout $U \in \Ob(\mathcal{C})$ et
tous $x, y\in \Ob(\mathcal{S}_U)$, le préfaisceau
$\mathit{Isom}(x, y)$ est un faisceau sur le site $\mathcal{C}/U$, et
\item pour tout recouvrement $\mathcal{U} = \{U_i \to U\}$ dans $\mathcal{C}$,
toute donnée de descente $(x_i, \phi_{ij})$ relative à $\mathcal{U}$ est effective.
\end{enumerate}
\end{definition}

\noindent
La troisième condition est généralement la plus difficile à vérifier.
Voici le lemme qui compare cette notion à celle de champ.

\begin{lemma}
\label{lemma-stack-in-groupoids-stack}
Soit $\mathcal{C}$ un site.
Soit $p : \mathcal{S} \to \mathcal{C}$ une catégorie au-dessus de $\mathcal{C}$.
Les assertions suivantes sont équivalentes :
\begin{enumerate}
\item $\mathcal{S}$ est un champ en groupoïdes sur $\mathcal{C}$,
\item $\mathcal{S}$ est un champ sur $\mathcal{C}$ et toutes ses
catégories fibres sont des groupoïdes, et
\item $\mathcal{S}$ est fibrée en groupoïdes au-dessus de $\mathcal{C}$
et est un champ sur $\mathcal{C}$.
\end{enumerate}
\end{lemma}

\begin{proof}
Démonstration omise, mais voir Catégories, lemme \ref{categories-lemma-fibred-groupoids}.
\end{proof}

\begin{lemma}
\label{lemma-stack-gives-stack-groupoids}
Soit $\mathcal{C}$ un site.
Soit $p : \mathcal{S} \to \mathcal{C}$ un champ.
Soit $p' : \mathcal{S}' \to \mathcal{C}$
la catégorie fibrée en groupoïdes associée à $\mathcal{S}$,
construite dans
Catégories, lemme \ref{categories-lemma-fibred-gives-fibred-groupoids}.
Alors $p' : \mathcal{S}' \to \mathcal{C}$ est un champ en groupoïdes.
\end{lemma}

\begin{proof}
Rappelons que les morphismes de $\mathcal{S}'$ sont exactement les
morphismes fortement cartésiens de $\mathcal{S}$ et que tout isomorphisme de
$\mathcal{S}$ est un tel morphisme. Ainsi, les données de descente dans $\mathcal{S}'$
sont exactement les mêmes que les données de descente dans $\mathcal{S}$. Appliquons maintenant le
lemme \ref{lemma-stack-equivalences}. Certains détails sont omis.
\end{proof}

\begin{lemma}
\label{lemma-stack-in-groupoids-equivalent}
Soit $\mathcal{C}$ un site.
Soient $\mathcal{S}_1$, $\mathcal{S}_2$ des catégories au-dessus de $\mathcal{C}$.
Supposons que $\mathcal{S}_1$ et $\mathcal{S}_2$ soient équivalentes
comme catégories au-dessus de $\mathcal{C}$.
Alors $\mathcal{S}_1$ est un champ en groupoïdes sur $\mathcal{C}$ si et seulement si
$\mathcal{S}_2$ est un champ en groupoïdes sur $\mathcal{C}$.
\end{lemma}

\begin{proof}
Cela résulte de la combinaison des
lemmes \ref{lemma-stack-in-groupoids-stack} et \ref{lemma-stack-equivalent}.
\end{proof}

\noindent
La $2$-catégorie des champs en groupoïdes sur $\mathcal{C}$
est définie comme suit.

\begin{definition}
\label{definition-stacks-in-groupoids-over-C}
Soit $\mathcal{C}$ un site.
La {\it $2$-catégorie des champs en groupoïdes sur $\mathcal{C}$}
est la sous-$2$-catégorie de la $2$-catégorie des champs
sur $\mathcal{C}$ (voir Définition \ref{definition-stacks-over-C})
définie comme suit :
\begin{enumerate}
\item Ses objets sont les champs en groupoïdes
$p : \mathcal{S} \to \mathcal{C}$.
\item Ses $1$-morphismes $(\mathcal{S}, p) \to (\mathcal{S}', p')$
sont les foncteurs $G : \mathcal{S} \to \mathcal{S}'$ tels que
$p' \circ G = p$. (Puisque tout morphisme est fortement cartésien,
tout foncteur les préserve.)
\item Ses $2$-morphismes $t : G \to H$, pour
$G, H : (\mathcal{S}, p) \to (\mathcal{S}', p')$,
sont les morphismes de foncteurs
tels que $p'(t_x) = \text{id}_{p(x)}$
pour tout $x \in \Ob(\mathcal{S})$.
\end{enumerate}
\end{definition}

\noindent
Remarquons que tout $2$-morphisme est automatiquement un isomorphisme, de sorte
que la $2$-catégorie des champs en groupoïdes sur $\mathcal{C}$
est en fait une $(2, 1)$-catégorie (stricte).

\begin{lemma}
\label{lemma-2-product-stacks-in-groupoids}
Soit $\mathcal{C}$ une catégorie.
La $2$-catégorie des champs en groupoïdes sur $\mathcal{C}$
admet des 2-produits fibrés, et ceux-ci sont décrits comme dans
Catégories, lemme \ref{categories-lemma-2-product-categories-over-C}.
\end{lemma}

\begin{proof}
Cela résulte immédiatement de
Catégories, lemme \ref{categories-lemma-2-product-fibred-categories},
et des lemmes \ref{lemma-stack-in-groupoids-stack}
et \ref{lemma-2-product-stacks}.
\end{proof}




\section{Champs en setoïdes}
\label{section-stacks-in-setoids}

\noindent
Cette brève section se contente d'indiquer qu'un champ en ensembles
est la même chose qu'un faisceau d'ensembles. Pour les notations, consulter
Catégories, section \ref{categories-section-fibred-in-setoids}.

\begin{definition}
\label{definition-stack-in-sets}
Soit $\mathcal{C}$ un site.
\begin{enumerate}
\item Un {\it champ en setoïdes} sur $\mathcal{C}$
est un champ sur $\mathcal{C}$ dont toutes les catégories fibres sont des
setoïdes.
\item Un {\it champ en ensembles}, ou un {\it champ en catégories discrètes},
est un champ sur $\mathcal{C}$ dont toutes les catégories fibres sont discrètes.
\end{enumerate}
\end{definition}

\noindent
D'après la discussion de la
section \ref{section-stacks-in-groupoids},
c'est la même chose qu'un champ en groupoïdes dont les catégories fibres
sont des setoïdes (resp.\ discrètes). De plus, c'est aussi la même chose
qu'une catégorie fibrée en setoïdes (resp.\ en ensembles) qui est un champ.

\begin{lemma}
\label{lemma-when-stack-in-sets}
Soit $\mathcal{C}$ un site. Sous l'équivalence
$$
\left\{
\begin{matrix}
\text{la catégorie des préfaisceaux}\\
\text{d'ensembles sur }\mathcal{C}
\end{matrix}
\right\}
\leftrightarrow
\left\{
\begin{matrix}
\text{la catégorie des catégories}\\
\text{fibrées en ensembles au-dessus de }\mathcal{C}
\end{matrix}
\right\}
$$
de
Catégories, lemme \ref{categories-lemma-2-category-fibred-sets},
les champs en ensembles correspondent précisément aux faisceaux.
\end{lemma}

\begin{proof}
Démonstration omise. Indication : montrer que l'effectivité de la descente correspond exactement à
la condition de faisceau.
\end{proof}

\begin{lemma}
\label{lemma-stack-in-setoids-characterize}
Soit $\mathcal{C}$ un site.
Soit $\mathcal{S}$ une catégorie fibrée en setoïdes au-dessus de $\mathcal{C}$.
Alors $\mathcal{S}$ est un champ en setoïdes si et seulement si l'unique
catégorie $\mathcal{S}'$ fibrée en ensembles qui lui est équivalente (voir
Catégories, lemme \ref{categories-lemma-setoid-fibres})
est un champ en ensembles. Autrement dit, si et seulement si le préfaisceau
$$
U \longmapsto \Ob(\mathcal{S}_U)/\!\!\cong
$$
est un faisceau.
\end{lemma}

\begin{proof}
Démonstration omise.
\end{proof}

\begin{lemma}
\label{lemma-stack-in-setoids-equivalent}
Soit $\mathcal{C}$ un site.
Soient $\mathcal{S}_1$, $\mathcal{S}_2$ des catégories au-dessus de $\mathcal{C}$.
Supposons que $\mathcal{S}_1$ et $\mathcal{S}_2$ soient équivalentes
comme catégories au-dessus de $\mathcal{C}$.
Alors $\mathcal{S}_1$ est un champ en setoïdes sur $\mathcal{C}$ si et seulement si
$\mathcal{S}_2$ est un champ en setoïdes sur $\mathcal{C}$.
\end{lemma}

\begin{proof}
D'après
Catégories, lemme \ref{categories-lemma-setoid-fibres},
nous voyons qu'une catégorie $\mathcal{S}$ au-dessus de $\mathcal{C}$ est fibrée en setoïdes
au-dessus de $\mathcal{C}$ si et seulement si elle est équivalente au-dessus de $\mathcal{C}$
à une catégorie fibrée en ensembles.
Nous voyons donc que $\mathcal{S}_1$ est fibrée en setoïdes au-dessus de $\mathcal{C}$
si et seulement si $\mathcal{S}_2$ est fibrée en setoïdes au-dessus de $\mathcal{C}$.
Le lemme résulte alors du
lemme \ref{lemma-stack-in-setoids-characterize}.
\end{proof}

\noindent
La $2$-catégorie des champs en setoïdes sur $\mathcal{C}$
est définie comme suit.

\begin{definition}
\label{definition-stacks-in-setoids-over-C}
Soit $\mathcal{C}$ un site.
La {\it $2$-catégorie des champs en setoïdes sur $\mathcal{C}$}
est la sous-$2$-catégorie de la $2$-catégorie des champs
sur $\mathcal{C}$ (voir Définition \ref{definition-stacks-over-C})
définie comme suit :
\begin{enumerate}
\item Ses objets sont les champs en setoïdes
$p : \mathcal{S} \to \mathcal{C}$.
\item Ses $1$-morphismes $(\mathcal{S}, p) \to (\mathcal{S}', p')$
sont les foncteurs $G : \mathcal{S} \to \mathcal{S}'$ tels que
$p' \circ G = p$. (Puisque tout morphisme est fortement cartésien,
tout foncteur les préserve.)
\item Ses $2$-morphismes $t : G \to H$, pour
$G, H : (\mathcal{S}, p) \to (\mathcal{S}', p')$,
sont les morphismes de foncteurs
tels que $p'(t_x) = \text{id}_{p(x)}$
pour tout $x \in \Ob(\mathcal{S})$.
\end{enumerate}
\end{definition}

\noindent
Remarquons que tout $2$-morphisme est automatiquement un isomorphisme, de sorte
que la $2$-catégorie des champs en setoïdes sur $\mathcal{C}$
est en fait une $(2, 1)$-catégorie (stricte).

\begin{lemma}
\label{lemma-2-product-stacks-in-setoids}
Soit $\mathcal{C}$ un site.
La $2$-catégorie des champs en setoïdes sur $\mathcal{C}$
admet des 2-produits fibrés, et ceux-ci sont décrits comme dans
Catégories, lemme \ref{categories-lemma-2-product-categories-over-C}.
\end{lemma}

\begin{proof}
Cela résulte immédiatement de
Catégories, lemmes \ref{categories-lemma-2-product-fibred-categories} et
\ref{categories-lemma-2-product-categories-fibred-setoids},
ainsi que des
lemmes \ref{lemma-stack-in-groupoids-stack} et
\ref{lemma-2-product-stacks}.
\end{proof}

\begin{lemma}
\label{lemma-2-fibre-product-gives-stack-in-setoids}
Soit $\mathcal{C}$ un site.
Soient $\mathcal{S}, \mathcal{T}$ des champs en groupoïdes sur $\mathcal{C}$
et soit $\mathcal{R}$ un champ en setoïdes sur $\mathcal{C}$.
Soient $f : \mathcal{T} \to \mathcal{S}$ et $g : \mathcal{R} \to \mathcal{S}$
des $1$-morphismes. Si $f$ est fidèle, alors le $2$-produit fibré
$$
\mathcal{T} \times_{f, \mathcal{S}, g} \mathcal{R}
$$
est un champ en setoïdes sur $\mathcal{C}$.
\end{lemma}

\begin{proof}
Cela résulte immédiatement de la description explicite du $2$-produit fibré dans
Catégories, lemme \ref{categories-lemma-2-product-categories-over-C}.
\end{proof}

\begin{lemma}
\label{lemma-2-fibre-product-stacks-in-setoids-over-stack-in-groupoids}
Soit $\mathcal{C}$ un site.
Soit $\mathcal{S}$ un champ en groupoïdes sur $\mathcal{C}$ et
soient $\mathcal{S}_i$, $i = 1, 2$, des champs en setoïdes sur $\mathcal{C}$.
Soient $f_i : \mathcal{S}_i \to \mathcal{S}$ des $1$-morphismes.
Alors le $2$-produit fibré
$$
\mathcal{S}_1 \times_{f_1, \mathcal{S}, f_2} \mathcal{S}_2
$$
est un champ en setoïdes sur $\mathcal{C}$.
\end{lemma}

\begin{proof}
C'est un cas particulier du
lemme \ref{lemma-2-fibre-product-gives-stack-in-setoids},
car $f_2$ est fidèle.
\end{proof}

\begin{lemma}
\label{lemma-faithful-descent}
Soit $\mathcal{C}$ un site. Soit
$$
\xymatrix{
\mathcal{T}_2 \ar[r] \ar[d]_{G'} & \mathcal{T}_1 \ar[d]^G \\
\mathcal{S}_2 \ar[r]^F & \mathcal{S}_1
}
$$
un diagramme $2$-cartésien de champs en groupoïdes sur $\mathcal{C}$.
Supposons que
\begin{enumerate}
\item pour tout $U \in \Ob(\mathcal{C})$ et tout
$x \in \Ob((\mathcal{S}_1)_U)$, il existe un recouvrement
$\{U_i \to U\}$ tel que $x|_{U_i}$ appartienne à l'image
essentielle de $F : (\mathcal{S}_2)_{U_i} \to (\mathcal{S}_1)_{U_i}$, et
\item $G'$ soit fidèle,
\end{enumerate}
alors $G$ est fidèle.
\end{lemma}

\begin{proof}
Nous pouvons supposer que $\mathcal{T}_2$ est la catégorie
$\mathcal{S}_2 \times_{\mathcal{S}_1} \mathcal{T}_1$
décrite dans
Catégories, lemme \ref{categories-lemma-2-product-categories-over-C}.
D'après
Catégories, lemme \ref{categories-lemma-equivalence-fibred-categories},
la fidélité de $G, G'$ peut être vérifiée sur les catégories fibres.
Supposons que $y, y'$ soient des objets de $\mathcal{T}_1$ au-dessus de l'objet $U$
de $\mathcal{C}$. Soient $\alpha, \beta : y \to y'$ des morphismes de
$(\mathcal{T}_1)_U$ tels que $G(\alpha) = G(\beta)$. Notre but est de
montrer que $\alpha = \beta$. En considérant plutôt
$\gamma = \alpha^{-1} \circ \beta$, nous voyons que $G(\gamma) = \text{id}_{G(y)}$
et nous devons montrer que $\gamma = \text{id}_y$. Par hypothèse, nous
pouvons trouver un recouvrement $\{U_i \to U\}$ tel que $G(y)|_{U_i}$ appartienne à
l'image essentielle de $F :(\mathcal{S}_2)_{U_i} \to (\mathcal{S}_1)_{U_i}$.
Puisqu'il suffit de montrer que $\gamma|_{U_i} = \text{id}$ pour tout $i$,
nous pouvons donc supposer qu'il existe
$f : F(x) \to G(y)$ pour un objet $x$ de $\mathcal{S}_2$ au-dessus de $U$,
la flèche $f$ étant un morphisme de $(\mathcal{S}_1)_U$. Dans ce cas, nous obtenons
un morphisme
$$
(1, \gamma) : (U, x, y, f) \longrightarrow (U, x, y, f)
$$
dans la catégorie fibre de $\mathcal{S}_2 \times_{\mathcal{S}_1} \mathcal{T}_1$
au-dessus de $U$, dont l'image par $G'$ dans $\mathcal{S}_1$ est $\text{id}_x$.
Comme $G'$ est fidèle, nous concluons que $\gamma = \text{id}_y$, ce qui termine la preuve.
\end{proof}

\begin{lemma}
\label{lemma-stack-in-setoids-descent}
Soit $\mathcal{C}$ un site.
Soit
$$
\xymatrix{
\mathcal{T}_2 \ar[r] \ar[d] & \mathcal{T}_1 \ar[d]^G \\
\mathcal{S}_2 \ar[r]^F & \mathcal{S}_1
}
$$
un diagramme $2$-cartésien de champs en groupoïdes sur $\mathcal{C}$.
Si
\begin{enumerate}
\item $F : \mathcal{S}_2 \to \mathcal{S}_1$ est pleinement fidèle,
\item pour tout $U \in \Ob(\mathcal{C})$ et tout
$x \in \Ob((\mathcal{S}_1)_U)$, il existe un recouvrement
$\{U_i \to U\}$ tel que $x|_{U_i}$ appartienne à l'image essentielle
de $F : (\mathcal{S}_2)_{U_i} \to (\mathcal{S}_1)_{U_i}$, et
\item $\mathcal{T}_2$ est un champ en setoïdes,
\end{enumerate}
alors $\mathcal{T}_1$ est un champ en setoïdes.
\end{lemma}

\begin{proof}
Nous pouvons supposer que $\mathcal{T}_2$ est la catégorie
$\mathcal{S}_2 \times_{\mathcal{S}_1} \mathcal{T}_1$
décrite dans
Catégories, lemme \ref{categories-lemma-2-product-categories-over-C}.
Choisissons $U \in \Ob(\mathcal{C})$ et
$y \in \Ob((\mathcal{T}_1)_U)$.
Nous devons montrer que le faisceau $\mathit{Aut}(y)$ sur $\mathcal{C}/U$
est trivial. Pour cela, nous pouvons remplacer $U$ par les membres d'un
recouvrement de $U$. Par l'hypothèse (2), nous pouvons donc supposer
qu'il existe un objet $x \in \Ob((\mathcal{S}_2)_U)$
et un isomorphisme $f : F(x) \to G(y)$.
Alors $y' = (U, x, y, f)$ est un objet de $\mathcal{T}_2$ au-dessus de $U$
qui est envoyé sur $y$ par la projection $\mathcal{T}_2 \to \mathcal{T}_1$.
Puisque $F$ est pleinement fidèle d'après (1), l'application
$\mathit{Aut}(y') \to \mathit{Aut}(y)$ est surjective, comme le montre la description
explicite des morphismes de $\mathcal{T}_2$ donnée dans
Catégories, lemme \ref{categories-lemma-2-product-categories-over-C}.
Puisque, par (3), le faisceau $\mathit{Aut}(y')$ est trivial,
nous obtenons le résultat du lemme.
\end{proof}

\begin{lemma}
\label{lemma-relative-sheaf-over-stack-is-stack}
Soit $\mathcal{C}$ un site. Soit $F : \mathcal{S} \to \mathcal{T}$
un $1$-morphisme de catégories fibrées en groupoïdes au-dessus de $\mathcal{C}$.
Supposons que
\begin{enumerate}
\item $\mathcal{T}$ soit un champ en groupoïdes sur $\mathcal{C}$,
\item pour tout $U \in \Ob(\mathcal{C})$, le foncteur
$\mathcal{S}_U \to \mathcal{T}_U$ entre catégories fibres soit fidèle,
\item pour tout $U$ et tout $y \in \Ob(\mathcal{T}_U)$, le préfaisceau
$$
(h : V \to U)
\longmapsto
\{(x, f) \mid x \in \Ob(\mathcal{S}_V), f : F(x) \to f^*y\text{ au-dessus de }V\}/\cong
$$
soit un faisceau sur $\mathcal{C}/U$.
\end{enumerate}
Alors $\mathcal{S}$ est un champ en groupoïdes sur $\mathcal{C}$.
\end{lemma}

\begin{proof}
Nous devons démontrer la descente des morphismes et la descente des objets.

\medskip\noindent
Descente des morphismes. Soit $\{U_i \to U\}$ un recouvrement de $\mathcal{C}$.
Soient $x, x'$ des objets de $\mathcal{S}$ au-dessus de $U$. Pour tout $i$,
soit $\alpha_i : x|_{U_i} \to x'|_{U_i}$ un morphisme au-dessus de $U_i$
tel que les restrictions de $\alpha_i$ et de $\alpha_j$ coïncident comme morphismes
$x|_{U_i \times_U U_j} \to x'|_{U_i \times_U U_j}$.
Puisque $\mathcal{T}$ est un champ en groupoïdes, il existe un morphisme
$\beta : F(x) \to F(x')$ au-dessus de $U$ dont la restriction à $U_i$ est $F(\alpha_i)$.
Nous pouvons alors considérer $\xi = (x, \beta)$ et $\xi' = (x', \text{id}_{F(x')})$
comme des sections du préfaisceau associé à $y = F(x')$ au-dessus de $U$
dans l'hypothèse (3). D'autre part, les restrictions de
$\xi$ et $\xi'$ à $U_i$ sont $(x|_{U_i}, F(\alpha_i))$
et $(x'|_{U_i}, \text{id}_{F(x'|_{U_i})})$.
Elles sont isomorphes l'une à l'autre par le morphisme $\alpha_i$.
Ainsi, $\xi$ et $\xi'$ sont isomorphes d'après l'hypothèse (3). Cela signifie qu'il existe un
morphisme $\alpha : x \to x'$ au-dessus de $U$ tel que $F(\alpha) = \beta$.
Puisque $F$ est fidèle sur les catégories fibres, nous obtenons
$\alpha|_{U_i} = \alpha_i$.

\medskip\noindent
Descente des objets. Soit $\{U_i \to U\}$ un recouvrement de $\mathcal{C}$.
Soit $(x_i, \varphi_{ij})$ une donnée de descente pour $\mathcal{S}$
relativement au recouvrement donné. Puisque $\mathcal{T}$ est un champ en groupoïdes,
il existe un objet $y$ de $\mathcal{T}_U$ et des isomorphismes
$\beta_i : F(x_i) \to y|_{U_i}$ tels que
$F(\varphi_{ij}) = \beta_j|_{U_i \times_U U_j} \circ
(\beta_i|_{U_i \times_U U_j})^{-1}$.
Alors les $(x_i, \beta_i)$ sont des sections du préfaisceau associé à
$y$ au-dessus de $U$ défini dans l'hypothèse (3).
De plus, $\varphi_{ij}$ définit un isomorphisme
de la paire $(x_i, \beta_i)|_{U_i \times_U U_j}$ vers
la paire $(x_j, \beta_j)|_{U_i \times_U U_j}$.
Ainsi, d'après l'hypothèse (3), il existe une paire $(x, \beta)$ au-dessus de $U$
dont la restriction à $U_i$ est isomorphe à $(x_i, \beta_i)$.
Cela signifie qu'il existe des morphismes $\alpha_i : x_i \to x|_{U_i}$
tels que $\beta_i = \beta|_{U_i} \circ F(\alpha_i)$.
Puisque $F$ est fidèle sur les catégories fibres, un calcul montre
que $\varphi_{ij} = \alpha_j|_{U_i \times_U U_j} \circ
(\alpha_i|_{U_i \times_U U_j})^{-1}$. Cela achève la démonstration.
\end{proof}










\section{Le champ d'inertie}
\label{section-the-inertia-stack}

\noindent
Soient
$p : \mathcal{S} \to \mathcal{C}$ et
$p' : \mathcal{S}' \to \mathcal{C}$
des catégories fibrées au-dessus de la catégorie $\mathcal{C}$.
Soit $F : \mathcal{S} \to \mathcal{S}'$ un $1$-morphisme de
catégories fibrées au-dessus de $\mathcal{C}$.
Rappelons que nous avons défini dans
Catégories, définition \ref{categories-definition-inertia-fibred-category},
une {\it catégorie fibrée d'inertie relative}
$\mathcal{I}_{\mathcal{S}/\mathcal{S}'} \to \mathcal{C}$ comme la
catégorie dont les objets sont les paires $(x , \alpha)$ où
$x \in \Ob(\mathcal{S})$ et $\alpha : x \to x$ vérifie
$F(\alpha) = \text{id}_{F(x)}$. Il existe aussi une version absolue,
à savoir l'{\it inertie} $\mathcal{I}_\mathcal{S}$ de $\mathcal{S}$.
Ces catégories d'inertie sont en fait des champs sur $\mathcal{C}$ dès que
$\mathcal{S}$ et $\mathcal{S}'$ sont des champs.

\begin{lemma}
\label{lemma-inertia}
Soit $\mathcal{C}$ un site. Soient
$p : \mathcal{S} \to \mathcal{C}$ et
$p' : \mathcal{S}' \to \mathcal{C}$
des champs sur le site $\mathcal{C}$.
Soit $F : \mathcal{S} \to \mathcal{S}'$ un $1$-morphisme de
champs sur $\mathcal{C}$.
\begin{enumerate}
\item Les inerties $\mathcal{I}_{\mathcal{S}/\mathcal{S}'}$ et
$\mathcal{I}_\mathcal{S}$ sont des champs sur $\mathcal{C}$.
\item Si $\mathcal{S}, \mathcal{S}'$ sont des champs en groupoïdes sur
$\mathcal{C}$, alors $\mathcal{I}_{\mathcal{S}/\mathcal{S}'}$ et
$\mathcal{I}_\mathcal{S}$ le sont également.
\item Si $\mathcal{S}, \mathcal{S}'$ sont des champs en setoïdes sur $\mathcal{C}$,
alors $\mathcal{I}_{\mathcal{S}/\mathcal{S}'}$ et
$\mathcal{I}_\mathcal{S}$ le sont également.
\end{enumerate}
\end{lemma}

\begin{proof}
Les trois assertions résultent des
lemmes \ref{lemma-2-product-stacks},
\ref{lemma-2-product-stacks-in-groupoids} et
\ref{lemma-2-product-stacks-in-setoids},
ainsi que de l'équivalence de
Catégories, lemme \ref{categories-lemma-inertia-fibred-category}, partie (1).
\end{proof}

\begin{lemma}
\label{lemma-characterize-stack-in-setoids}
Soit $\mathcal{C}$ un site.
Si $\mathcal{S}$ est un champ en groupoïdes, alors le
$1$-morphisme canonique $\mathcal{I}_\mathcal{S} \to \mathcal{S}$
est une équivalence si et seulement si $\mathcal{S}$ est un champ en setoïdes.
\end{lemma}

\begin{proof}
Cela résulte directement de
Catégories, lemme \ref{categories-lemma-characterize-fibred-setoids-inertia}.
\end{proof}

\section{Champification des catégories fibrées}
\label{section-stackify}

\noindent
Le résultat de la procédure du lemme suivant sera appelé la
{\it champification} d'une catégorie fibrée sur un site.

\begin{lemma}
\label{lemma-stackify}
Soit $\mathcal{C}$ un site.
Soit $p : \mathcal{S} \to \mathcal{C}$ une catégorie fibrée sur $\mathcal{C}$.
Il existe un champ $p' : \mathcal{S}' \to \mathcal{C}$ et un
$1$-morphisme $G : \mathcal{S} \to \mathcal{S}'$
de catégories fibrées sur $\mathcal{C}$ (voir
Catégories, Définition \ref{categories-definition-fibred-categories-over-C})
tels que
\begin{enumerate}
\item pour tout $U \in \Ob(\mathcal{C})$ et tous
$x, y \in \Ob(\mathcal{S}_U)$, l'application
$$
\mathit{Mor}(x, y) \longrightarrow \mathit{Mor}(G(x), G(y))
$$
induite par $G$ identifie le membre de droite à la faisceautisation
du membre de gauche, et
\item pour tout $U \in \Ob(\mathcal{C})$ et tout
$x' \in \Ob(\mathcal{S}'_U)$, il existe un recouvrement
$\{U_i \to U\}_{i \in I}$ tel que, pour tout $i \in I$,
l'objet $x'|_{U_i}$ appartienne à l'image essentielle du
foncteur $G : \mathcal{S}_{U_i} \to \mathcal{S}'_{U_i}$.
\end{enumerate}
De plus, ces conditions déterminent le champ $\mathcal{S}'$ à un unique
$2$-isomorphisme près.
\end{lemma}

\begin{proof}[Démonstration par la méthode naïve]
Dans cette méthode de démonstration, nous procédons par étapes :

\medskip\noindent
Tout d'abord, étant donnés $x$ au-dessus de $U$ et un objet quelconque $y$ de
$\mathcal{S}$, on dit que deux morphismes
$a, b : x \to y$ de $\mathcal{S}$
au-dessus d'une même flèche de $\mathcal{C}$
sont {\it localement égaux}
s'il existe un recouvrement $\{f_i : U_i \to U\}$ de $\mathcal{C}$
tel que les composées
$$
f_i^*x \to x \xrightarrow{a} y,
\quad
f_i^*x \to x \xrightarrow{b} y
$$
soient égales. Cela définit une relation d'équivalence $\sim$
sur les flèches de $\mathcal{S}$. Si $b \sim b'$, alors
$a \circ b \circ c \sim a \circ b' \circ c$ (vérification omise).
On peut donc quotienter par cette relation d'équivalence pour
obtenir une nouvelle catégorie $\mathcal{S}^1$ sur $\mathcal{C}$,
ainsi qu'un morphisme $G^1 : \mathcal{S} \to \mathcal{S}^1$.

\medskip\noindent
On vérifie que $G^1$ préserve les morphismes fortement cartésiens
et que $\mathcal{S}^1$ est une catégorie fibrée sur $\mathcal{C}$.
Les vérifications sont omises. On se ramène ainsi au cas où des
morphismes localement égaux sont égaux.

\medskip\noindent
Ensuite, on ajoute des morphismes comme suit. Étant donnés
$x$ au-dessus de $U$ et un objet quelconque $y$ au-dessus de $V$,
un {\it morphisme de $x$ vers $y$ défini localement} consiste en
\begin{enumerate}
\item un morphisme $f : U \to V$,
\item un recouvrement $\{f_i : U_i \to U\}$ de $U$, et
\item des morphismes $a_i : f_i^*x \to y$ avec $p(a_i) = f \circ f_i$
\end{enumerate}
ayant la propriété que les composées
$$
(f_i \times f_j)^*x \to f_i^*x \xrightarrow{a_i} y,
\quad
(f_i \times f_j)^*x \to f_j^*x \xrightarrow{a_j} y
$$
sont égales. Remarquons qu'un morphisme usuel $a : x \to y$ fournit le
morphisme défini localement $(p(a) : U \to V, \{\text{id}_U\}, a)$.
On dit que deux morphismes définis localement
$(f, \{f_i : U_i \to U\}, a_i)$ et $(g, \{g_j : U'_j \to U\}, b_j)$
sont {\it égaux} si $f = g$ et si les composées
$$
(f_i \times g_j)^*x \to f_i^*x \xrightarrow{a_i} y,
\quad
(f_i \times g_j)^*x \to g_j^*x \xrightarrow{b_j} y
$$
sont égales (c'est la bonne condition puisque nous sommes dans le
cas où des morphismes localement égaux sont égaux).
Pour composer les morphismes définis localement
$(f, \{f_i : U_i \to U\}, a_i)$ de $x$ vers $y$ et
$(g, \{g_j : V_j \to V\}, b_j)$ de $y$ vers $z$ au-dessus de $W$,
il suffit de prendre $g \circ f : U \to W$, le recouvrement
$\{U_i \times_V V_j \to U\}$ et, comme applications, les composées
$$
x|_{U_i \times_V V_j}
\xrightarrow{\text{pr}_0^*a_i}
y|_{V_j}
\xrightarrow{b_j}
z
$$
Nous omettons la vérification qu'il s'agit d'un morphisme défini localement.

\medskip\noindent
On vérifie que $\mathcal{S}^2$, qui a les mêmes objets que
$\mathcal{S}$ et pour morphismes les morphismes définis localement,
est une catégorie sur $\mathcal{C}$, qu'il existe un foncteur
$G^2 : \mathcal{S} \to \mathcal{S}^2$ sur $\mathcal{C}$,
que ce foncteur préserve les objets fortement cartésiens
et que $\mathcal{S}^2$ est une catégorie fibrée sur $\mathcal{C}$.
Les vérifications sont omises. On se ramène ainsi au cas où tous les
préfaisceaux de morphismes de $\mathcal{S}$ sont des faisceaux, en
vérifiant que l'emploi de morphismes définis localement a pour effet de
faisceautiser les préfaisceaux (séparés) de morphismes.

\medskip\noindent
Enfin, lorsque les préfaisceaux de morphismes sont tous des faisceaux,
il faut ajouter des objets afin que les données de descente soient
effectives dans le résultat final. La façon la plus simple consiste à
considérer la catégorie $\mathcal{S}'$ dont les objets sont les
paires $(\mathcal{U}, \xi)$ où
$\mathcal{U} = \{U_i \to U\}$ est un recouvrement de $\mathcal{C}$ et
$\xi = (X_i, \varphi_{ii'})$ est une donnée de descente relativement à $\mathcal{U}$.
Supposons données deux telles données
$(\mathcal{U}, \xi) = (\{f_i : U_i \to U\},  x_i, \varphi_{ii'})$ et
$(\mathcal{V}, \eta) = (\{g_j : V_j \to V\},  y_j, \psi_{jj'})$.
On définit
$$
\Mor_{\mathcal{S}'}((\mathcal{U}, \xi), (\mathcal{V}, \eta))
$$
comme l'ensemble des $(f, a_{ij})$, où $f : U \to V$ et
$$
a_{ij} :
x_i|_{U_i \times_V V_j}
\longrightarrow
y_j
$$
sont des morphismes de $\mathcal{S}$ au-dessus de $U_i \times_V V_j \to V_j$.
Ils doivent satisfaire la condition suivante : pour tous
$i, i' \in I$ et $j, j' \in J$, posons
$W = (U_i \times_U U_{i'}) \times_V (V_j \times_V V_{j'})$. Alors
$$
\xymatrix{
x_i|_W \ar[r]_{a_{ij}|_W} \ar[d]_{\varphi_{ii'}|_W} &
y_j|_W \ar[d]^{\psi_{jj'}|_W} \\
x_{i'}|_W \ar[r]^{a_{i'j'}|_W} &
y_{j'}|_W
}
$$
est commutatif. Il faut alors vérifier les points suivants :
\begin{enumerate}
\item les morphismes ci-dessus admettent une composition bien définie,
\item cela fait de $\mathcal{S}'$ une catégorie sur $\mathcal{C}$,
\item il existe un foncteur $G : \mathcal{S} \to \mathcal{S}'$ sur $\mathcal{C}$,
\item si $x, y$ sont des objets de $\mathcal{S}$, alors
$\Mor_\mathcal{S}(x, y) = \Mor_{\mathcal{S}'}(G(x), G(y))$,
\item tout objet de $\mathcal{S}'$ provient localement d'un objet de
$\mathcal{S}$, autrement dit l'assertion (2) du lemme est satisfaite,
\item $G$ préserve les morphismes fortement cartésiens,
\item $\mathcal{S}'$ est une catégorie fibrée sur $\mathcal{C}$, et
\item $\mathcal{S}'$ est un champ sur $\mathcal{C}$.
\end{enumerate}
Rien de tout cela n'est difficile, mais les vérifications sont nombreuses. Détails omis.
\end{proof}

\begin{proof}[Démonstration moins naïve]
Voici une démonstration moins naïve.
D'après Catégories, Lemme \ref{categories-lemma-fibred-strict},
il existe une équivalence de
catégories fibrées $\mathcal{S} \to \mathcal{S}'$ telle que $\mathcal{S}'$
soit une catégorie fibrée scindée, c'est-à-dire une catégorie dans laquelle les foncteurs
d'image inverse se composent au sens strict. Le lemme pour $\mathcal{S}'$
implique évidemment le lemme pour $\mathcal{S}$. On peut donc considérer $\mathcal{S}$
comme un préfaisceau en catégories.

\medskip\noindent
Considérons temporairement la $2$-catégorie $\textit{Cat}$ comme une
catégorie, en oubliant les $2$-morphismes.
Considérons une catégorie comme un quintuplet
$(\text{Ob}, \text{Arrows}, s, t, \circ)$ comme dans
Catégories, Section \ref{categories-section-definition-categories}.
Considérons le foncteur d'oubli
$$
forget : \textit{Cat} \to \textit{Sets} \times \textit{Sets}, \quad
(\text{Ob}, \text{Arrows}, s, t, \circ)
\mapsto
(\text{Ob}, \text{Arrows}).
$$
Alors $forget$ est fidèle, $\textit{Cat}$ admet des limites et
$forget$ commute à celles-ci, $\textit{Cat}$ admet des colimites filtrantes et
$forget$ commute à celles-ci, et $forget$ reflète les isomorphismes.
On peut faisceautiser les préfaisceaux à valeurs dans $\textit{Cat}$ et,
par un argument analogue à celui de la première partie de
Sites, Section \ref{sites-section-sheaves-algebraic-structures},
le résultat commute avec $forget$. En appliquant ceci à
$\mathcal{S}$, on obtient une faisceautisation $\mathcal{S}^\#$
qui possède un faisceau d'objets et un faisceau de morphismes,
tous deux faisceautisations des préfaisceaux correspondants
de $\mathcal{S}$. Dans ce cas, il est assez
facile de voir que l'application $\mathcal{S} \to \mathcal{S}^\#$
possède les propriétés (1) et (2) du lemme.

\medskip\noindent
Cependant, la catégorie $\mathcal{S}^\#$ n'est pas nécessairement encore un
champ car, bien que le préfaisceau d'objets soit un faisceau,
la condition de descente peut ne pas être satisfaite.
Pour y remédier, il faut ajouter davantage d'objets. Mais l'argument
ci-dessus nous ramène au cas où $\mathcal{S} = \mathcal{S}_F$
pour un faisceau (!) de catégories
$F : \mathcal{C}^{opp} \to \textit{Cat}$. Dans ce cas, considérons le foncteur
$F' : \mathcal{C}^{opp} \to \textit{Cat}$ défini comme suit :
\begin{enumerate}
\item L'ensemble $\Ob(F'(U))$ est l'ensemble des paires
$(\mathcal{U}, \xi)$ où $\mathcal{U} = \{U_i \to U\}$
est un recouvrement de $U$ et $\xi = (x_i, \varphi_{ii'})$ est
une donnée de descente relativement à $\mathcal{U}$.
\item Un morphisme dans $F'(U)$ de
$(\mathcal{U}, \xi)$ vers $(\mathcal{V}, \eta)$
est un élément de
$$
\colim \Mor_{DD(\mathcal{W})}(a^*\xi, b^*\eta)
$$
où la colimite porte sur tous les raffinements communs
$a : \mathcal{W} \to \mathcal{U}$, $b : \mathcal{W} \to \mathcal{V}$.
Cette colimite est filtrante (vérification omise).
La composition des morphismes dans $F(U)$ est donc définie en
choisissant un raffinement commun puis en composant dans $DD(\mathcal{W})$.
\item Étant donnés $h : V \to U$ et un objet
$(\mathcal{U}, \xi)$ de $F'(U)$, on définit $F'(h)(\mathcal{U}, \xi)$
comme $(V \times_U \mathcal{U}, \text{pr}_1^*\xi)$.
Plus précisément, si $\mathcal{U} = \{U_i \to U\}$
et $\xi = (x_i, \varphi_{ii'})$, alors
$V \times_U \mathcal{U} = \{V \times_U U_i \to V\}$,
muni d'un morphisme canonique
$\text{pr}_1 : V \times_U \mathcal{U} \to \mathcal{U}$, et
$\text{pr}_1^*\xi$ est l'image inverse de $\xi$ par
ce morphisme (voir la Définition \ref{definition-pullback-functor}).
\item Soient $h : V \to U$, des objets $(\mathcal{U}, \xi)$
et $(\mathcal{V}, \eta)$, ainsi qu'un morphisme entre eux représenté par
$a : \mathcal{W} \to \mathcal{U}$, $b : \mathcal{W} \to \mathcal{V}$
et $\alpha : a^*\xi \to b^*\eta$ ; alors $F'(h)(\alpha)$ est
représenté par
$a' : V \times_U\mathcal{W} \to V \times_U\mathcal{U}$,
$b' : V \times_U\mathcal{W} \to V \times_U\mathcal{V}$
et l'image inverse $\alpha'$ du morphisme $\alpha$ par
l'application $V \times_U \mathcal{W} \to \mathcal{W}$. Cela fonctionne
car les images inverses dans $\mathcal{S}_F$ commutent au sens strict.
\end{enumerate}
Il existe une application $F \to F'$ qui associe à
un objet $x$ de $F(U)$ l'objet $(\{U \to U\}, (x, triv))$ de
$F'(U)$. Il reste alors à vérifier que le foncteur correspondant
$\mathcal{S}_F \to \mathcal{S}_{F'}$ possède les propriétés (1)
et (2) du lemme, et enfin que $\mathcal{S}_{F'}$ est
un champ. Détails omis.
\end{proof}

\begin{lemma}
\label{lemma-stackify-universal-property}
Soit $\mathcal{C}$ un site.
Soit $p : \mathcal{S} \to \mathcal{C}$ une catégorie fibrée sur $\mathcal{C}$.
Soient $p' : \mathcal{S}' \to \mathcal{C}$ et $G : \mathcal{S} \to \mathcal{S}'$
respectivement le champ et le $1$-morphisme construits dans le Lemme \ref{lemma-stackify}.
Cette construction possède la propriété universelle suivante : étant donnés un champ
$q : \mathcal{X} \to \mathcal{C}$ et un $1$-morphisme
$F : \mathcal{S} \to \mathcal{X}$ de catégories fibrées sur $\mathcal{C}$,
il existe un $1$-morphisme $H : \mathcal{S}' \to \mathcal{X}$
tel que le diagramme
$$
\xymatrix{
\mathcal{S} \ar[rr]_F \ar[rd]_G & & \mathcal{X} \\
& \mathcal{S}' \ar[ru]_H
}
$$
soit $2$-commutatif.
\end{lemma}

\begin{proof}
Omis. Indication : supposons que $x' \in \Ob(\mathcal{S}'_U)$.
D'après le Lemme \ref{lemma-stackify},
il existe un recouvrement $\{U_i \to U\}_{i \in I}$
tel que $x'|_{U_i} = G(x_i)$ pour un certain $x_i \in \Ob(\mathcal{S}_{U_i})$.
De plus, il existe des recouvrements $\{U_{ijk} \to U_i \times_U U_j\}$
et des isomorphismes $\alpha_{ijk} : x_i|_{U_{ijk}} \to x_j|_{U_{ijk}}$
tels que $G(\alpha_{ijk}) = \text{id}_{x'|_{U_{ijk}}}$. Posons $y_i = F(x_i)$.
On peut alors vérifier que les morphismes
$$
F(\alpha_{ijk}) : y_i|_{U_{ijk}} \to y_j|_{U_{ijk}}
$$
coïncident sur les intersections et définissent donc (puisque $\mathcal{X}$ est un champ)
un morphisme $\beta_{ij} : y_i|_{U_i \times_U U_j} \to y_j|_{U_i \times_U U_j}$.
On vérifie ensuite que les $\beta_{ij}$ définissent une donnée de descente. Comme
$\mathcal{X}$ est un champ, ces données de descente sont effectives et l'on trouve
un objet $y$ de $\mathcal{X}_U$ coïncidant avec $G(x_i)$ au-dessus de $U_i$.
L'indication consiste à poser $H(x') = y$.
\end{proof}

\begin{lemma}
\label{lemma-stackify-universal-property-more}
Les notations et les hypothèses sont celles du
Lemme \ref{lemma-stackify-universal-property}.
Il existe une équivalence canonique de catégories
$$
\Mor_{\textit{Fib}/\mathcal{C}}(\mathcal{S}, \mathcal{X})
=
\Mor_{\textit{Stacks}/\mathcal{C}}(\mathcal{S}', \mathcal{X})
$$
donnée par les constructions de la démonstration du lemme précité.
\end{lemma}

\begin{proof}
Omis.
\end{proof}

\begin{lemma}
\label{lemma-stackification-fibre-product-fibred-categories}
Soit $\mathcal{C}$ un site.
Soient $f : \mathcal{X} \to \mathcal{Y}$ et $g : \mathcal{Z} \to \mathcal{Y}$
des morphismes de catégories fibrées sur $\mathcal{C}$.
Alors la champification du $2$-produit fibré est le $2$-produit
fibré des champifications.
\end{lemma}

\begin{proof}
Notons $\mathcal{X}', \mathcal{Y}', \mathcal{Z}'$ les champifications
et $\mathcal{W}$ la champification de
$\mathcal{X} \times_\mathcal{Y} \mathcal{Z}$. La construction des $2$-produits
fibrés fournit un $1$-morphisme canonique
$\mathcal{X} \times_\mathcal{Y} \mathcal{Z} \to
\mathcal{X}' \times_{\mathcal{Y}'} \mathcal{Z}'$.
Comme le second $2$-produit fibré est un champ (voir le
Lemme \ref{lemma-2-product-stacks}),
ce $1$-morphisme induit un $1$-morphisme
$h : \mathcal{W} \to \mathcal{X}' \times_{\mathcal{Y}'} \mathcal{Z}'$
par la propriété universelle de la champification, voir le
Lemme \ref{lemma-stackify-universal-property}.
Or $h$ est un morphisme de champs, et l'on peut vérifier que c'est une
équivalence à l'aide des
Lemmes \ref{lemma-characterize-ff} et
\ref{lemma-characterize-essentially-surjective-when-ff}.

\medskip\noindent
Montrons donc d'abord que $h$ induit des isomorphismes de faisceaux $\mathit{Mor}$.
Soient $\xi, \xi'$ des objets de $\mathcal{W}$ au-dessus de
$U \in \Ob(\mathcal{C})$. Nous voulons montrer que
$$
h : \mathit{Mor}(\xi, \xi') \longrightarrow \mathit{Mor}(h(\xi), h(\xi'))
$$
est un isomorphisme. Pour cela, on peut travailler localement sur $U$ (voir
Sites, Section \ref{sites-section-glueing-sheaves}).
Par la construction de $\mathcal{W}$ (voir le
Lemme \ref{lemma-stackify}),
on peut donc supposer que $\xi, \xi'$ proviennent effectivement
des objets $(x, z, \alpha)$ et $(x', z', \alpha')$
de $\mathcal{X} \times_\mathcal{Y} \mathcal{Z}$ au-dessus de $U$.
Le même lemme montre encore que, dans ce cas,
$\mathit{Mor}(\xi, \xi')$ est la faisceautisation de
$$
V/U \longmapsto
\Mor_{\mathcal{X}_V}(x|_V, x'|_V)
\times_{\Mor_{\mathcal{Y}_V}(f(x)|_V, f(x')|_V)}
\Mor_{\mathcal{Z}_V}(z|_V, z'|_V)
$$
et que $\mathit{Mor}(h(\xi), h(\xi'))$ est égal au produit fibré
$$
\mathit{Mor}(i(x), i(x'))
\times_{\mathit{Mor}(j(f(x)), j(f(x'))}
\mathit{Mor}(k(z), k(z'))
$$
où $i : \mathcal{X} \to \mathcal{X}'$,
$j : \mathcal{Y} \to \mathcal{Y}'$ et
$k : \mathcal{Z} \to \mathcal{Z}'$ sont les foncteurs canoniques.
La première application affichée de ce paragraphe est donc un isomorphisme, car
la faisceautisation est exacte (et la faisceautisation d'un produit fibré
de préfaisceaux est par conséquent le produit fibré des faisceautisations).

\medskip\noindent
Enfin, il faut vérifier que tout objet de
$\mathcal{X}' \times_{\mathcal{Y}'} \mathcal{Z}'$
au-dessus de $U$ appartient localement sur $U$ à l'image essentielle de $h$.
Écrivons un tel objet sous la forme d'un triplet $(x', z', \alpha)$.
Alors $x'$ provient localement d'un objet de $\mathcal{X}$,
$z'$ provient localement d'un objet de $\mathcal{Z}$ et,
après avoir effectué des remplacements convenables de $x'$ et $z'$, le morphisme
$\alpha$ de $\mathcal{Y}'_U$ provient localement d'un morphisme de
$\mathcal{Y}$. Autrement dit, nous avons montré que tout objet de
$\mathcal{X}' \times_{\mathcal{Y}'} \mathcal{Z}'$
au-dessus de $U$ appartient localement sur $U$ à l'image essentielle de
$\mathcal{X} \times_\mathcal{Y} \mathcal{Z} \to
\mathcal{X}' \times_{\mathcal{Y}'} \mathcal{Z}'$ ; il appartient donc
a fortiori localement à l'image essentielle de $h$.
\end{proof}

\begin{lemma}
\label{lemma-stackification-inertia}
Soit $\mathcal{C}$ un site.
Soit $\mathcal{X}$ une catégorie fibrée sur $\mathcal{C}$.
La champification de la catégorie fibrée d'inertie $\mathcal{I}_\mathcal{X}$
est l'inertie de la champification de $\mathcal{X}$.
\end{lemma}

\begin{proof}
Cela résulte du fait que la champification est compatible
avec les $2$-produits fibrés d'après le
Lemme \ref{lemma-stackification-fibre-product-fibred-categories},
et du fait que l'inertie s'exprime en termes de
$2$-produits fibrés de catégories sur $\mathcal{C}$, voir
Catégories, Lemme \ref{categories-lemma-inertia-fibred-category}.
\end{proof}






\section{Champification des catégories fibrées en groupoïdes}
\label{section-stackify-groupoids}

\noindent
Voici le résultat.

\begin{lemma}
\label{lemma-stackify-groupoids}
Soit $\mathcal{C}$ un site.
Soit $p : \mathcal{S} \to \mathcal{C}$ une catégorie
fibrée en groupoïdes sur $\mathcal{C}$.
Il existe un champ en groupoïdes
$p' : \mathcal{S}' \to \mathcal{C}$ et un
$1$-morphisme $G : \mathcal{S} \to \mathcal{S}'$
de catégories fibrées en groupoïdes sur $\mathcal{C}$ (voir
Catégories, Définition
\ref{categories-definition-categories-fibred-in-groupoids-over-C})
tels que
\begin{enumerate}
\item pour tout $U \in \Ob(\mathcal{C})$ et tous
$x, y \in \Ob(\mathcal{S}_U)$, l'application
$$
\mathit{Mor}(x, y) \longrightarrow \mathit{Mor}(G(x), G(y))
$$
induite par $G$ identifie le membre de droite à la faisceautisation
du membre de gauche, et
\item pour tout $U \in \Ob(\mathcal{C})$ et tout
$x' \in \Ob(\mathcal{S}'_U)$, il existe un recouvrement
$\{U_i \to U\}_{i \in I}$ tel que, pour tout $i \in I$,
l'objet $x'|_{U_i}$ appartienne à l'image essentielle du
foncteur $G : \mathcal{S}_{U_i} \to \mathcal{S}'_{U_i}$.
\end{enumerate}
De plus, ces conditions déterminent le champ en groupoïdes $\mathcal{S}'$
à un unique $2$-isomorphisme près.
\end{lemma}

\begin{proof}
Appliquons le Lemme \ref{lemma-stackify}. Le résultat est un
champ en groupoïdes par le Lemme \ref{lemma-stack-in-groupoids-stack}.
\end{proof}

\begin{lemma}
\label{lemma-stackify-groupoids-universal-property}
Soit $\mathcal{C}$ un site.
Soit $p : \mathcal{S} \to \mathcal{C}$ une catégorie fibrée en groupoïdes
sur $\mathcal{C}$. Soient $p' : \mathcal{S}' \to \mathcal{C}$ et
$G : \mathcal{S} \to \mathcal{S}'$
respectivement le champ en groupoïdes et le $1$-morphisme construits dans le
Lemme \ref{lemma-stackify-groupoids}.
Cette construction possède la propriété universelle suivante : étant donnés un champ
en groupoïdes $q : \mathcal{X} \to \mathcal{C}$ et un $1$-morphisme
$F : \mathcal{S} \to \mathcal{X}$ de catégories sur $\mathcal{C}$,
il existe un $1$-morphisme $H : \mathcal{S}' \to \mathcal{X}$
tel que le diagramme
$$
\xymatrix{
\mathcal{S} \ar[rr]_F \ar[rd]_G & & \mathcal{X} \\
& \mathcal{S}' \ar[ru]_H
}
$$
soit $2$-commutatif.
\end{lemma}

\begin{proof}
C'est un cas particulier du
Lemme \ref{lemma-stackify-universal-property}.
\end{proof}

\begin{lemma}
\label{lemma-stackification-fibre-product-categories-fibred-in-groupoids}
Soit $\mathcal{C}$ un site.
Soient $f : \mathcal{X} \to \mathcal{Y}$ et $g : \mathcal{Z} \to \mathcal{Y}$
des morphismes de catégories fibrées en groupoïdes sur $\mathcal{C}$.
Alors la champification du $2$-produit fibré est le $2$-produit
fibré des champifications.
\end{lemma}

\begin{proof}
C'est un cas particulier du
Lemme \ref{lemma-stackification-fibre-product-fibred-categories}.
\end{proof}





\section{Topologies héritées}
\label{section-topology}

\noindent
Il s'avère qu'une catégorie fibrée sur un site hérite d'une topologie
canonique du site sous-jacent.

\begin{lemma}
\label{lemma-topology-inherited}
Soit $\mathcal{C}$ un site. Soit $p : \mathcal{S} \to \mathcal{C}$
une catégorie fibrée. Soit $\text{Cov}(\mathcal{S})$
l'ensemble des familles $\{x_i \to x\}_{i \in I}$ de morphismes de $\mathcal{S}$
de but fixé telles que (a) chaque $x_i \to x$ soit fortement cartésien
et (b) $\{p(x_i) \to p(x)\}_{i \in I}$ soit un recouvrement de $\mathcal{C}$.
Alors $(\mathcal{S}, \text{Cov}(\mathcal{S}))$ est un site.
\end{lemma}

\begin{proof}
Il faut vérifier les trois conditions de
Sites, Définition \ref{sites-definition-site}.
\begin{enumerate}
\item Si $x \to y$ est un isomorphisme de $\mathcal{S}$, alors
il est fortement cartésien d'après
Catégories, Lemme \ref{categories-lemma-composition-cartesian},
et $p(x) \to p(y)$ est un isomorphisme de $\mathcal{C}$. Par conséquent,
$\{p(x) \to p(y)\}$ est un recouvrement de $\mathcal{C}$, d'où
$\{x \to y\} \in \text{Cov}(\mathcal{S})$.
\item Si $\{x_i \to x\}_{i\in I} \in \text{Cov}(\mathcal{S})$ et si, pour tout
$i$, on a $\{y_{ij} \to x_i\}_{j\in J_i} \in \text{Cov}(\mathcal{S})$, alors
chaque composée $y_{ij} \to x$ est fortement cartésienne d'après
Catégories, Lemme \ref{categories-lemma-composition-cartesian},
et $\{p(y_{ij}) \to p(x)\}_{i \in I, j\in J_i} \in \text{Cov}(\mathcal{C})$.
On a donc aussi $\{y_{ij} \to x\}_{i \in I, j\in J_i} \in \text{Cov}(\mathcal{S})$.
\item Supposons $\{x_i \to x\}_{i\in I}\in \text{Cov}(\mathcal{S})$
et soit $y \to x$ un morphisme de $\mathcal{S}$. Puisque $\{p(x_i) \to p(x)\}$
est un recouvrement de $\mathcal{C}$, le produit fibré
$p(x_i) \times_{p(x)} p(y)$ existe. Par conséquent,
Catégories, Lemme \ref{categories-lemma-fibred-category-representable-goes-up}
implique que $x_i \times_x y$ existe, que $p(x_i \times_x y) =
p(x_i) \times_{p(x)} p(y)$ et que $x_i \times_x y \to y$ est
fortement cartésien. Comme, de plus,
$\{p(x_i) \times_{p(x)} p(y) \to p(y) \}_{i\in I} \in \text{Cov}(\mathcal{C})$,
on en conclut que
$\{x_i \times_x y \to y \}_{i\in I} \in \text{Cov}(\mathcal{S})$.
\end{enumerate}
Cela achève la démonstration.
\end{proof}

\noindent
Notons que, si $p : \mathcal{S} \to \mathcal{C}$ est fibré en groupoïdes, alors
les recouvrements du site $\mathcal{S}$ du
Lemme \ref{lemma-topology-inherited}
se caractérisent par
$$
\{x_i \to x\} \in \text{Cov}(\mathcal{S})
\Leftrightarrow
\{p(x_i) \to p(x)\} \in \text{Cov}(\mathcal{C})
$$
car tout morphisme de $\mathcal{S}$ est fortement cartésien.

\begin{definition}
\label{definition-topology-inherited}
Soit $\mathcal{C}$ un site. Soit $p : \mathcal{S} \to \mathcal{C}$ une
catégorie fibrée. On dit que $(\mathcal{S}, \text{Cov}(\mathcal{S}))$, comme dans le
Lemme \ref{lemma-topology-inherited},
est la {\it structure de site sur $\mathcal{S}$ héritée de $\mathcal{C}$}.
On l'exprime parfois en disant que
{\it $\mathcal{S}$ est munie de la topologie héritée de $\mathcal{C}$}.
\end{definition}

\noindent
On obtient en particulier, dans cette situation, un topos de faisceaux
$\Sh(\mathcal{S})$. Il s'avère que ce topos est fonctoriel par rapport
aux $1$-morphismes de catégories fibrées.

\begin{lemma}
\label{lemma-topology-inherited-functorial}
Soit $\mathcal{C}$ un site. Soit $F : \mathcal{X} \to \mathcal{Y}$
un $1$-morphisme de catégories fibrées sur $\mathcal{C}$.
Alors $F$ est un foncteur continu et cocontinu entre les structures
de site héritées de $\mathcal{C}$. Ainsi, $F$ induit un morphisme de topos
$f : \Sh(\mathcal{X}) \to \Sh(\mathcal{Y})$ avec
$f_* = {}_sF = {}_pF$ et $f^{-1} = F^s = F^p$. En particulier,
$f^{-1}(\mathcal{G})(x) = \mathcal{G}(F(x))$
pour tout faisceau $\mathcal{G}$ sur $\mathcal{Y}$ et tout objet $x$ de $\mathcal{X}$.
\end{lemma}

\begin{proof}
Montrons d'abord que $F$ est continu.
Soit $\{x_i \to x\}_{i \in I}$ un recouvrement de $\mathcal{X}$. D'après
Catégories, Définition \ref{categories-definition-fibred-categories-over-C},
le foncteur $F$ transforme les morphismes fortement cartésiens en morphismes
fortement cartésiens ; ainsi, $\{F(x_i) \to F(x)\}_{i \in I}$ est un recouvrement
de $\mathcal{Y}$. Cela prouve la partie (1) de
Sites, Définition \ref{sites-definition-continuous}.
De plus, soit $x' \to x$ un morphisme de $\mathcal{X}$. D'après
Catégories, Lemme \ref{categories-lemma-fibred-category-representable-goes-up},
le produit fibré $x_i \times_x x'$ existe et $x_i \times_x x' \to x'$
est fortement cartésien. Ainsi, $F(x_i \times_x x') \to F(x')$ est fortement
cartésien. D'après
Catégories, Lemme \ref{categories-lemma-fibred-category-representable-goes-up},
appliqué à $\mathcal{Y}$, cela signifie que
$F(x_i \times_x x') = F(x_i) \times_{F(x)} F(x')$.
Cela prouve la partie (2) de
Sites, Définition \ref{sites-definition-continuous},
et l'on en conclut que $F$ est continu.

\medskip\noindent
Montrons ensuite que $F$ est cocontinu.
Soit $x \in \Ob(\mathcal{X})$ et soit $\{y_i \to F(x)\}_{i \in I}$
un recouvrement de $\mathcal{Y}$. Notons $\{U_i \to U\}_{i \in I}$ le
recouvrement correspondant de $\mathcal{C}$. Pour chaque $i$, choisissons un morphisme
fortement cartésien $x_i \to x$ de $\mathcal{X}$ au-dessus de $U_i \to U$.
Alors $F(x_i) \to F(x)$ et $y_i \to F(x)$ sont tous deux des morphismes
fortement cartésiens de $\mathcal{Y}$ au-dessus de $U_i \to U$. Il existe donc
un unique isomorphisme $F(x_i) \to y_i$ dans $\mathcal{Y}_{U_i}$ compatible
avec les applications vers $F(x)$. Ainsi, $\{x_i \to x\}_{i \in I}$ est un recouvrement de
$\mathcal{X}$ tel que $\{F(x_i) \to F(x)\}_{i \in I}$ soit isomorphe à
$\{y_i \to F(x)\}_{i \in I}$. Le foncteur $F$ est donc cocontinu, voir
Sites, Définition \ref{sites-definition-cocontinuous}.

\medskip\noindent
L'assertion finale découle des deux premières, voir
Sites, Lemmes
\ref{sites-lemma-cocontinuous-morphism-topoi},
\ref{sites-lemma-pu-sheaf} et
\ref{sites-lemma-when-shriek}.
\end{proof}

\begin{lemma}
\label{lemma-localizing}
Soit $\mathcal{C}$ un site. Soit $p : \mathcal{X} \to \mathcal{C}$
une catégorie fibrée en groupoïdes. Soit $x \in \Ob(\mathcal{X})$
au-dessus de $U = p(x)$. Le foncteur $p$ induit une équivalence de sites
$\mathcal{X}/x \to \mathcal{C}/U$, où $\mathcal{X}$ est munie de
la topologie héritée de $\mathcal{C}$.
\end{lemma}

\begin{proof}
Ici, $\mathcal{C}/U$ est la localisation du site
$\mathcal{C}$ en l'objet $U$, et de même pour $\mathcal{X}/x$.
Il résulte de
Catégories, Définition \ref{categories-definition-fibred-groupoids},
que la règle $x'/x \mapsto p(x')/p(x)$ définit une équivalence de
catégories $\mathcal{X}/x \to \mathcal{C}/U$. Il résulte alors de la
Définition \ref{definition-topology-inherited}
que les recouvrements de $x'$ dans $\mathcal{X}/x$ correspondent bijectivement
aux recouvrements de $p(x')$ dans $\mathcal{C}/U$.
\end{proof}

\begin{lemma}
\label{lemma-stack-in-groupoids-over-stack-in-groupoids}
Soit $\mathcal{C}$ un site. Soient $p : \mathcal{X} \to \mathcal{C}$
et $q : \mathcal{Y} \to \mathcal{C}$
des champs en groupoïdes. Soit $F : \mathcal{X} \to \mathcal{Y}$
un $1$-morphisme de catégories sur $\mathcal{C}$. Si $F$ fait de
$\mathcal{X}$ une catégorie fibrée en groupoïdes sur $\mathcal{Y}$,
alors $\mathcal{X}$ est un champ en groupoïdes sur $\mathcal{Y}$ (pour la
topologie héritée de $\mathcal{C}$).
\end{lemma}

\begin{proof}
Montrons la descente des objets. Soit $\{y_i \to y\}$ un recouvrement
de $\mathcal{Y}$. Soit $(x_i, \varphi_{ij})$ une donnée de descente dans
$\mathcal{X}$ relativement à ce recouvrement. Alors $(x_i, \varphi_{ij})$ est aussi une donnée
de descente relativement au recouvrement $\{q(y_i) \to q(y)\}$ de
$\mathcal{C}$. Comme $\mathcal{X}$ est un champ en groupoïdes, on obtient un
objet $x$ au-dessus de $q(y)$ et des isomorphismes
$\psi_i : x|_{q(y_i)} \to x_i$
au-dessus de $q(y_i)$, compatibles avec les $\varphi_{ij}$, c'est-à-dire tels que
$$
\varphi_{ij} = \psi_j|_{q(y_i) \times_{q(y)} q(y_j)}
\circ \psi_i^{-1}|_{q(y_i) \times_{q(y)} q(y_j)}.
$$
Considérons le faisceau $\mathit{I} = \mathit{Isom}_\mathcal{Y}(F(x), y)$
sur $\mathcal{C}/p(x)$. Notons que $s_i = F(\psi_i) \in \mathit{I}(q(x_i))$
car $F(x_i) = y_i$. Puisque $F(\varphi_{ij}) = \text{id}$ (car nous sommes partis
d'une donnée de descente sur $\{y_i \to y\}$), la formule affichée montre
que $s_i|_{q(y_i) \times_{q(y)} q(y_j)} = s_j|_{q(y_i) \times_{q(y)} q(y_j)}$.
Les sections locales $s_i$ se recollent donc en $s : F(x) \to y$. Comme $F$ est fibré en
groupoïdes, on voit que $x$ est isomorphe à un objet $x'$ tel que $F(x') = y$.
Nous omettons de vérifier que $x'$, dans la catégorie fibre de $\mathcal{X}$
au-dessus de $y$, résout le problème de descente posé par la donnée de descente
$(x_i, \varphi_{ij})$. Nous omettons aussi la démonstration de la propriété de faisceau
des préfaisceaux $\mathit{Isom}$ de $\mathcal{X}/\mathcal{Y}$.
\end{proof}

\begin{lemma}
\label{lemma-stack-over-stack}
Soit $\mathcal{C}$ un site. Soit $p : \mathcal{X} \to \mathcal{C}$
un champ. Munissons $\mathcal{X}$ de la topologie héritée de
$\mathcal{C}$ et soit $q : \mathcal{Y} \to \mathcal{X}$ un champ.
Alors $\mathcal{Y}$ est un champ sur $\mathcal{C}$.
Si $p$ et $q$ définissent des champs en groupoïdes, alors
$\mathcal{Y}$ est un champ en groupoïdes sur $\mathcal{C}$.
\end{lemma}

\begin{proof}
Vérifions les trois conditions de la Définition \ref{definition-stack}
pour montrer que $\mathcal{Y}$ est un champ sur $\mathcal{C}$.
D'après Catégories, Lemme \ref{categories-lemma-fibred-over-fibred},
$\mathcal{Y}$ est une catégorie fibrée sur $\mathcal{C}$.
La condition (1) est donc satisfaite.

\medskip\noindent
Soit $U$ un objet de $\mathcal{C}$ et soient $y_1, y_2$ des objets
de $\mathcal{Y}$ au-dessus de $U$. Notons $x_i = q(y_i)$ dans $\mathcal{X}$.
Considérons l'application de préfaisceaux
$$
q : \mathit{Mor}_{\mathcal{Y}/\mathcal{C}}(y_1, y_2)
\longrightarrow
\mathit{Mor}_{\mathcal{X}/\mathcal{C}}(x_1, x_2)
$$
sur $\mathcal{C}/U$, voir le Lemme \ref{lemma-presheaf-mor-map-fibred-categories}.
Soit $\{U_i \to U\}$ un recouvrement et soient les $\varphi_i$ des sections
du préfaisceau de gauche sur les $U_i$, telles que $\varphi_i$ et
$\varphi_j$ aient la même restriction sur $U_i \times_U U_j$.
Il faut trouver un morphisme $\varphi : y_1 \to y_2$ dont la restriction soit $\varphi_i$.
Notons que $q(\varphi_i) = \psi|_{U_i}$ pour un certain morphisme
$\psi : x_1 \to x_2$ au-dessus de $U$, car le second préfaisceau est un faisceau
(par hypothèse). Soit $y_{12} \to y_2$ le morphisme fortement
$\mathcal{X}$-cartésien de $\mathcal{Y}$ au-dessus de $\psi$. Alors $\varphi_i$ correspond
à un morphisme $\varphi'_i : y_1|_{U_i} \to y_{12}|_{U_i}$ au-dessus de $x_1|_{U_i}$.
Autrement dit, les $\varphi'_i$ définissent maintenant des sections locales du préfaisceau
$$
\mathit{Mor}_{\mathcal{Y}/\mathcal{X}}(y_1, y_{12})
$$
sur les membres du recouvrement $\{x_1|_{U_i} \to x_1\}$. Par hypothèse, celles-ci
se recollent en un unique morphisme $y_1 \to y_{12}$ qui, composé avec le morphisme
donné $y_{12} \to y_2$, fournit le morphisme recherché $y_1 \to y_2$.

\medskip\noindent
Enfin, montrons que les données de descente sont effectives. Soit $\{f_i : U_i \to U\}$
un recouvrement de $\mathcal{C}$ et soit $(y_i, \varphi_{ij})$ une donnée de descente
relativement à ce recouvrement (Définition \ref{definition-descent-data}).
En posant $x_i = q(y_i)$ et $\psi_{ij} = q(\varphi_{ij})$,
on obtient une donnée de descente $(x_i, \psi_{ij})$ pour le recouvrement dans $\mathcal{X}$.
Par l'hypothèse sur $\mathcal{X}$, on peut supposer que $x_i = x|_{U_i}$
et que les $\psi_{ij}$ sont égaux à la donnée de descente canonique
(Définition \ref{definition-effective-descent-datum}).
Dans ce cas, $\{x|_{U_i} \to x\}$ est un recouvrement et l'on peut considérer
$(y_i, \varphi_{ij})$ comme une donnée de descente relativement à ce recouvrement.
Notre hypothèse selon laquelle $\mathcal{Y}$ est un champ sur $\mathcal{C}$
montre qu'elle est effective, ce qui achève la démonstration de la condition (3).

\medskip\noindent
L'assertion finale résulte de ce que $\mathcal{Y}$ est un champ sur
$\mathcal{C}$ et est fibrée en groupoïdes d'après
Catégories, Lemme
\ref{categories-lemma-fibred-in-groupoids-over-fibred-in-groupoids}.
\end{proof}





\section{Gerbes}
\label{section-gerbes}

\noindent
Les gerbes sont une classe particulière de champs en groupoïdes.

\begin{definition}
\label{definition-gerbe}
Une {\it gerbe} sur un site $\mathcal{C}$ est une catégorie
$p : \mathcal{S} \to \mathcal{C}$ sur $\mathcal{C}$ telle que
\begin{enumerate}
\item $p : \mathcal{S} \to \mathcal{C}$ soit un champ
en groupoïdes sur $\mathcal{C}$ (voir la
Définition \ref{definition-stack-in-groupoids}),
\item pour tout $U \in \Ob(\mathcal{C})$, il existe
un recouvrement $\{U_i \to U\}$ dans $\mathcal{C}$ tel que
$\mathcal{S}_{U_i}$ soit non vide, et
\item pour tout $U \in \Ob(\mathcal{C})$ et tous
$x, y \in \Ob(\mathcal{S}_U)$, il existe
un recouvrement $\{U_i \to U\}$ dans $\mathcal{C}$ tel que
$x|_{U_i} \cong y|_{U_i}$ dans $\mathcal{S}_{U_i}$.
\end{enumerate}
\end{definition}

\noindent
Autrement dit, une gerbe est un champ en groupoïdes dans lequel deux objets
quelconques sont localement isomorphes et où des objets existent localement.

\begin{lemma}
\label{lemma-gerbe-equivalent}
Soit $\mathcal{C}$ un site.
Soient $\mathcal{S}_1$ et $\mathcal{S}_2$ des catégories sur $\mathcal{C}$.
Supposons que $\mathcal{S}_1$ et $\mathcal{S}_2$ soient équivalentes
en tant que catégories sur $\mathcal{C}$.
Alors $\mathcal{S}_1$ est une gerbe sur $\mathcal{C}$ si et seulement si
$\mathcal{S}_2$ est une gerbe sur $\mathcal{C}$.
\end{lemma}

\begin{proof}
Supposons que $\mathcal{S}_1$ soit une gerbe sur $\mathcal{C}$. D'après le
Lemme \ref{lemma-stack-in-groupoids-equivalent},
$\mathcal{S}_2$ est un champ en groupoïdes sur $\mathcal{C}$.
Soient $F : \mathcal{S}_1 \to \mathcal{S}_2$ et
$G : \mathcal{S}_2 \to \mathcal{S}_1$ des équivalences de catégories sur
$\mathcal{C}$. Étant donné $U \in \Ob(\mathcal{C})$, il existe
un recouvrement $\{U_i \to U\}$ tel que $(\mathcal{S}_1)_{U_i}$ soit
non vide. En appliquant $F$, on voit que $(\mathcal{S}_2)_{U_i}$ est non vide.
Étant donnés $U \in \Ob(\mathcal{C})$ et
$x, y \in \Ob((\mathcal{S}_2)_U)$, il existe
un recouvrement $\{U_i \to U\}$ dans $\mathcal{C}$ tel que
$G(x)|_{U_i} \cong G(y)|_{U_i}$ dans $(\mathcal{S}_1)_{U_i}$. D'après
Catégories, Lemme \ref{categories-lemma-equivalence-fibred-categories},
cela implique que $x|_{U_i} \cong y|_{U_i}$ dans $(\mathcal{S}_2)_{U_i}$.
\end{proof}

\noindent
Nous voulons généraliser quelque peu la définition des gerbes. Soit donc
$F : \mathcal{X} \to \mathcal{Y}$ un $1$-morphisme de champs en groupoïdes
sur un site $\mathcal{C}$. Nous voulons préciser ce que signifie
pour $\mathcal{X}$ d'être une gerbe sur $\mathcal{Y}$. D'après la
Section \ref{section-topology},
la catégorie $\mathcal{Y}$ hérite de $\mathcal{C}$ d'une structure de site.
Une idée naïve serait d'exiger simplement que
$\mathcal{X} \to \mathcal{Y}$ {\it soit} une gerbe au sens ci-dessus. Toutefois,
la notion ainsi obtenue n'est pas invariante lorsque l'on remplace $\mathcal{X}$
par un champ en groupoïdes équivalent sur $\mathcal{C}$ ; c'est même le
cas de la propriété d'être fibrée en groupoïdes sur $\mathcal{Y}$.
Il s'avère cependant que l'on peut remplacer $\mathcal{X}$ par un champ
en groupoïdes équivalent sur $\mathcal{C}$ qui soit fibré en groupoïdes sur
$\mathcal{Y}$, et que la propriété d'être une gerbe sur $\mathcal{Y}$
est alors indépendante de ce choix. Voici l'énoncé précis.

\begin{lemma}
\label{lemma-when-gerbe}
Soit $\mathcal{C}$ un site. Soient $p : \mathcal{X} \to \mathcal{C}$
et $q : \mathcal{Y} \to \mathcal{C}$ des champs en groupoïdes.
Soit $F : \mathcal{X} \to \mathcal{Y}$ un $1$-morphisme de catégories
sur $\mathcal{C}$. Les assertions suivantes sont équivalentes :
\begin{enumerate}
\item Pour une certaine factorisation (de manière équivalente, pour toute factorisation)
$F = F' \circ a$ où
$a : \mathcal{X} \to \mathcal{X}'$ est une équivalence de catégories sur
$\mathcal{C}$ et où $F'$ est fibré en groupoïdes, l'application
$F' : \mathcal{X}' \to \mathcal{Y}$ est une gerbe (pour la topologie
sur $\mathcal{Y}$ héritée de $\mathcal{C}$).
\item Les deux conditions suivantes sont satisfaites :
\begin{enumerate}
\item pour tout $y \in \Ob(\mathcal{Y})$ au-dessus de
$U \in \Ob(\mathcal{C})$, il existe un recouvrement
$\{U_i \to U\}$ dans $\mathcal{C}$ et des objets $x_i$ de $\mathcal{X}$
au-dessus de $U_i$ tels que $F(x_i) \cong y|_{U_i}$ dans $\mathcal{Y}_{U_i}$, et
\item pour $U \in \Ob(\mathcal{C})$,
$x, x' \in \Ob(\mathcal{X}_U)$ et $b : F(x) \to F(x')$ dans
$\mathcal{Y}_U$, il existe
un recouvrement $\{U_i \to U\}$ dans $\mathcal{C}$ et des morphismes
$a_i : x|_{U_i} \to x'|_{U_i}$ dans $\mathcal{X}_{U_i}$ tels que
$F(a_i) = b|_{U_i}$.
\end{enumerate}
\end{enumerate}
\end{lemma}

\begin{proof}
D'après
Catégories, Lemme
\ref{categories-lemma-ameliorate-morphism-categories-fibred-groupoids},
il existe une factorisation $F = F' \circ a$ où
$a : \mathcal{X} \to \mathcal{X}'$ est une équivalence de catégories sur
$\mathcal{C}$ et où $F'$ est fibré en groupoïdes. D'après
Catégories, Lemme \ref{categories-lemma-amelioration-unique},
étant données deux telles factorisations $F = F' \circ a = F'' \circ b$,
on a une équivalence de $\mathcal{X}'$ et $\mathcal{X}''$ en tant que
catégories sur $\mathcal{Y}$. Par conséquent, le
Lemme \ref{lemma-gerbe-equivalent}
garantit que la condition (1) ne dépend pas du choix de la
factorisation. De plus, cela signifie que l'on peut supposer
$\mathcal{X}' = \mathcal{X} \times_{F, \mathcal{Y}, \text{id}} \mathcal{Y}$
comme dans la démonstration de
Catégories, Lemme
\ref{categories-lemma-ameliorate-morphism-categories-fibred-groupoids}.

\medskip\noindent
Montrons que (a) et (b) impliquent que $\mathcal{X}' \to \mathcal{Y}$
est une gerbe. Tout d'abord, d'après le
Lemme \ref{lemma-stack-in-groupoids-over-stack-in-groupoids},
$\mathcal{X}' \to \mathcal{Y}$ est un champ en groupoïdes.
Ensuite, soit $y$ un objet de $\mathcal{Y}$ au-dessus de
$U \in \Ob(\mathcal{C})$. D'après (a), on peut trouver un recouvrement
$\{U_i \to U\}$ dans $\mathcal{C}$, des objets $x_i$ de $\mathcal{X}$
au-dessus de $U_i$ et des isomorphismes $f_i : F(x_i) \to y|_{U_i}$ dans
$\mathcal{Y}_{U_i}$. Alors les $(U_i, x_i, y|_{U_i}, f_i)$ sont des objets
de $\mathcal{X}'_{U_i}$ ; la deuxième condition de la
Définition \ref{definition-gerbe}
est donc satisfaite. Enfin, soient $(U, x, y, f)$ et $(U, x', y, f')$ des objets
de $\mathcal{X}'$ au-dessus du même objet $y \in \Ob(\mathcal{Y})$.
Posons $b = (f')^{-1} \circ f$. D'après la condition (b), on peut trouver un recouvrement
$\{U_i \to U\}$ et des isomorphismes $a_i : x|_{U_i} \to x'|_{U_i}$
dans $\mathcal{X}_{U_i}$ tels que $F(a_i) = b|_{U_i}$. Alors
$$
(a_i, \text{id}) : (U, x, y, f)|_{U_i} \to (U, x', y, f')|_{U_i}
$$
est un morphisme dans $\mathcal{X}'_{U_i}$ comme voulu. Cela prouve que
(2) implique (1).

\medskip\noindent
Pour montrer que (1) implique (2), il suffit de lire à rebours les arguments
du paragraphe précédent. Détails omis.
\end{proof}

\begin{definition}
\label{definition-gerbe-over-stack-in-groupoids}
Soit $\mathcal{C}$ un site. Soient $\mathcal{X}$
et $\mathcal{Y}$ des champs en groupoïdes sur $\mathcal{C}$.
Soit $F : \mathcal{X} \to \mathcal{Y}$ un $1$-morphisme de catégories
sur $\mathcal{C}$. On dit que $\mathcal{X}$ est une {\it gerbe sur} $\mathcal{Y}$
si les conditions équivalentes du
Lemme \ref{lemma-when-gerbe}
sont satisfaites.
\end{definition}

\noindent
Cette définition n'entre pas en conflit avec la
Définition \ref{definition-gerbe}
lorsque $\mathcal{Y} = \mathcal{C}$, car on peut alors prendre
$\mathcal{X}' = \mathcal{X}$ dans la partie (1) du
Lemme \ref{lemma-when-gerbe}.
Notons que les conditions (2)(a) et (2)(b) du
Lemme \ref{lemma-when-gerbe}
sont très proches, dans leur esprit, des conditions (2) et (3) de la
Définition \ref{definition-gerbe}.
En effet, (2)(a) affirme que l'application des préfaisceaux de classes
d'isomorphisme d'objets devient une surjection après faisceautisation.
De plus, (2)(b) affirme que
$$
\mathit{Isom}_\mathcal{X}(x, x')
\longrightarrow
\mathit{Isom}_\mathcal{Y}(F(x), F(x'))
$$
est une surjection de faisceaux sur $\mathcal{C}/U$, pour tous $U$ et
$x, x' \in \Ob(\mathcal{X}_U)$.

\begin{lemma}
\label{lemma-base-change-gerbe}
Soit $\mathcal{C}$ un site. Soit
$$
\xymatrix{
\mathcal{X}' \ar[r]_{G'} \ar[d]_{F'} & \mathcal{X} \ar[d]^F \\
\mathcal{Y}' \ar[r]^G & \mathcal{Y}
}
$$
un $2$-produit fibré de champs en groupoïdes sur $\mathcal{C}$.
Si $\mathcal{X}$ est une gerbe sur $\mathcal{Y}$, alors
$\mathcal{X}'$ est une gerbe sur $\mathcal{Y}'$.
\end{lemma}

\begin{proof}
Par la propriété d'unicité d'un $2$-produit fibré, on peut supposer que
$\mathcal{X}' = \mathcal{Y}' \times_\mathcal{Y} \mathcal{X}$
comme dans
Catégories, Lemme \ref{categories-lemma-2-product-categories-over-C}.
Montrons les propriétés (2)(a) et (2)(b) du
Lemme \ref{lemma-when-gerbe}
pour $\mathcal{Y}' \times_\mathcal{Y} \mathcal{X} \to \mathcal{Y}'$.

\medskip\noindent
Soit $y'$ un objet de $\mathcal{Y}'$ au-dessus de l'objet $U$
de $\mathcal{C}$. Par hypothèse, il existe
un recouvrement $\{U_i \to U\}$ de $U$ et des objets
$x_i \in \mathcal{X}_{U_i}$ munis d'isomorphismes
$\alpha_i : G(y')|_{U_i} \to F(x_i)$.
Alors $(U_i, y'|_{U_i}, x_i, \alpha_i)$ est un objet de
$\mathcal{Y}' \times_\mathcal{Y} \mathcal{X}$ au-dessus de $U_i$
dont l'image dans $\mathcal{Y}'$ est $y'|_{U_i}$. La condition (2)(a) est donc satisfaite.

\medskip\noindent
Soit $U \in \Ob(\mathcal{C})$, soient $x'_1, x'_2$ des objets de
$\mathcal{Y}' \times_\mathcal{Y} \mathcal{X}$ au-dessus de $U$, et soit
$b' : F'(x'_1) \to F'(x'_2)$ un morphisme dans $\mathcal{Y}'_U$.
Écrivons $x'_i = (U, y'_i, x_i, \alpha_i)$. Notons que $F'(x'_i) = x_i$
et $G'(x'_i) = y'_i$. Par hypothèse, il existe
un recouvrement $\{U_i \to U\}$ dans $\mathcal{C}$ et des morphismes
$a_i : x_1|_{U_i} \to x_2|_{U_i}$ dans $\mathcal{X}_{U_i}$ tels que
$F(a_i) = G(b')|_{U_i}$. Alors $(b'|_{U_i}, a_i)$ est un morphisme
$x'_1|_{U_i} \to x'_2|_{U_i}$ comme l'exige (2)(b).
\end{proof}

\begin{lemma}
\label{lemma-composition-gerbe}
Soit $\mathcal{C}$ un site. Soient
$F : \mathcal{X} \to \mathcal{Y}$ et $G : \mathcal{Y} \to \mathcal{Z}$
des $1$-morphismes de champs en groupoïdes sur $\mathcal{C}$.
Si $\mathcal{X}$ est une gerbe sur $\mathcal{Y}$ et si
$\mathcal{Y}$ est une gerbe sur $\mathcal{Z}$, alors
$\mathcal{X}$ est une gerbe sur $\mathcal{Z}$.
\end{lemma}

\begin{proof}
Montrons les propriétés (2)(a) et (2)(b) du
Lemme \ref{lemma-when-gerbe}
pour $\mathcal{X} \to \mathcal{Z}$.

\medskip\noindent
Soit $z$ un objet de $\mathcal{Z}$ au-dessus de l'objet $U$
de $\mathcal{C}$. Par l'hypothèse sur $G$, il existe
un recouvrement $\{U_i \to U\}$ de $U$ et des objets
$y_i \in \mathcal{Y}_{U_i}$ tels que $G(y_i) \cong z|_{U_i}$.
Par l'hypothèse sur $F$, il existe des recouvrements $\{U_{ij} \to U_i\}$
et des objets $x_{ij} \in \mathcal{X}_{U_{ij}}$ tels que
$F(x_{ij}) \cong y_i|_{U_{ij}}$.
Alors $\{U_{ij} \to U\}$ est un recouvrement de $\mathcal{C}$
et $(G \circ F)(x_{ij}) \cong z|_{U_{ij}}$. La condition (2)(a) est donc satisfaite.

\medskip\noindent
Soit $U \in \Ob(\mathcal{C})$, soient $x_1, x_2$ des objets de
$\mathcal{X}$ au-dessus de $U$, et soit
$c : (G \circ F)(x_1) \to (G \circ F)(x_2)$ un morphisme dans
$\mathcal{Z}_U$. Par l'hypothèse sur $G$, il existe
un recouvrement $\{U_i \to U\}$ de $U$ et des morphismes
$b_i : F(x_1)|_{U_i} \to F(x_2)|_{U_i}$ dans $\mathcal{Y}_{U_i}$
tels que $G(b_i) = c|_{U_i}$.
Par l'hypothèse sur $F$, il existe des recouvrements $\{U_{ij} \to U_i\}$
et des morphismes $a_{ij} : x_1|_{U_{ij}} \to x_2|_{U_{ij}}$ dans
$\mathcal{X}_{U_{ij}}$ tels que $F(a_{ij}) = b_i|_{U_{ij}}$.
Alors $\{U_{ij} \to U\}$ est un recouvrement de $\mathcal{C}$
et $(G \circ F)(a_{ij}) = c|_{U_{ij}}$, comme l'exige (2)(b).
\end{proof}

\begin{lemma}
\label{lemma-gerbe-descent}
Soit $\mathcal{C}$ un site. Soit
$$
\xymatrix{
\mathcal{X}' \ar[r]_{G'} \ar[d]_{F'} & \mathcal{X} \ar[d]^F \\
\mathcal{Y}' \ar[r]^G & \mathcal{Y}
}
$$
un diagramme $2$-cartésien de champs en groupoïdes sur $\mathcal{C}$.
Si, pour tout $U \in \Ob(\mathcal{C})$ et tout
$x \in \Ob(\mathcal{Y}_U)$, il existe un recouvrement
$\{U_i \to U\}$ tel que $x|_{U_i}$ appartienne à l'image
essentielle de $G : \mathcal{Y}'_{U_i} \to \mathcal{Y}_{U_i}$ et si
$\mathcal{X}'$ est une gerbe sur $\mathcal{Y}'$, alors
$\mathcal{X}$ est une gerbe sur $\mathcal{Y}$.
\end{lemma}

\begin{proof}
Par la propriété d'unicité d'un $2$-produit fibré, on peut supposer que
$\mathcal{X}' = \mathcal{Y}' \times_\mathcal{Y} \mathcal{X}$
comme dans
Catégories, Lemme \ref{categories-lemma-2-product-categories-over-C}.
Montrons les propriétés (2)(a) et (2)(b) du
Lemme \ref{lemma-when-gerbe}
pour $\mathcal{X} \to \mathcal{Y}$.

\medskip\noindent
Soit $y$ un objet de $\mathcal{Y}$ au-dessus de l'objet $U$
de $\mathcal{C}$. Par hypothèse, il existe
un recouvrement $\{U_i \to U\}$ de $U$ et des objets
$y'_i \in \mathcal{Y}'_{U_i}$ tels que $G(y'_i) \cong y|_{U_i}$.
D'après (2)(a) pour $\mathcal{X}' \to \mathcal{Y}'$, il existe des
recouvrements $\{U_{ij} \to U_i\}$ et des objets $x'_{ij}$ de
$\mathcal{X}'$ au-dessus de $U_{ij}$ tels que $F'(x'_{ij})$ soit isomorphe
à la restriction de $y'_i$ à $U_{ij}$. Alors
$\{U_{ij} \to U\}$ est un recouvrement de $\mathcal{C}$ et les
$G'(x'_{ij})$ sont des objets de $\mathcal{X}$ au-dessus de $U_{ij}$
dont les images dans $\mathcal{Y}$ sont isomorphes aux restrictions
$y|_{U_{ij}}$. Cela prouve (2)(a) pour $\mathcal{X} \to \mathcal{Y}$.

\medskip\noindent
Soit $U \in \Ob(\mathcal{C})$, soient $x_1, x_2$ des objets de
$\mathcal{X}$ au-dessus de $U$, et soit $b : F(x_1) \to F(x_2)$ un morphisme
dans $\mathcal{Y}_U$. Par hypothèse, on peut choisir un recouvrement
$\{U_i \to U\}$ et des objets $y'_i$ de $\mathcal{Y}'$
au-dessus de $U_i$ tels qu'il existe des isomorphismes
$\alpha_i : G(y'_i) \to F(x_1)|_{U_i}$.
On obtient alors des objets
$$
x'_{1i} = (U_i, y'_i, x_1|_{U_i}, \alpha_i)
\quad\text{et}\quad
x'_{2i} = (U_i, y'_i, x_2|_{U_i}, b|_{U_i} \circ \alpha_i)
$$
de $\mathcal{X}'$ au-dessus de $U_i$. Le morphisme identité de $y'_i$ est un
morphisme $F'(x'_{1i}) \to F'(x'_{2i})$. D'après (2)(b) pour
$\mathcal{X}' \to \mathcal{Y}'$, il existe des recouvrements
$\{U_{ij} \to U_i\}$ et des morphismes
$a'_{ij} : x'_{1i}|_{U_{ij}} \to x'_{2i}|_{U_{ij}}$
tels que $F'(a'_{ij}) = \text{id}_{y'_i}|_{U_{ij}}$. En explicitant la définition
des morphismes de $\mathcal{Y}' \times_\mathcal{Y} \mathcal{X}$,
on voit que les $G'(a'_{ij}) : x_1|_{U_{ij}} \to x_2|_{U_{ij}}$ sont
les morphismes recherchés, c'est-à-dire que (2)(b) est satisfaite pour
$\mathcal{X} \to \mathcal{Y}$.
\end{proof}

\noindent
Les gerbes dont tous les faisceaux d'automorphismes sont abéliens jouent un rôle
important en géométrie algébrique.

\begin{lemma}
\label{lemma-gerbe-abelian-auts}
Soit $p : \mathcal{S} \to \mathcal{C}$ une gerbe sur un site $\mathcal{C}$.
Supposons que, pour tous $U \in \Ob(\mathcal{C})$ et $x \in \Ob(\mathcal{S}_U)$,
le faisceau de groupes $\mathit{Aut}(x) = \mathit{Isom}(x, x)$ sur $\mathcal{C}/U$
soit abélien. Alors il existe
\begin{enumerate}
\item un faisceau $\mathcal{G}$ de groupes abéliens sur $\mathcal{C}$,
\item pour tout $U \in \Ob(\mathcal{C})$ et tout $x \in \Ob(\mathcal{S}_U)$,
un isomorphisme $\mathcal{G}|_U \to \mathit{Aut}(x)$
\end{enumerate}
tel que, pour tout $U$ et tout morphisme $\varphi : x \to y$
dans $\mathcal{S}_U$, le diagramme
$$
\xymatrix{
\mathcal{G}|_U \ar[d] \ar@{=}[rr] & & \mathcal{G}|_U \ar[d] \\
\mathit{Aut}(x)
\ar[rr]^{\alpha \mapsto \varphi \circ \alpha \circ \varphi^{-1}} & &
\mathit{Aut}(y)
}
$$
soit commutatif.
\end{lemma}

\begin{proof}
Soient $x, y$ deux objets de $\mathcal{S}$ avec $U = p(x) = p(y)$.

\medskip\noindent
S'il existe un morphisme $\varphi : x \to y$ au-dessus de $U$, c'est un
isomorphisme et l'on obtient bien un isomorphisme
$\mathit{Aut}(x) \to \mathit{Aut}(y)$ qui envoie
$\alpha$ sur $\varphi \circ \alpha \circ \varphi^{-1}$.
De plus, comme on suppose $\mathit{Aut}(x)$
commutatif, cet isomorphisme ne dépend pas du choix
de $\varphi$, par un calcul immédiat : en effet, si $\psi$ est
une seconde application de ce type, alors
$$
\varphi \circ \alpha \circ \varphi^{-1} =
\psi \circ \psi^{-1} \circ \varphi \circ \alpha \circ \varphi^{-1} =
\psi \circ \alpha \circ \psi^{-1} \circ \varphi \circ \varphi^{-1} =
\psi \circ \alpha \circ \psi^{-1}
$$
Il en résulte un isomorphisme canonique de faisceaux
$\mathit{Aut}(x) \to \mathit{Aut}(y)$. En outre, s'il existe
un troisième objet $z$ et un morphisme $y \to z$ (et donc
aussi un morphisme $x \to z$), alors les
isomorphismes canoniques $\mathit{Aut}(x) \to \mathit{Aut}(y)$,
$\mathit{Aut}(y) \to \mathit{Aut}(z)$ et
$\mathit{Aut}(x) \to \mathit{Aut}(z)$ sont compatibles au sens
où le diagramme
$$
\xymatrix{
\mathit{Aut}(x) \ar[rd] \ar[rr] & &  \mathit{Aut}(z) \\
& \mathit{Aut}(y) \ar[ru]
}
$$
est commutatif.

\medskip\noindent
S'il n'existe aucun morphisme de $x$ vers $y$ au-dessus de $U$, on peut
choisir un recouvrement $\{U_i \to U\}$ tel qu'il existe des morphismes
$x|_{U_i} \to y|_{U_i}$. On obtient ainsi des isomorphismes canoniques
$$
\mathit{Aut}(x)|_{U_i} \longrightarrow \mathit{Aut}(y)|_{U_i}
$$
qui coïncident sur $U_i \times_U U_j$ (par canonicité). Par recollement des
faisceaux (Sites, Lemme \ref{sites-lemma-glue-maps}), on obtient un unique
isomorphisme $\mathit{Aut}(x) \to \mathit{Aut}(y)$ dont la restriction
à chaque $U_i$ est l'isomorphisme canonique du paragraphe précédent.
Comme ci-dessus, ces isomorphismes canoniques satisfont une
compatibilité en présence d'un troisième objet au-dessus de $U$.

\medskip\noindent
Que se passe-t-il si la catégorie fibre de $\mathcal{S}$ au-dessus de $U$ est vide ?
Dans ce cas, on peut trouver un recouvrement $\{U_i \to U\}$
et des objets $x_i$ de $\mathcal{S}$ au-dessus de $U_i$. Posons alors
$\mathcal{G}_i = \mathit{Aut}(x_i)$. D'après ce qui précède, on obtient
des isomorphismes canoniques
$$
\varphi_{ij} :
\mathcal{G}_i|_{U_i \times_U U_j}
\longrightarrow
\mathcal{G}_j|_{U_i \times_U U_j}
$$
dont les restrictions à $U_i \times_U U_j \times_U U_k$ satisfont
la condition de cocycle décrite dans Sites, Section
\ref{sites-section-glueing-sheaves}.
D'après Sites, Lemme \ref{sites-lemma-glue-sheaves},
on obtient un faisceau $\mathcal{G}$ sur $U$ dont la
restriction à $U_i$ donne $\mathcal{G}_i$
d'une manière compatible avec les applications de recollement $\varphi_{ij}$.

\medskip\noindent
Si $\mathcal{C}$ possède un objet final $U$, cela achève la démonstration,
car on peut prendre pour $\mathcal{G}$ le faisceau que l'on vient de construire.
Dans le cas général, il faut vérifier que les faisceaux $\mathcal{G}$
construits sur les différents $U$ sont canoniquement compatibles.
Cette vérification est omise.
\end{proof}




\section{Fonctorialité des champs}
\label{section-inverse-image}

\noindent
Dans cette section, nous étudions ce qui se passe lorsque l'on veut changer
le site de base d'un champ. Cette section peut être omise en première lecture.

\medskip\noindent
Soit $u : \mathcal{C} \to \mathcal{D}$ un foncteur entre catégories.
Soit $p : \mathcal{S} \to \mathcal{D}$ une catégorie au-dessus de $\mathcal{D}$.
Dans cette situation, nous notons $u^p\mathcal{S}$ la catégorie au-dessus de $\mathcal{C}$
définie comme suit :
\begin{enumerate}
\item Un objet de $u^p\mathcal{S}$ est un couple $(U, y)$ formé
d'un objet $U$ de $\mathcal{C}$ et d'un objet $y$ de $\mathcal{S}_{u(U)}$.
\item Un morphisme $(a, \beta) : (U, y) \to (U', y')$ est donné par
un morphisme $a : U \to U'$ de $\mathcal{C}$ et un morphisme $\beta : y \to y'$
de $\mathcal{S}$ tels que $p(\beta) = u(a)$.
\end{enumerate}
Remarquons qu'avec ces définitions, la catégorie fibre de
$u^p\mathcal{S}$ au-dessus de $U$ est égale à la catégorie fibre de
$\mathcal{S}$ au-dessus de $u(U)$.

\begin{lemma}
\label{lemma-fibred-category-pushforward}
Dans la situation ci-dessus, si $\mathcal{S}$ est une catégorie fibrée au-dessus de
$\mathcal{D}$, alors $u^p\mathcal{S}$ est une catégorie fibrée au-dessus de $\mathcal{C}$.
\end{lemma}

\begin{proof}
Avant de lire cette démonstration, prière de se reporter à la discussion qui entoure
Catégories, Définitions \ref{categories-definition-cartesian-over-C} et
\ref{categories-definition-fibred-category}.
Soit $(a, \beta) : (U, y) \to (U', y')$ un morphisme de $u^p\mathcal{S}$.
Nous affirmons que $(a, \beta)$ est fortement cartésien si et seulement si
$\beta$ est fortement cartésien. Supposons d'abord que $\beta$ soit fortement cartésien.
Considérons un autre morphisme quelconque
$(a_1, \beta_1) : (U_1, y_1) \to (U', y')$ de $u^p\mathcal{S}$. Alors
\begin{align*}
& \Mor_{u^p\mathcal{S}}((U_1, y_1), (U, y)) \\
& =
\Mor_\mathcal{C}(U_1, U)
\times_{\Mor_\mathcal{D}(u(U_1), u(U))}
\Mor_\mathcal{S}(y_1, y) \\
& =
\Mor_\mathcal{C}(U_1, U)
\times_{\Mor_\mathcal{D}(u(U_1), u(U))}
\Mor_\mathcal{S}(y_1, y')
\times_{\Mor_\mathcal{D}(u(U_1), u(U'))}
\Mor_\mathcal{D}(u(U_1), u(U)) \\
& =
\Mor_\mathcal{S}(y_1, y')
\times_{\Mor_\mathcal{D}(u(U_1), u(U'))}
\Mor_\mathcal{C}(U_1, U) \\
& =
\Mor_{u^p\mathcal{S}}((U_1, y_1), (U', y'))
\times_{\Mor_\mathcal{C}(U_1, U')}
\Mor_\mathcal{C}(U_1, U)
\end{align*}
la deuxième égalité résultant de ce que $\beta$ est fortement cartésien. Nous voyons donc
que $(a, \beta)$ est bien fortement cartésien. Réciproquement, supposons que
$(a, \beta)$ soit fortement cartésien. Choisissons un morphisme fortement cartésien
$\beta' : y'' \to y'$ dans $\mathcal{S}$ tel que $p(\beta') = u(a)$.
Alors les deux morphismes $(a, \beta) : (U, y) \to (U, y')$ et
$(a, \beta') : (U, y'') \to (U, y)$ sont fortement cartésiens et
relèvent $a$. Par l'unicité des morphismes fortement cartésiens
(voir la discussion dans
Catégories, Section \ref{categories-section-fibred-categories}),
il existe donc un isomorphisme $\iota : y \to y''$ dans $\mathcal{S}_{u(U)}$
tel que $\beta = \beta' \circ \iota$, ce qui implique que
$\beta$ est fortement cartésien dans $\mathcal{S}$ d'après
Catégories, Lemme \ref{categories-lemma-composition-cartesian}.

\medskip\noindent
Enfin, il nous faut montrer
qu'étant donnés $(U', y')$ et $U \to U'$, on peut trouver un morphisme
fortement cartésien $(U, y) \to (U', y')$ dans $u^p\mathcal{S}$
relevant le morphisme $U \to U'$. Cela résulte de ce qui précède puisque, par hypothèse,
on peut trouver un morphisme fortement cartésien $y \to y'$ relevant le
morphisme $u(U) \to u(U')$.
\end{proof}

\begin{lemma}
\label{lemma-stack-pushforward}
Soit $u : \mathcal{C} \to \mathcal{D}$ un foncteur continu de sites.
Soit $p : \mathcal{S} \to \mathcal{D}$ un champ sur $\mathcal{D}$.
Alors $u^p\mathcal{S}$ est un champ sur $\mathcal{C}$.
\end{lemma}

\begin{proof}
Nous avons vu dans le
Lemme \ref{lemma-fibred-category-pushforward}
que $u^p\mathcal{S}$ est une catégorie fibrée au-dessus de $\mathcal{C}$.
De plus, dans la démonstration de ce lemme, nous avons vu qu'un morphisme
$(a, \beta)$ de $u^p\mathcal{S}$ est fortement cartésien si et seulement si
$\beta$ est fortement cartésien dans $\mathcal{S}$. Ainsi, étant donné
un morphisme $a : U \to U'$ de $\mathcal{C}$, non seulement
avons-nous les égalités
$(u^p\mathcal{S})_U = \mathcal{S}_U$ et
$(u^p\mathcal{S})_{U'} = \mathcal{S}_{U'}$, mais, par ces
égalités, les foncteurs de changement de base coïncident ; en formule,
$a^*(U', y') = (U, u(a)^*y')$.

\medskip\noindent
Ceci étant dit, soit $\mathcal{U} = \{U_i \to U\}$ un recouvrement
de $\mathcal{C}$. Comme $u$ est continu, nous voyons que
$\mathcal{V} = \{u(U_i) \to u(U)\}$ est un recouvrement de $\mathcal{D}$,
que $u(U_i \times_U U_j) = u(U_i) \times_{u(U)} u(U_j)$, et de même
pour les produits fibrés triples $U_i \times_U U_j \times_U U_k$.
Compte tenu des identifications des catégories fibres et des changements de base,
les données de descente relatives à $\mathcal{U}$ sont identiques
aux données de descente relatives à $\mathcal{V}$. Puisque, par hypothèse, la descente
est effective dans $\mathcal{S}$, nous concluons qu'il en est de même dans
$u^p\mathcal{S}$.
\end{proof}

\begin{lemma}
\label{lemma-stack-in-groupoids-pushforward}
Soit $u : \mathcal{C} \to \mathcal{D}$ un foncteur continu de sites.
Soit $p : \mathcal{S} \to \mathcal{D}$ un champ en groupoïdes
sur $\mathcal{D}$. Alors $u^p\mathcal{S}$ est un champ en groupoïdes
sur $\mathcal{C}$.
\end{lemma}

\begin{proof}
Cela résulte immédiatement du
Lemme \ref{lemma-stack-pushforward}
et du fait que toutes les catégories fibres sont des groupoïdes.
\end{proof}

\begin{definition}
\label{definition-pushforward-stack}
Soit $f : \mathcal{D} \to \mathcal{C}$ un morphisme de sites
donné par le foncteur continu $u : \mathcal{C} \to \mathcal{D}$.
Soit $\mathcal{S}$ une catégorie fibrée au-dessus de $\mathcal{D}$.
Dans ce cadre, nous notons {\it $f_*\mathcal{S}$} la catégorie
fibrée $u^p\mathcal{S}$ définie ci-dessus. Nous disons que
$f_*\mathcal{S}$ est l'{\it image directe de $\mathcal{S}$ par $f$}.
\end{definition}

\noindent
D'après les résultats précédents, nous savons que $f_*\mathcal{S}$ est un champ (en groupoïdes)
si $\mathcal{S}$ est un champ (en groupoïdes).
Il est plus difficile de définir l'image inverse d'un champ (et notre construction
particulière nécessitera des hypothèses supplémentaires -- le lecteur peut rédiger et
soumettre une construction plus générale). Nous procédons en plusieurs étapes.

\medskip\noindent
Soit $u : \mathcal{C} \to \mathcal{D}$ un foncteur entre catégories.
Soit $p : \mathcal{S} \to \mathcal{C}$ une catégorie au-dessus de $\mathcal{C}$.
Dans ce cadre, nous définissons une catégorie $u_{pp}\mathcal{S}$ comme suit :
\begin{enumerate}
\item Un objet de $u_{pp}\mathcal{S}$ est un triplet
$(U, \phi : V \to u(U), x)$ où $U \in \Ob(\mathcal{C})$,
l'application $\phi : V \to u(U)$ est un morphisme de $\mathcal{D}$,
et $x \in \Ob(\mathcal{S}_U)$.
\item Un morphisme
$$
(U_1, \phi_1 : V_1 \to u(U_1), x_1)
\longrightarrow
(U_2, \phi_2 : V_2 \to u(U_2), x_2)
$$
de $u_{pp}\mathcal{S}$ est donné par un
triplet $(a, b, \alpha)$ où $a : U_1 \to U_2$ est un morphisme de $\mathcal{C}$,
$b : V_1 \to V_2$ est un morphisme de $\mathcal{D}$, et
$\alpha : x_1 \to x_2$ est un morphisme de $\mathcal{S}$,
tels que $p(\alpha) = a$ et que le diagramme
$$
\xymatrix{
V_1 \ar[d]_{\phi_1} \ar[r]_b & V_2 \ar[d]^{\phi_2} \\
u(U_1) \ar[r]^{u(a)} & u(U_2)
}
$$
soit commutatif dans $\mathcal{D}$.
\end{enumerate}
Nous considérons $u_{pp}\mathcal{S}$ comme une catégorie au-dessus de $\mathcal{D}$ au moyen de
$$
p_{pp} : u_{pp}\mathcal{S} \longrightarrow \mathcal{D},
\quad
(U, \phi : V \to u(U), x) \longmapsto V.
$$
La catégorie fibre de $u_{pp}\mathcal{S}$ au-dessus d'un objet $V$ de $\mathcal{D}$
n'admet pas de description simple.

\begin{lemma}
\label{lemma-right-multiplicative-system}
Dans la situation ci-dessus, supposons que
\begin{enumerate}
\item $p : \mathcal{S} \to \mathcal{C}$ soit une catégorie fibrée,
\item $\mathcal{C}$ possède les limites finies non vides, et que
\item $u : \mathcal{C} \to \mathcal{D}$ commute aux limites finies non vides.
\end{enumerate}
Considérons l'ensemble $R \subset \text{Flèches}(u_{pp}\mathcal{S})$ des morphismes
de la forme
$$
(a, \text{id}_V, \alpha) :
(U', \phi' : V \to u(U'), x')
\longrightarrow
(U, \phi : V \to u(U), x)
$$
où $\alpha$ est fortement cartésien. Alors $R$ est un système multiplicatif à droite.
\end{lemma}

\begin{proof}
D'après
Catégories, Définition \ref{categories-definition-multiplicative-system},
nous devons vérifier RMS1, RMS2 et RMS3.
La condition RMS1 est satisfaite, car une composée de morphismes fortement cartésiens est
fortement cartésienne ; voir
Catégories, Lemme \ref{categories-lemma-composition-cartesian}.

\medskip\noindent
Pour vérifier RMS2, supposons donnés un morphisme
$$
(a, b, \alpha) :
(U_1, \phi_1 : V_1 \to u(U_1), x_1)
\longrightarrow
(U, \phi : V \to u(U), x)
$$
de $u_{pp}\mathcal{S}$ et un morphisme
$$
(c, \text{id}_V, \gamma) :
(U', \phi' : V \to u(U'), x')
\longrightarrow
(U, \phi : V \to u(U), x)
$$
où $\gamma$ est fortement cartésien, appartenant à $R$. Dans cette situation, posons
$U'_1 = U_1 \times_U U'$, et notons $a' : U'_1 \to U'$ et
$c' : U'_1 \to U_1$ les projections.
Puisque $u(U'_1) = u(U_1) \times_{u(U)} u(U')$,
$\phi'_1 = (\phi_1, \phi') : V_1 \to u(U'_1)$ est
un morphisme de $\mathcal{D}$. Soit $\gamma_1 : x_1' \to x_1$
un morphisme fortement cartésien de $\mathcal{S}$ tel que
$p(\gamma_1) = \phi'_1$ (il existe puisque $\mathcal{S}$ est une
catégorie fibrée au-dessus de $\mathcal{C}$). Comme $\gamma : x' \to x$ est
fortement cartésien, il existe alors un unique morphisme
$\alpha' : x'_1 \to x'$ tel que $p(\alpha') = a'$.
À ce stade, nous voyons que
$$
(a', b, \alpha') :
(U_1, \phi_1 : V_1 \to u(U'_1), x'_1)
\longrightarrow
(U, \phi : V \to u(U'), x')
$$
est un morphisme et que
$$
(c', \text{id}_{V_1}, \gamma_1) :
(U'_1, \phi'_1 : V_1 \to u(U'_1), x'_1)
\longrightarrow
(U_1, \phi : V_1 \to u(U_1), x_1)
$$
est un élément de $R$ ; ensemble, ces morphismes donnent une solution au problème
d'existence posé par RMS2.

\medskip\noindent
Enfin, supposons que
$$
(a, b, \alpha), (a', b', \alpha') :
(U_1, \phi_1 : V_1 \to u(U_1), x_1)
\longrightarrow
(U, \phi : V \to u(U), x)
$$
soient deux morphismes de $u_{pp}\mathcal{S}$ et supposons que
$$
(c, \text{id}_V, \gamma) :
(U, \phi : V \to u(U), x)
\longrightarrow
(U', \phi : V \to u(U'), x')
$$
soit un élément de $R$ qui égalise les morphismes
$(a, b, \alpha)$ et $(a', b', \alpha')$. Cela implique en particulier
que $b = b'$. Soit $d : U_2 \to U_1$ l'égalisateur de $a, a'$, lequel
existe (voir
Catégories, Lemme \ref{categories-lemma-almost-finite-limits-exist}).
En outre, $u(d) : u(U_2) \to u(U_1)$ est l'égalisateur de $u(a), u(a')$ ;
il existe donc (puisque $b = b'$) un morphisme $\phi_2 : V_1 \to u(U_2)$ tel que
$\phi_1 = u(d) \circ \phi_1$. Soit $\delta : x_2 \to x_1$ un morphisme fortement
cartésien de $\mathcal{S}$ tel que $p(\delta) = u(d)$.
Nous affirmons maintenant que $\alpha \circ \delta = \alpha' \circ \delta$.
C'est vrai parce que
$\gamma$ est fortement cartésien,
$\gamma \circ \alpha \circ \delta = \gamma \circ \alpha' \circ \delta$, et
$p(\alpha \circ \delta) = p(\alpha' \circ \delta)$.
Par conséquent, la flèche
$$
(d, \text{id}_{V_1}, \delta) :
(U_2, \phi_2 : V_1 \to u(U_2), x_2)
\longrightarrow
(U_1, \phi_1 : V_1 \to u(U_1), x_1)
$$
est un élément de $R$ et égalise $(a, b, \alpha)$ et $(a', b', \alpha')$.
Ainsi, $R$ satisfait également RMS3.
\end{proof}

\begin{lemma}
\label{lemma-fibred-category-pullback}
Avec les notations et les hypothèses du
Lemme \ref{lemma-right-multiplicative-system},
posons $u_p\mathcal{S} = R^{-1}u_{pp}\mathcal{S}$ ; voir
Catégories, Section \ref{categories-section-localization}.
Alors $u_p\mathcal{S}$ est une catégorie fibrée au-dessus de $\mathcal{D}$.
\end{lemma}

\begin{proof}
Nous utilisons la description de $u_p\mathcal{S}$ donnée juste avant
Catégories, Lemme \ref{categories-lemma-right-localization}.
Remarquons que le foncteur $p_{pp} : u_{pp}\mathcal{S} \to \mathcal{D}$
transforme tout élément de $R$ en un morphisme identité.
Ainsi, d'après
Catégories, Lemme \ref{categories-lemma-properties-right-localization},
nous obtenons un foncteur canonique $p_p : u_p\mathcal{S} \to \mathcal{D}$
qui prolonge le foncteur donné. C'est ainsi que nous considérons
$u_p\mathcal{S}$ comme une catégorie au-dessus de $\mathcal{D}$.

\medskip\noindent
Nous voulons d'abord caractériser les morphismes fortement cartésiens au-dessus de $\mathcal{D}$
dans $u_p\mathcal{S}$.
Un morphisme $f : X \to Y$ de $u_p\mathcal{S}$ est la classe d'équivalence
d'un couple $(f' : X' \to Y, r : X' \to X)$ avec $r \in R$.
En effet, dans $u_p\mathcal{S}$, nous avons $f = (f', 1) \circ (r, 1)^{-1}$
avec les notations évidentes.
Remarquons qu'un isomorphisme est toujours fortement cartésien, de même qu'une
composée de morphismes fortement cartésiens ; voir
Catégories, Lemme \ref{categories-lemma-composition-cartesian}.
Ainsi, $f$ est fortement cartésien si et seulement si $(f', 1)$ l'est.
L'affirmation suivante caractérise donc complètement les morphismes fortement cartésiens.
Affirmation : un morphisme
$$
(a, b, \alpha) :
X_1 = (U_1, \phi_1 : V_1 \to u(U_1), x_1)
\longrightarrow
(U_2, \phi_2 : V_2 \to u(U_2), x_2) = X_2
$$
de $u_{pp}\mathcal{S}$ a une image $f = ((a, b, \alpha), 1)$
fortement cartésienne dans $u_p\mathcal{S}$ si et seulement si $\alpha$
est un morphisme fortement cartésien de $\mathcal{S}$.

\medskip\noindent
Supposons $\alpha$ fortement cartésien.
Soit $X = (U, \phi : V \to u(U), x)$ un autre objet, et soit
$f_2 : X \to X_2$ un morphisme de $u_p\mathcal{S}$
tel que $p_p(f_2) = b \circ b_1$ pour un certain $b_1 : U \to U_1$.
Pour montrer que $f$ est fortement cartésien, nous devons montrer
qu'il existe un unique morphisme $f_1 : X \to X_1$ dans $u_p\mathcal{S}$
tel que $p_p(f_1) = b_1$ et $f_2 = f \circ f_1$ dans $u_p\mathcal{S}$.
Écrivons $f_2 = (f'_2 : X' \to X_2, r : X' \to X)$. Là encore, nous pouvons écrire
$f_2 = (f'_2, 1) \circ (r, 1)^{-1}$ dans $u_p\mathcal{S}$. Puisque $(r, 1)$
est un isomorphisme dont l'image dans $\mathcal{D}$ est une identité, nous
voyons que trouver un morphisme $f_1 : X \to X_1$ possédant les propriétés
requises revient à trouver un morphisme $f'_1 : X' \to X_1$
dans $u_p\mathcal{S}$ tel que $p(f'_1) = b_1$ et $f'_2 = f \circ f'_1$.
Nous pouvons donc supposer que $f_2$ est de la forme
$f_2 = ((a_2, b_2, \alpha_2), 1)$ avec $b_2 = b \circ b_1$. Voici
un diagramme :
$$
\xymatrix{
& & (U_1, V_1 \to u(U_1), x_1) \ar[d]^{(a, b, \alpha)} \\
(U, V \to u(U), x) \ar[rr]^{(a_2, b_2, \alpha_2)} & &
(U_2, V_2 \to u(U_2), x_2)
}
$$
La construction du morphisme $f_1$ est maintenant claire.
Posons en effet $U' = U \times_{U_2} U_1$, avec pour projections
$c : U' \to U$ et $a_1 : U' \to U_1$.
Choisissons un morphisme fortement cartésien $\gamma : x' \to x$
relevant le morphisme $c$. Puisque $b_2 = b \circ b_1$
et que $u(U') = u(U) \times_{u(U_2)} u(U_1)$, nous avons
$\phi' = (\phi, \phi_1 \circ b_1) : V \to u(U')$. Comme
$\alpha$ est fortement cartésien et que
$a \circ a_1 = a_2 \circ c = p(\alpha_2 \circ \gamma)$,
il existe un morphisme
$\alpha_1 : x' \to x_1$ relevant $a_1$ tel que
$\alpha \circ \alpha_1 = \alpha_2 \circ \gamma$.
Posons $X' = (U', \phi' : V \to u(U'), x')$.
Nous voyons ainsi que
$$
f_1 =
((a_1, b_1, \alpha_1) : X' \to X_1, (c, \text{id}_V, \gamma) : X' \to X) :
X \longrightarrow X_1
$$
convient ; en effet, le diagramme
$$
\xymatrix{
(U', \phi' : V \to u(U'), x')
\ar[d]_{(c, \text{id}_V, \gamma)}
\ar[rr]_{(a_1, b_1, \alpha_1)}
& & (U_1, V_1 \to u(U_1), x_1) \ar[d]^{(a, b, \alpha)} \\
(U, V \to u(U), x) \ar[rr]^{(a_2, b_2, \alpha_2)} & &
(U_2, V_2 \to u(U_2), x_2)
}
$$
est commutatif par construction. Cela démontre l'existence.

\medskip\noindent
Démontrons ensuite l'unicité, toujours dans le cas particulier
$f = ((a, b, \alpha), 1)$ et $f_2 = ((a_2, b_2, \alpha_2), 1)$.
Nous recommandons vivement au lecteur de passer cette partie.
Supposons que $g_1, g'_1 : X \to X_1$ soient deux morphismes de
$u_p\mathcal{S}$ tels que $p_p(g_1) = p_p(g'_1) = b_1$ et
$f_2 = f \circ g_1 = f \circ g'_1$. Notre but est de montrer que
$g_1 = g'_1$. D'après
Catégories, Lemme \ref{categories-lemma-morphisms-right-localization},
nous pouvons représenter $g_1$ et $g'_1$ par les classes d'équivalence de
$(f_1 : X' \to X_1, r : X' \to X)$ et
$(f'_1 : X' \to X_1, r : X' \to X)$ pour un certain $r \in R$. D'après
Catégories, Lemme \ref{categories-lemma-equality-morphisms-right-localization},
l'égalité $f_2 = f \circ g_1 = f \circ g'_1$ signifie qu'il
existe un morphisme $r' : X'' \to X'$ dans $u_{pp}\mathcal{S}$ tel que
$r' \circ r \in R$ et
$$
(a, b, \alpha) \circ f_1 \circ r' =
(a, b, \alpha) \circ f'_1 \circ r' =
(a_2, b_2, \alpha_2) \circ r'
$$
dans $u_{pp}\mathcal{S}$. Remarquons que
$g_1$ est maintenant représenté par $(f_1 \circ r', r \circ r')$ et
de même pour $g'_1$. Nous pouvons donc supposer que
$$
(a, b, \alpha) \circ f_1 =
(a, b, \alpha) \circ f'_1 =
(a_2, b_2, \alpha_2).
$$
Écrivons $r = (c, \text{id}_V, \gamma) : (U', \phi' : V \to u(U'), x')$,
$f_1 = (a_1, b_1, \alpha_1)$ et $f'_1 = (a'_1, b_1, \alpha'_1)$.
Nous avons utilisé ici la condition $p_p(g_1) = p_p(g'_1)$.
Les égalités précédentes équivalent maintenant à
$a \circ a_1 = a \circ a'_1 = a_2 \circ c$ et
$\alpha \circ \alpha_1 = \alpha \circ \alpha'_1 = \alpha_2 \circ \gamma$.
Dans cette situation, il n'est pas nécessaire que $a_1 = a'_1$.
Nous devons donc précomposer par un morphisme supplémentaire appartenant à $R$.
Soit en effet $U'' = \text{Eq}(a_1, a'_1)$ l'égalisateur de
$a_1$ et $a'_1$, qui est un sous-objet de $U'$. Notons
$c' : U'' \to U'$ le monomorphisme canonique.
En raison des relations entre les morphismes
ci-dessus, nous voyons que $V \to u(U')$ se factorise par
$u(U'') = u(\text{Eq}(a_1, a'_1)) = \text{Eq}(u(a_1), u(a'_1))$.
Nous obtenons donc un nouvel objet $(U'', \phi'' : V \to u(U''), x'')$, où
$\gamma' : x'' \to x'$ est un morphisme fortement cartésien relevant $\gamma$.
Nous voyons alors que nous pouvons précomposer $f_1$ et $f'_1$ par l'élément
$(c', \text{id}_V, \gamma')$ de $R$. Après cette opération, c'est-à-dire après avoir remplacé
$(U', \phi' : V \to u(U'), x')$ par
$(U'', \phi'' : V \to u(U''), x'')$, nous nous retrouvons dans la situation
précédente, avec en outre $a_1 = a'_1$.
Dans ce cas, il résulte formellement du fait que $\alpha$
est fortement cartésien (!) que $\alpha_1 = \alpha'_1$.
Cela montre que $g_1 = g'_1$, comme souhaité.

\medskip\noindent
Nous omettons la démonstration du fait que, pour tout morphisme fortement cartésien
de $u_p\mathcal{S}$ de la forme $((a, b, \alpha), 1)$, le morphisme
$\alpha$ est fortement cartésien dans $\mathcal{S}$.
(Nous n'avons pas besoin de cette caractérisation des morphismes fortement cartésiens
dans la suite de la démonstration, bien que nous l'utilisions plus loin dans cette section.)

\medskip\noindent
Soit $(U, \phi : V \to u(U), x)$ un objet de $u_p\mathcal{S}$.
Soit $b : V' \to V$ un morphisme de $\mathcal{D}$. Alors le morphisme
$$
(\text{id}_U, b, \text{id}_x) :
(U, \phi \circ b : V' \to u(U), x)
\longrightarrow
(U, \phi : V \to u(U), x)
$$
est fortement cartésien d'après le résultat des paragraphes précédents, ce qui conclut la démonstration.
\end{proof}

\begin{lemma}
\label{lemma-fibred-groupoids-category-pullback}
Avec les notations et les hypothèses du
Lemme \ref{lemma-fibred-category-pullback},
si $\mathcal{S}$ est fibrée en groupoïdes, alors $u_p\mathcal{S}$ est fibrée
en groupoïdes.
\end{lemma}

\begin{proof}
D'après le
Lemme \ref{lemma-fibred-category-pullback},
nous savons que $u_p\mathcal{S}$ est une catégorie fibrée.
Soit $f : X \to Y$ un morphisme de $u_p\mathcal{S}$ tel que
$p_p(f) = \text{id}_V$. Il suffit de montrer que $f$ est inversible ; voir
Catégories, Lemme \ref{categories-lemma-fibred-groupoids}.
Écrivons $f$ comme la classe d'équivalence
d'un couple $((a, b, \alpha), r)$ avec $r \in R$. Alors $p_p(r) = \text{id}_V$,
donc $p_{pp}((a, b, \alpha)) = \text{id}_V$. Par conséquent, $b = \text{id}_V$.
Or tout morphisme de $\mathcal{S}$ est fortement cartésien ; voir
Catégories, Lemme \ref{categories-lemma-fibred-groupoids}.
Ainsi, $(a, b, \alpha) \in R$ est inversible
dans $u_p\mathcal{S}$, comme souhaité.
\end{proof}

\begin{lemma}
\label{lemma-adjointness-pullback-pushforward}
Soit $u : \mathcal{C} \to \mathcal{D}$ un foncteur.
Soient $p : \mathcal{S} \to \mathcal{C}$ et $q : \mathcal{T} \to \mathcal{D}$
des catégories au-dessus de $\mathcal{C}$ et de $\mathcal{D}$. Supposons que
\begin{enumerate}
\item $p : \mathcal{S} \to \mathcal{C}$ soit une catégorie fibrée,
\item $q : \mathcal{T} \to \mathcal{D}$ soit une catégorie fibrée,
\item $\mathcal{C}$ possède les limites finies non vides, et que
\item $u : \mathcal{C} \to \mathcal{D}$ commute aux limites finies non vides.
\end{enumerate}
Alors nous avons une équivalence canonique de catégories
$$
\Mor_{\textit{Fib}/\mathcal{C}}(\mathcal{S}, u^p\mathcal{T})
=
\Mor_{\textit{Fib}/\mathcal{D}}(u_p\mathcal{S}, \mathcal{T})
$$
entre les catégories de morphismes.
\end{lemma}

\begin{proof}
Dans cette démonstration, nous employons la notation $x/U$ pour désigner un objet
$x$ de $\mathcal{S}$ situé au-dessus de $U$ dans $\mathcal{C}$.
De même, $y/V$ désigne un objet $y$ de $\mathcal{T}$
situé au-dessus de $V$ dans $\mathcal{D}$. Dans le même esprit,
$\alpha/a : x/U \to x'/U'$ désigne le morphisme
$\alpha : x \to x'$ dont l'image est $a : U \to U'$ dans $\mathcal{C}$.

\medskip\noindent
Soit $G : u_p\mathcal{S} \to \mathcal{T}$ un $1$-morphisme de catégories fibrées
au-dessus de $\mathcal{D}$. Notons $G' : u_{pp}\mathcal{S} \to \mathcal{T}$
la composée de $G$ avec le foncteur canonique (de localisation)
$u_{pp}\mathcal{S} \to u_p\mathcal{S}$. Considérons alors le foncteur
$H : \mathcal{S} \to u^p\mathcal{T}$ donné par
$$
H(x/U) = (U, G'(U, \text{id}_{u(U)} : u(U) \to u(U), x))
$$
sur les objets et par
$$
H((\alpha, a) : x/U \to x'/U') = G'(a, u(a), \alpha)
$$
sur les morphismes. Comme $G$ transforme les morphismes fortement cartésiens
en morphismes fortement cartésiens, nous voyons que, si $\alpha$ est fortement
cartésien, alors $H(\alpha)$ est fortement cartésien.
En effet, nous avons vu dans la démonstration du
Lemme \ref{lemma-fibred-category-pullback}
que, dans ce cas, l'application $(a, u(a), \alpha)$ devient
fortement cartésienne dans $u_p\mathcal{S}$. Cette construction est manifestement
fonctorielle en $G$, et nous obtenons un foncteur
$$
A :
\Mor_{\textit{Fib}/\mathcal{D}}(u_p\mathcal{S}, \mathcal{T})
\longrightarrow
\Mor_{\textit{Fib}/\mathcal{C}}(\mathcal{S}, u^p\mathcal{T})
$$

\medskip\noindent
Réciproquement, soit $H : \mathcal{S} \to u^p\mathcal{T}$ un $1$-morphisme
de catégories fibrées. Rappelons qu'un objet de $u^p\mathcal{T}$ est
un couple $(U, y)$ avec $y \in \Ob(\mathcal{T}_{u(U)})$. Notons
$\text{pr} : u^p\mathcal{T} \to \mathcal{T}$ le foncteur $(U, y) \mapsto y$.
Dans ce cas, nous définissons un foncteur
$G' : u_{pp}\mathcal{S} \to \mathcal{T}$ par les règles
$$
G'(U, \phi : V \to u(U), x) = \phi^*\text{pr}(H(x))
$$
sur les objets et, sur les morphismes, par
$$
G'((a, b, \alpha) : (U, \phi : V \to u(U), x)
\to (U', \phi' : V' \to u(U'), x'))
=
\beta
$$
où $\beta : \phi^*\text{pr}(H(x)) \to (\phi')^*\text{pr}(H(x'))$
est l'unique morphisme
tel que $q(\beta) = b$ et que le diagramme
$$
\xymatrix{
\phi^*\text{pr}(H(x)) \ar[d] \ar[r]_-{\beta} &
(\phi')^*\text{pr}(H(x')) \ar[d] \\
\text{pr}(H(x)) \ar[r]^{\text{pr}(H(a, \alpha))} & \text{pr}(H(x'))
}
$$
soit commutatif. Un tel morphisme existe et est unique parce que $\mathcal{T}$ est une catégorie
fibrée.

\medskip\noindent
Vérifions que $G'(r)$ est un isomorphisme si $r \in R$.
En effet, si
$$
(a, \text{id}_V, \alpha) :
(U', \phi' : V \to u(U'), x')
\longrightarrow
(U, \phi : V \to u(U), x)
$$
avec $\alpha$ fortement cartésien est un élément du système multiplicatif à
droite $R$ du
Lemme \ref{lemma-right-multiplicative-system},
alors $H(\alpha)$ est fortement cartésien, et
$\text{pr}(H(\alpha))$ est fortement cartésien ; voir la démonstration du
Lemme \ref{lemma-fibred-category-pushforward}.
Dans ce cas, le morphisme $\beta$ vérifie donc $q(\beta) = \text{id}_V$
et est fortement cartésien. Ainsi, $\beta$ est un isomorphisme d'après
Catégories, Lemme \ref{categories-lemma-composition-cartesian}.
Par conséquent, d'après
Catégories, Lemme \ref{categories-lemma-properties-right-localization},
nous obtenons un prolongement canonique $G : u_p\mathcal{S} \to \mathcal{T}$.

\medskip\noindent
Montrons ensuite que $G$ transforme les morphismes fortement cartésiens
en morphismes fortement cartésiens.
Supposons que $f : X \to Y$ soit fortement cartésien. D'après la caractérisation
des morphismes fortement cartésiens dans $u_p\mathcal{S}$, nous pouvons écrire $f$ sous la forme
$((a, b, \alpha) : X' \to Y, r : X' \to Y)$, où $r \in R$ et où
$\alpha$ est fortement cartésien dans $\mathcal{S}$. D'après ce qui précède, il suffit
de montrer que $G'(a, b \alpha)$ est fortement cartésien. Comme précédemment,
la condition que $\alpha$ soit fortement cartésien implique que
$\text{pr}(H(a, \alpha)) : \text{pr}(H(x)) \to \text{pr}(H(x'))$
est fortement cartésien dans $\mathcal{T}$. Puisque, dans le carré commutatif
ci-dessus, toutes les flèches sauf peut-être $\beta$ sont maintenant fortement cartésiennes,
il s'ensuit que $\beta$ est lui aussi fortement cartésien, comme souhaité.
La construction $H \mapsto G$ est manifestement fonctorielle en $H$, et nous
obtenons un foncteur
$$
B :
\Mor_{\textit{Fib}/\mathcal{C}}(\mathcal{S}, u^p\mathcal{T})
\longrightarrow
\Mor_{\textit{Fib}/\mathcal{D}}(u_p\mathcal{S}, \mathcal{T})
$$
Pour achever la démonstration du lemme, il faut montrer que les foncteurs
$A$ et $B$ sont quasi-inverses l'un de l'autre. Nous omettons les vérifications.
\end{proof}

\begin{definition}
\label{definition-pullback-stack}
Soit $f : \mathcal{D} \to \mathcal{C}$ un morphisme de sites
donné par un foncteur continu $u : \mathcal{C} \to \mathcal{D}$
qui satisfait les hypothèses et les conclusions de
Sites, Proposition \ref{sites-proposition-get-morphism}.
Soit $\mathcal{S}$ un champ sur $\mathcal{C}$.
Dans ce cadre, nous notons {\it $f^{-1}\mathcal{S}$} la champification
de la catégorie fibrée $u_p\mathcal{S}$ au-dessus de $\mathcal{D}$ construite
ci-dessus. Nous disons que $f^{-1}\mathcal{S}$ est
l'{\it image inverse de $\mathcal{S}$ par $f$}.
\end{definition}

\noindent
Bien entendu, si $\mathcal{S}$ est un champ en groupoïdes, alors $f^{-1}\mathcal{S}$
est un champ en groupoïdes d'après les
Lemmes \ref{lemma-stackify-groupoids} et
\ref{lemma-fibred-groupoids-category-pullback}.

\begin{lemma}
\label{lemma-adjointness-pullback-pushforward-stacks}
Soit $f : \mathcal{D} \to \mathcal{C}$ un morphisme de sites
donné par un foncteur continu $u : \mathcal{C} \to \mathcal{D}$
qui satisfait les hypothèses et les conclusions de
Sites, Proposition \ref{sites-proposition-get-morphism}.
Soient $p : \mathcal{S} \to \mathcal{C}$ et
$q : \mathcal{T} \to \mathcal{D}$ des champs.
Alors nous avons une équivalence canonique de catégories
$$
\Mor_{\textit{Stacks}/\mathcal{C}}(\mathcal{S}, f_*\mathcal{T})
=
\Mor_{\textit{Stacks}/\mathcal{D}}(f^{-1}\mathcal{S}, \mathcal{T})
$$
entre les catégories de morphismes.
\end{lemma}

\begin{proof}
Pour $i = 1, 2$, un $i$-morphisme de champs est la même chose qu'un
$i$-morphisme de catégories fibrées ; voir la
Définition \ref{definition-stacks-over-C}.
D'après le
Lemme \ref{lemma-adjointness-pullback-pushforward},
nous avons déjà
$$
\Mor_{\textit{Fib}/\mathcal{C}}(\mathcal{S}, u^p\mathcal{T})
=
\Mor_{\textit{Fib}/\mathcal{D}}(u_p\mathcal{S}, \mathcal{T})
$$
Le résultat découle donc du
Lemme \ref{lemma-stackify-universal-property-more},
puisque $u^p\mathcal{T} = f_*\mathcal{T}$ et que $f^{-1}\mathcal{S}$ est la
champification de $u_p\mathcal{S}$.
\end{proof}

\begin{lemma}
\label{lemma-technical-up}
Soit $f : \mathcal{D} \to \mathcal{C}$ un morphisme de sites
donné par un foncteur continu $u : \mathcal{C} \to \mathcal{D}$
qui satisfait les hypothèses et les conclusions de
Sites, Proposition \ref{sites-proposition-get-morphism}.
Soit $\mathcal{S} \to \mathcal{C}$ une catégorie fibrée, et
soit $\mathcal{S} \to \mathcal{S}'$ la champification de $\mathcal{S}$.
Alors $f^{-1}\mathcal{S}'$ est la champification de
$u_p\mathcal{S}$.
\end{lemma}

\begin{proof}
Omis. Indication :
C'est l'analogue de Sites, Lemme \ref{sites-lemma-technical-up}.
\end{proof}

\noindent
Le lemme suivant nous dit que la $2$-catégorie des champs
sur $\Sch_{fppf}$ est une ``2-sous-catégorie pleine'' de la
$2$-catégorie des champs sur $\Sch'_{fppf}$ dès que
$\Sch'_{fppf}$ contient $\Sch_{fppf}$ (voir
Topologies, Section \ref{topologies-section-change-alpha}).

\begin{lemma}
\label{lemma-bigger-site}
Soient $\mathcal{C}$ et $\mathcal{D}$ des sites.
Soit $u : \mathcal{C} \to \mathcal{D}$ un foncteur satisfaisant aux
hypothèses de
Sites, Lemme \ref{sites-lemma-bigger-site}.
Soit $f : \mathcal{D} \to \mathcal{C}$ le morphisme de sites
correspondant. Alors
\begin{enumerate}
\item pour tout champ $p : \mathcal{S} \to \mathcal{C}$, le
foncteur canonique $\mathcal{S} \to f_*f^{-1}\mathcal{S}$ est
une équivalence de champs ;
\item étant donnés des champs $\mathcal{S}, \mathcal{S}'$ sur $\mathcal{C}$,
la construction $f^{-1}$
induit une équivalence
$$
\Mor_{\textit{Stacks}/\mathcal{C}}(\mathcal{S}, \mathcal{S}')
\longrightarrow
\Mor_{\textit{Stacks}/\mathcal{D}}(f^{-1}\mathcal{S}, f^{-1}\mathcal{S}')
$$
entre les catégories de morphismes.
\end{enumerate}
\end{lemma}

\begin{proof}
Remarquons que, d'après le
Lemme \ref{lemma-adjointness-pullback-pushforward-stacks},
nous avons une équivalence de catégories
$$
\Mor_{\textit{Stacks}/\mathcal{D}}(f^{-1}\mathcal{S}, f^{-1}\mathcal{S}')
=
\Mor_{\textit{Stacks}/\mathcal{C}}(\mathcal{S}, f_*f^{-1}\mathcal{S}')
$$
Ainsi, (2) découle de (1).

\medskip\noindent
Pour démontrer (1), nous allons utiliser le
Lemme \ref{lemma-characterize-essentially-surjective-when-ff}.
Ce lemme nous dit qu'il faut montrer que
$can : \mathcal{S} \to f_*f^{-1}\mathcal{S}$ est
pleinement fidèle et que tous les objets de $f_*f^{-1}\mathcal{S}$
sont localement dans l'image essentielle.

\medskip\noindent
Décrivons rapidement le foncteur $can$ ; voir la démonstration du
Lemme \ref{lemma-adjointness-pullback-pushforward}.
Pour ce faire, introduisons le foncteur
$c'' : \mathcal{S} \to u_{pp}\mathcal{S}$ défini par
$c''(x/U) = (U, \text{id} : u(U) \to u(U), x)$ et
$c''(\alpha/a) = (a, u(a), \alpha)$. Nous définissons
$c' : \mathcal{S} \to u_p\mathcal{S}$ comme la composée
de $c''$ et du foncteur canonique $u_{pp}\mathcal{S} \to u_p\mathcal{S}$.
Nous définissons $c : \mathcal{S} \to f^{-1}\mathcal{S}$ comme la composée
de $c'$ et du foncteur canonique $u_p\mathcal{S} \to f^{-1}\mathcal{S}$.
Alors $can : \mathcal{S} \to f_*f^{-1}\mathcal{S}$ est le foncteur
qui associe à $x/U$ le couple $(U, c(x))$ et à
$\alpha/a$ le morphisme $(a, c(\alpha))$.

\medskip\noindent
Pleine fidélité. Pour la démontrer, nous allons utiliser le
Lemme \ref{lemma-characterize-ff}.
Soit $U \in \Ob(\mathcal{C})$.
Soient $x, y \in \mathcal{S}_U$.
Tout d'abord, comme $u$ est pleinement fidèle, nous avons
$$
\Mor_{(f_*f^{-1}\mathcal{S})_U}(can(x), can(y))
=
\Mor_{(f^{-1}\mathcal{S})_{u(U)}}(c(x), c(y))
$$
directement par la définition de $f_*$. Il en va de même après
changement de base à tout $U'/U$. Puisque $f^{-1}\mathcal{S}$ est la
champification de $u_p\mathcal{S}$, et puisque $u$ est continu
et cocontinu, le préfaisceau
$$
U'/U \longmapsto
\Mor_{(f^{-1}\mathcal{S})_{u(U')}}(c(x|_{U'}), c(y|_{U'}))
$$
est le faisceau associé au préfaisceau
$$
U'/U \longmapsto
\Mor_{(u_p\mathcal{S})_{u(U')}}(c'(x|_{U'}), c'(y|_{U'}))
$$
Ainsi, pour achever la démonstration de la pleine fidélité, il suffit
de montrer que, quels que soient $U$ et $x, y$, l'application
$$
\Mor_{\mathcal{S}_U}(x, y)
\longrightarrow
\Mor_{(u_p\mathcal{S})_U}(c'(x), c'(y))
$$
est bijective. Un morphisme $f : x \to y$ dans $u_p\mathcal{S}$ au-dessus de $u(U)$
est donné par une classe d'équivalence de diagrammes
$$
\xymatrix{
(U', \phi : u(U) \to u(U'), x')
\ar[d]_{(c, \text{id}_{u(U)}, \gamma)}
\ar[r]_{(a, b, \alpha)} &
(U, \text{id} : u(U) \to u(U), y)
\\
(U, \text{id} : u(U) \to u(U), x)
}
$$
avec $\gamma$ fortement cartésien et $b = \text{id}_{u(U)}$.
Mais, puisque $u$ est pleinement fidèle, nous pouvons écrire $\phi = u(c')$ pour un
morphisme $c' : U \to U'$ ;
nous voyons alors que $a \circ c' = \text{id}_U$ et
$c \circ c' = \text{id}_{U'}$. Comme $\gamma$ est fortement cartésien,
nous pouvons trouver un morphisme $\gamma' : x \to x'$ relevant $c'$ tel que
$\gamma \circ \gamma' = \text{id}_x$. Par définition des classes
d'équivalence qui définissent les morphismes dans $u_p\mathcal{S}$, il s'ensuit que
le morphisme
$$
\xymatrix{
(U, \text{id} : u(U) \to u(U), x)
\ar[rr]_{(\text{id}, \text{id}, \alpha \circ \gamma')} & &
(U, \text{id} : u(U) \to u(U), y)
}
$$
de $u_{pp}\mathcal{S}$ induit le morphisme $f$ dans $u_p\mathcal{S}$.
Cela prouve que l'application est surjective. Nous omettons la démonstration de son
injectivité.

\medskip\noindent
Enfin, nous devons montrer que tout objet de $f_*f^{-1}\mathcal{S}$
provient localement d'un objet de $\mathcal{S}$. Cela ressort des
constructions (détails omis).
\end{proof}




\section{Champs et localisation}
\label{section-localize}

\noindent
Soit $\mathcal{C}$ un site.
Soit $U$ un objet de $\mathcal{C}$.
Nous voulons interpréter les champs sur $\mathcal{C}/U$ comme des champs
sur $\mathcal{C}$ munis d'un morphisme vers $U$.
Le lemme suivant explique pourquoi cela est plus facile
lorsque le préfaisceau $h_U$ est un faisceau.

\begin{lemma}
\label{lemma-when-localization-stack}
Soit $\mathcal{C}$ un site. Soit $U \in \Ob(\mathcal{C})$.
Alors $j_U : \mathcal{C}/U \to \mathcal{C}$ est un champ sur $\mathcal{C}$
si et seulement si $h_U$ est un faisceau.
\end{lemma}

\begin{proof}
Combiner le
Lemme \ref{lemma-stack-in-setoids-characterize}
avec
Catégories,
Exemple \ref{categories-example-fibred-category-from-functor-of-points}.
\end{proof}

\noindent
Supposons que $\mathcal{C}$ soit un site
et que $U$ soit un objet de $\mathcal{C}$ dont le préfaisceau représentable
associé soit un faisceau. Nous notons $j : \mathcal{C}/U \to \mathcal{C}$
le foncteur de localisation.

\medskip\noindent
{\bf Construction A.} Soit $p : \mathcal{S} \to \mathcal{C}/U$ un champ
sur le site $\mathcal{C}/U$. Nous définissons un champ
$j_!p : j_!\mathcal{S} \to \mathcal{C}$ comme suit :
\begin{enumerate}
\item En tant que catégorie, $j_!\mathcal{S} = \mathcal{S}$, et
\item le foncteur $j_!p : j_!\mathcal{S} \to \mathcal{C}$
est simplement la composée $j \circ p$.
\end{enumerate}
Nous omettons la vérification qu'il s'agit d'un champ (indication : utiliser le fait que
$h_U$ est un faisceau pour recoller les morphismes vers $U$). Il existe un foncteur
canonique
$$
j_!\mathcal{S} \longrightarrow \mathcal{C}/U
$$
à savoir le foncteur $p$, qui est un $1$-morphisme de champs sur $\mathcal{C}$.

\medskip\noindent
{\bf Construction B.}
Soit $q : \mathcal{T} \to \mathcal{C}$ un champ sur $\mathcal{C}$
muni d'un morphisme de champs $p : \mathcal{T} \to \mathcal{C}/U$
sur $\mathcal{C}$. Alors $p : \mathcal{T} \to \mathcal{C}/U$
est automatiquement un champ sur $\mathcal{C}/U$.

\begin{lemma}
\label{lemma-localize-stacks}
Supposons que $\mathcal{C}$ soit un site
et que $U$ soit un objet de $\mathcal{C}$ dont le préfaisceau représentable
associé soit un faisceau. Les constructions A et B ci-dessus définissent
des foncteurs de $2$-catégories mutuellement inverses (!)
$$
\left\{
\begin{matrix}
2\text{-catégorie des}\\
\text{champs sur }\mathcal{C}/U
\end{matrix}
\right\}
\leftrightarrow
\left\{
\begin{matrix}
2\text{-catégorie des couples }(\mathcal{T}, p)
\text{ constitués} \\
\text{d'un champ }\mathcal{T}\text{ sur }\mathcal{C}\text{ et d'un morphisme} \\
p : \mathcal{T} \to \mathcal{C}/U\text{ de champs sur }\mathcal{C}
\end{matrix}
\right\}
$$
\end{lemma}

\begin{proof}
C'est clair.
\end{proof}



\begin{multicols}{2}[\section{Autres chapitres}]
\noindent
Préliminaires
\begin{enumerate}
\item \hyperref[introduction-section-phantom]{Introduction}
\item \hyperref[conventions-section-phantom]{Conventions}
\item \hyperref[sets-section-phantom]{Théorie des ensembles}
\item \hyperref[categories-section-phantom]{Catégories}
\item \hyperref[topology-section-phantom]{Topologie}
\item \hyperref[sheaves-section-phantom]{Faisceaux sur les espaces}
\item \hyperref[sites-section-phantom]{Sites et faisceaux}
\item \hyperref[stacks-section-phantom]{Champs}
\item \hyperref[fields-section-phantom]{Corps}
\item \hyperref[algebra-section-phantom]{Algèbre commutative}
\item \hyperref[brauer-section-phantom]{Groupes de Brauer}
\item \hyperref[homology-section-phantom]{Algèbre homologique}
\item \hyperref[derived-section-phantom]{Catégories dérivées}
\item \hyperref[simplicial-section-phantom]{Méthodes simpliciales}
\item \hyperref[more-algebra-section-phantom]{Compléments d'algèbre}
\item \hyperref[smoothing-section-phantom]{Lissage des morphismes d'anneaux}
\item \hyperref[modules-section-phantom]{Faisceaux de modules}
\item \hyperref[sites-modules-section-phantom]{Modules sur les sites}
\item \hyperref[injectives-section-phantom]{Objets injectifs}
\item \hyperref[cohomology-section-phantom]{Cohomologie des faisceaux}
\item \hyperref[sites-cohomology-section-phantom]{Cohomologie sur les sites}
\item \hyperref[dga-section-phantom]{Algèbre différentielle graduée}
\item \hyperref[dpa-section-phantom]{Algèbre à puissances divisées}
\item \hyperref[sdga-section-phantom]{Faisceaux différentiels gradués}
\item \hyperref[hypercovering-section-phantom]{Hyperrecouvrements}
\end{enumerate}
Schémas
\begin{enumerate}
\setcounter{enumi}{25}
\item \hyperref[schemes-section-phantom]{Schémas}
\item \hyperref[constructions-section-phantom]{Constructions de schémas}
\item \hyperref[properties-section-phantom]{Propriétés des schémas}
\item \hyperref[morphisms-section-phantom]{Morphismes de schémas}
\item \hyperref[coherent-section-phantom]{Cohomologie des schémas}
\item \hyperref[divisors-section-phantom]{Diviseurs}
\item \hyperref[limits-section-phantom]{Limites de schémas}
\item \hyperref[varieties-section-phantom]{Variétés}
\item \hyperref[topologies-section-phantom]{Topologies sur les schémas}
\item \hyperref[descent-section-phantom]{Descente}
\item \hyperref[perfect-section-phantom]{Catégories dérivées de schémas}
\item \hyperref[more-morphisms-section-phantom]{Compléments sur les morphismes}
\item \hyperref[flat-section-phantom]{Compléments sur la platitude}
\item \hyperref[groupoids-section-phantom]{Schémas en groupoïdes}
\item \hyperref[more-groupoids-section-phantom]{Compléments sur les schémas en groupoïdes}
\item \hyperref[etale-section-phantom]{Morphismes étales de schémas}
\end{enumerate}
Thèmes de théorie des schémas
\begin{enumerate}
\setcounter{enumi}{41}
\item \hyperref[chow-section-phantom]{Homologie de Chow}
\item \hyperref[intersection-section-phantom]{Théorie de l'intersection}
\item \hyperref[pic-section-phantom]{Schémas de Picard des courbes}
\item \hyperref[weil-section-phantom]{Théories cohomologiques de Weil}
\item \hyperref[adequate-section-phantom]{Modules adéquats}
\item \hyperref[dualizing-section-phantom]{Complexes dualisants}
\item \hyperref[duality-section-phantom]{Dualité pour les schémas}
\item \hyperref[discriminant-section-phantom]{Discriminants et différentes}
\item \hyperref[derham-section-phantom]{Cohomologie de de Rham}
\item \hyperref[local-cohomology-section-phantom]{Cohomologie locale}
\item \hyperref[algebraization-section-phantom]{Géométrie algébrique et formelle}
\item \hyperref[curves-section-phantom]{Courbes algébriques}
\item \hyperref[resolve-section-phantom]{Résolution des surfaces}
\item \hyperref[models-section-phantom]{Réduction semi-stable}
\item \hyperref[functors-section-phantom]{Foncteurs et morphismes}
\item \hyperref[equiv-section-phantom]{Catégories dérivées de variétés}
\item \hyperref[pione-section-phantom]{Groupes fondamentaux de schémas}
\item \hyperref[etale-cohomology-section-phantom]{Cohomologie étale}
\item \hyperref[crystalline-section-phantom]{Cohomologie cristalline}
\item \hyperref[proetale-section-phantom]{Cohomologie pro-étale}
\item \hyperref[relative-cycles-section-phantom]{Cycles relatifs}
\item \hyperref[more-etale-section-phantom]{Compléments de cohomologie étale}
\item \hyperref[trace-section-phantom]{La formule des traces}
\end{enumerate}
Espaces algébriques
\begin{enumerate}
\setcounter{enumi}{64}
\item \hyperref[spaces-section-phantom]{Espaces algébriques}
\item \hyperref[spaces-properties-section-phantom]{Propriétés des espaces algébriques}
\item \hyperref[spaces-morphisms-section-phantom]{Morphismes d'espaces algébriques}
\item \hyperref[decent-spaces-section-phantom]{Espaces algébriques décents}
\item \hyperref[spaces-cohomology-section-phantom]{Cohomologie des espaces algébriques}
\item \hyperref[spaces-limits-section-phantom]{Limites d'espaces algébriques}
\item \hyperref[spaces-divisors-section-phantom]{Diviseurs sur les espaces algébriques}
\item \hyperref[spaces-over-fields-section-phantom]{Espaces algébriques sur des corps}
\item \hyperref[spaces-topologies-section-phantom]{Topologies sur les espaces algébriques}
\item \hyperref[spaces-descent-section-phantom]{Descente et espaces algébriques}
\item \hyperref[spaces-perfect-section-phantom]{Catégories dérivées d'espaces}
\item \hyperref[spaces-more-morphisms-section-phantom]{Compléments sur les morphismes d'espaces}
\item \hyperref[spaces-flat-section-phantom]{Platitude sur les espaces algébriques}
\item \hyperref[spaces-groupoids-section-phantom]{Groupoïdes dans les espaces algébriques}
\item \hyperref[spaces-more-groupoids-section-phantom]{Compléments sur les groupoïdes dans les espaces}
\item \hyperref[bootstrap-section-phantom]{Amorçage}
\item \hyperref[spaces-pushouts-section-phantom]{Sommes amalgamées d'espaces algébriques}
\end{enumerate}
Thèmes de géométrie
\begin{enumerate}
\setcounter{enumi}{81}
\item \hyperref[spaces-chow-section-phantom]{Groupes de Chow des espaces}
\item \hyperref[groupoids-quotients-section-phantom]{Quotients de groupoïdes}
\item \hyperref[spaces-more-cohomology-section-phantom]{Compléments sur la cohomologie des espaces}
\item \hyperref[spaces-simplicial-section-phantom]{Espaces simpliciaux}
\item \hyperref[spaces-duality-section-phantom]{Dualité pour les espaces}
\item \hyperref[formal-spaces-section-phantom]{Espaces algébriques formels}
\item \hyperref[restricted-section-phantom]{Algébrisation des espaces formels}
\item \hyperref[spaces-resolve-section-phantom]{Résolution des surfaces revisitée}
\end{enumerate}
Théorie des déformations
\begin{enumerate}
\setcounter{enumi}{89}
\item \hyperref[formal-defos-section-phantom]{Théorie formelle des déformations}
\item \hyperref[defos-section-phantom]{Théorie des déformations}
\item \hyperref[cotangent-section-phantom]{Le complexe cotangent}
\item \hyperref[examples-defos-section-phantom]{Problèmes de déformation}
\end{enumerate}
Champs algébriques
\begin{enumerate}
\setcounter{enumi}{93}
\item \hyperref[algebraic-section-phantom]{Champs algébriques}
\item \hyperref[examples-stacks-section-phantom]{Exemples de champs}
\item \hyperref[stacks-sheaves-section-phantom]{Faisceaux sur les champs algébriques}
\item \hyperref[criteria-section-phantom]{Critères de représentabilité}
\item \hyperref[artin-section-phantom]{Axiomes d'Artin}
\item \hyperref[quot-section-phantom]{Espaces de Quot et de Hilbert}
\item \hyperref[stacks-properties-section-phantom]{Propriétés des champs algébriques}
\item \hyperref[stacks-morphisms-section-phantom]{Morphismes de champs algébriques}
\item \hyperref[stacks-limits-section-phantom]{Limites de champs algébriques}
\item \hyperref[stacks-cohomology-section-phantom]{Cohomologie des champs algébriques}
\item \hyperref[stacks-perfect-section-phantom]{Catégories dérivées de champs}
\item \hyperref[stacks-introduction-section-phantom]{Introduction aux champs algébriques}
\item \hyperref[stacks-more-morphisms-section-phantom]{Compléments sur les morphismes de champs}
\item \hyperref[stacks-geometry-section-phantom]{La géométrie des champs}
\end{enumerate}
Thèmes de théorie des espaces de modules
\begin{enumerate}
\setcounter{enumi}{107}
\item \hyperref[moduli-section-phantom]{Champs de modules}
\item \hyperref[moduli-curves-section-phantom]{Espaces de modules de courbes}
\end{enumerate}
Divers
\begin{enumerate}
\setcounter{enumi}{109}
\item \hyperref[examples-section-phantom]{Exemples}
\item \hyperref[exercises-section-phantom]{Exercices}
\item \hyperref[guide-section-phantom]{Guide bibliographique}
\item \hyperref[desirables-section-phantom]{Desiderata}
\item \hyperref[coding-section-phantom]{Style de programmation}
\item \hyperref[obsolete-section-phantom]{Obsolète}
\item \hyperref[fdl-section-phantom]{Licence de documentation libre GNU}
\item \hyperref[index-section-phantom]{Index généré automatiquement}
\end{enumerate}
\end{multicols}



\bibliography{my}
\bibliographystyle{amsalpha}

\end{document}
